%% 
%% Copyright 2019-2024 Elsevier Ltd
%% 
%% Version 2.4
%% 
%% This file is part of the 'CAS Bundle'.
%% --------------------------------------
%% 
%% It may be distributed under the conditions of the LaTeX Project Public
%% License, either version 1.2 of this license or (at your option) any
%% later version.  The latest version of this license is in
%%    http://www.latex-project.org/lppl.txt
%% and version 1.2 or later is part of all distributions of LaTeX
%% version 1999/12/01 or later.
%% 
%% The list of all files belonging to the 'CAS Bundle' is
%% given in the file `manifest.txt'.
%% 
%% Template article for cas-sc documentclass for 
%% single column output.

%\documentclass[a4paper,fleqn,longmktitle]{cas-sc}
\documentclass[a4paper,fleqn]{cas-sc}


\usepackage[numbers]{natbib}
%\usepackage[authoryear]{natbib}
% \usepackage[authoryear,longnamesfirst]{natbib}
\usepackage{soul}
\usepackage{float}
\usepackage{placeins}
%%Pseudocode
\usepackage{algorithm}
\usepackage{algorithmic}
\usepackage{amsmath}
\usepackage[utf8]{inputenc}
\usepackage{subcaption}
\usepackage{graphicx}
\usepackage{pdflscape}
%\usepackage[margin=1in]{geometry} % Adjust margins to ensure the table fits
\usepackage{tabularx} % Allows for table width to be set
\usepackage{ltablex} % Allows tables to break across pages
%\usepackage{authblk}
\usepackage{cleveref}
%imported extra
\usepackage[utf8]{inputenc} % allow utf-8 input
\usepackage[T1]{fontenc}    % use 8-bit T1 fonts
% \usepackage[colorlinks, citecolor=green, linkcolor=red]{hyperref}      % hyperlinks
% \usepackage{url}            % simple URL typesetting
% \usepackage{booktabs}       % professional-quality tables
% \usepackage{amsfonts}       % blackboard math symbols
% \usepackage{nicefrac}       % compact symbols for 1/2, etc.
% \usepackage{microtype}      % microtypography
% \usepackage{lipsum}		  % Can be removed after putting your text content
% \usepackage{graphicx}
% \usepackage{doi}
% \usepackage{authblk}
% \usepackage{framed,multirow}

%% The amssymb package provides various useful mathematical symbols

% \usepackage{amssymb}asa
% \usepackage{latexsym}
% \usepackage{graphicx}
% \usepackage{amsmath}
% \usepackage{amssymb}
% \usepackage{booktabs} % for better looking tables
% \usepackage{tabularx}
\usepackage{array}
% \usepackage{tabularx, multirow}
\usepackage{tabularx, multirow, booktabs}
% \usepackage{subcaption}
% \usepackage{algorithm}
% \usepackage{algorithmic}
% \usepackage{float}
\usepackage{verbatim}
\usepackage{ragged2e}
\usepackage{longtable} 
\usepackage{xcolor}
%%%Author macros
\def\tsc#1{\csdef{#1}{\textsc{\lowercase{#1}}\xspace}}
\tsc{WGM}
\tsc{QE}
\tsc{EP}
\tsc{PMS}
\tsc{BEC}
\tsc{DE}
%%%

\begin{document}
\let\WriteBookmarks\relax
\def\floatpagepagefraction{1}
\def\textpagefraction{.001}
\shorttitle{LSC-XAI}
\shortauthors{A.R. Noor et~al.}
%\begin{frontmatter}

% \title [mode = title]{This is a specimen $a_b$ title}                      
% \title [mode = title]{LSC-XAI: A Hybrid Explainable AI Framework for Predicting and Mitigating Healthcare Worker Attrition}                     

\title [mode = title]{An Interpretable Machine Learning Framework for Predicting and Managing Healthcare Worker Attrition}                     

% \tnotemark[1,2]

% \tnotetext[1]{This document is the results of the research
%    project funded by the National Science Foundation.}

% \tnotetext[2]{The second title footnote which is a longer text matter
%    to fill through the whole text width and overflow into
%    another line in the footnotes area of the first page.}

% \affiliation[1]{organization={Department of Computer Science and Engineering, Ahsanullah University of Science and Technology}, 
%                 city={Dhaka},
%                 country={Bangladesh}}
                

                
% \author[1,2]{Abrar Rafid Noor}
% \author[1,2]{Farhan Sharukh Hasan}
% \author[3]{Md Shawmoon Azad}[orcid=0009-0005-1318-072X]
% \author[3, 4]{M.R.C. Mahdy}[orcid=0000-0003-3737-0315]
% \cormark[1]
% \ead{mahdy.chowdhury@northsouth.edu}
% \cortext[cor1]{Corresponding author; email: mahdy.chowdhury@northsouth.edu}

% \affiliation[2]{organization={Mahdy Research Academy},
%     addressline={Bashundhara R/A},
%     city={Dhaka},
%     postcode={1229},
%     country={Bangladesh}}
% \affiliation[3]{organization={Department of Electrical and Computer Engineering, North South University},
%     addressline={Bashundhara R/A},
%     city={Dhaka},
%     postcode={1229},
%     country={Bangladesh}}

% \affiliation[4]{organization={NSU Center of Quantum Computing, North South University},
%     addressline={Bashundhara R/A},
%     city={Dhaka},
%     postcode={1229},
%     country={Bangladesh}}    

\begin{abstract}
Healthcare worker attrition poses a serious threat to workforce stability, quality of patient care and organizational efficiency. High turnover disrupts service delivery, increases operational costs and contributes to staff burnout. This study introduces the LIME SHAP Counterfactual Explainable Artificial Intelligence framework (LSC-XAI), a novel hybrid explainable artificial intelligence (XAI) system designed to predict and mitigate healthcare worker attrition. The framework integrates Local Interpretable Model-agnostic Explanations (LIME), Shapley Additive Explanations (SHAP) and counterfactual reasoning to identify the key factors influencing attrition and to generate actionable retention strategies. Fourteen ensemble machine learning and deep learning models are evaluated using a comprehensive healthcare employee attrition dataset, with rigorous preprocessing that includes feature selection, distributional transformation and data upsampling to ensure robust and reliable performance. \textcolor{blue}{Logistic Regression served as the framework's base model and achieved the highest mean cross-validated macro-averaged recall among the evaluated models (83.11 $\pm$ 1.34\%; mean $\pm$ standard deviation). CatBoost achieved the highest mean accuracy (91.41 $\pm$ 1.63\%), compared with a majority-class baseline accuracy of 88.13\% obtained by always predicting non-attrition (1,477 of 1,676 records); this baseline has zero attrition recall.} Using the LSC-XAI framework, we identify critical drivers of attrition, including job level, shift structure and overtime. The framework also produces frequency-based and interpretable recommendations that support targeted managerial interventions. These data-driven insights enable healthcare organizations and policymakers to design more effective retention strategies, improve workforce satisfaction and optimize overall operational efficiency.
\end{abstract}


% \begin{graphicalabstract}
% \includegraphics{figs/cas-grabs.pdf}
% \end{graphicalabstract}

\begin{highlights}
\item Develop an explainable analytics framework for predicting healthcare worker attrition.
\item Integrate machine learning and interpretability to support healthcare workforce decisions.
\item Identify key job and workload drivers of employee turnover in healthcare.
\item Generate interpretable recommendations for improving staff retention and stability.
\item Support data-driven workforce planning in healthcare organizations.
\end{highlights}

% \begin{keywords}
% Employee attrition  \sep Machine learning \sep Explainable AI \sep LIME \sep SHAP \sep Counterfactual
% \end{keywords}

\begin{keywords}
Human resource analytics \sep Explainable artificial intelligence \sep Machine learning \sep Healthcare workforce retention \sep Predictive modeling \sep Employee attrition
\end{keywords}

\maketitle

\section{Introduction}
\label{sec:intr}
Attrition among healthcare workers has become one of the most pressing challenges confronting contemporary healthcare systems across the globe. In addition to disrupting organizational continuity, the steady decline in healthcare personnel jeopardizes healthcare organizations' core capacity to provide high-quality patient care. The 2024 NSI National Healthcare Retention and RN Staffing Report \cite{NSI2024} states that the hospital turnover rate is 20.7\% and the registered nurse turnover rate is 18.4\%.

\textcolor{black}{High staff turnover in healthcare settings is well-documented to increase medication errors, infection rates and hospital length of stay, as demonstrated by several key studies in patient safety and workforce literature. For instance, Dong et al. \cite{dong2022} found that nursing home staff turnover directly correlates with higher infection control citations due to disrupted care continuity and inexperienced staffing, while Castle \cite{castle2005} linked turnover to increased re-hospitalizations and quality deficiencies like restraint use that prolong stays. Similarly, Aiken et al. \cite{aiken2014} showed in a multi-country analysis that nurse shortages from turnover elevate adverse outcomes including medication errors and healthcare-associated infections}. Additionally, the remaining employees frequently deal with more work and stress, which can start a risky cycle that lowers care quality and increases turnover.
There are many interrelated contributing factors to attrition among healthcare workers. The primary contributing factors include professional burnout — characterized by emotional exhaustion and depersonalization — workplace stress exacerbated by insufficient staffing ratios, limited career advancement opportunities, inadequate compensation relative to workload and difficulties in maintaining work–life balance \cite{Maslach2016, NAM2017}. These challenges are further compounded by the distinctive characteristics of the healthcare sector, such as its continuous operations, emotionally taxing work demands and high-stakes decision-making settings.
The attrition problem has become even more complex as a result of demographic changes. Healthcare systems must deal with both the retirement of seasoned healthcare professionals and the growing demand for services as the world's population ages. Concurrently, younger healthcare professionals tend to place greater value on work–life balance and career development, reflecting a broader shift in workforce expectations \cite{LavoieTremblay2010}.
The COVID-19 pandemic has acted as a trigger, escalating already-existing workforce issues and generating new ones. Unprecedented levels of stress, trauma and professional demands have been experienced by healthcare workers, prompting many to reconsider their career trajectories \cite{Shah2020}. The pandemic has also exposed systemic deficiencies in healthcare workforce planning and retention strategies, compelling organizations to reconsider their methods for supporting and engaging their workforce.
Changes in healthcare worker attrition by geography pose additional difficulties for policymakers and workforce planners. Urban and rural regions frequently have different turnover patterns, with rural facilities usually having more serious staffing issues because they have fewer resources and recruitment pools. Similarly, attrition patterns vary considerably between acute care hospitals and long-term care facilities, necessitating customized retention strategies \cite{Kovner2014}.
Although previous studies have examined various dimensions of healthcare workforce attrition, there remains a pressing need for a comprehensive and integrative analysis that accounts for the interplay among systemic, organizational and individual-level factors. Understanding these interdependencies is vital for designing effective and evidence-based retention strategies.

The mitigation of healthcare professional attrition remains a critical imperative for ensuring the longitudinal resilience of healthcare systems, a challenge recently prioritized in \textit{Healthcare Analytics} through the development of multicriteria decision models for workforce retention \cite{DICAPRIO2025100403}. While these foundational models provide a robust framework for identifying retention priorities based on employee satisfaction, there is an evolving need for predictive architectures capable of navigating the high-dimensional, non-linear complexities of clinician turnover. Crucially, as demonstrated by contemporary methodological advancements in this journal, the utility of such predictive tools is fundamentally contingent upon their interpretability \cite{BAIRY2025100407}. By synthesizing deep learning frameworks with Explainable AI (XAI) techniques, this study directly addresses the journal’s call for transparent, clinically-actionable analytics that bridge the gap between complex algorithmic outputs and the pragmatic decision-making requirements of hospital administrators.

This study aims to address this gap by investigating the multifaceted phenomenon of employee attrition in healthcare environments through the lens of predictive analytics and explainable artificial intelligence (XAI). Specifically, we employ state-of-the-art machine learning techniques to forecast employee turnover and propose a novel hybrid XAI framework (shown in Fig. \ref{fig:lcsxai_flow}) to elucidate the most salient predictors influencing attrition. By doing so, we seek to generate actionable insights that inform policy interventions and organizational strategies aimed at improving workforce stability in healthcare institutions.
The research is guided by the following questions:
\begin{itemize}
 
    \item [RQ1:] How do different categories of machine learning models, namely shallow learners, ensemble methods and deep learning models, compare in their effectiveness at predicting healthcare employee attrition?

    \item [RQ2:] Can a hybrid interpretability framework be developed by integrating multiple XAI techniques (e.g., SHAP, LIME, counterfactual explanations) to robustly identify key drivers of employee turnover and provide domain-relevant recommendations for mitigation?
    
    \item [RQ3:] Can the hybrid interpretable framework provide better direction compared to the existing XAI frameworks?
\end{itemize}

For healthcare administrators, stakeholders and HR specialists tasked with creating long-term healthcare workforces, the study's conclusions have important implications. Organizations can create more focused and efficient retention strategies to improve patient care delivery quality and consistency by better understanding the dynamics of healthcare worker attrition.
\begin{figure}[pos=H]
    \centering
    \includegraphics[width=0.6\linewidth, height=7em]{xai_flow_2.png}
    \caption{
        Overview of the proposed LSC-XAI framework. The pipeline integrates three stages of interpretability: local explanations via LIME for true positive predictions, global feature attributions via SHAP for both true positives and false negatives and actionable counterfactual reasoning on the intersection of feasible features. The framework outputs frequency-ranked, modifiable feature suggestions to support transparent, instance-level and population-level insights for attrition mitigation.
    }
    \label{fig:lcsxai_flow}
\end{figure}

\par\noindent\textbf{Contributions.} This study makes the following key contributions:
% \textcolor{black}{improve the presentation of the contributions}
\begin{itemize}
    \item [(I)] We conduct a comprehensive empirical evaluation of predictive modeling approaches—including shallow learners, ensemble methods and deep neural architectures—on employee attrition data from the healthcare sector. Our analysis benchmarks model performance across multiple dimensions of accuracy and interpretability. We further evaluate the framework on real-world employee attrition data related to medical sector.
    
    \item[(II)] We introduce \textbf{LSC-XAI}, a novel hybrid explainability framework (shown in \cref{fig:lcsxai_flow}) that sequentially integrates LIME, SHAP and counterfactual reasoning to enhance the interpretability of black-box models. By intersecting local and global importance measures and filtering through feasibility constraints, LSC-XAI generates \textcolor{blue}{compact, feasibility-filtered feature subsets, with the quantitative status of the framework reported in Section~5.6}. A frequency-based aggregation of counterfactual modifications is further applied to prioritize consistent and actionable intervention strategies.

    \item[(III)] \textcolor{black}{We provide a} \textcolor{blue}{precise positioning of the suggested LSC-XAI framework relative to the state-of-the-art XAI frameworks. Rather than claiming a general performance advantage, we specify what the framework adds: suggestions restricted to features jointly endorsed by two independent explainers (LIME and SHAP) and deemed institutionally feasible, aggregated into frequency-ranked recommendations. Its quantitative status against standalone LIME, SHAP and counterfactual explanations is reported in Section~5.6, and the properties that follow by design rather than by measurement are identified explicitly.}  
\end{itemize}

The structure of this paper is as follows: Section \ref{sec:related_work} reviews relevant literature to contextualize our approach. \textcolor{black}{Section \ref{sec:method} outlines the methodology, discusses the dataset and the preprocessing steps applied to it, along with the machine learning models' description. Section \ref{sec:interpretable_models} features the mathematical framework and model explainability for predicting employee attrition. Section \ref{sec:exp} details the experimental setup, results of the models' training and key findings with interpretative analysis. Section \ref{sec:disc} discusses the results, ethical considerations regarding this framework and outlines the study’s limitations}. Section \ref{sec:conclusion} concludes the paper and proposes future research directions. \textcolor{black}{Appendix provides supplementary technical and experimental visualizations and tables}.

\section{Related Works}
\label{sec:related_work}
\subsection{\textcolor{black}{Statistical analysis}}
This study \cite{Ocean2023} examined the UK NHS (National Health Service) workforce to investigate their work-life balance issues. The study found that satisfaction with leisure time and income has remained lower for healthcare workers than other public sector workers. By 2020, workers with low satisfaction with leisure time were less likely to remain in healthcare. The study also suggested that overall job satisfaction had fallen faster in other areas of the public sector than in healthcare, suggesting wider issues within the UK public sector. The study did not incorporate any prediction of employee attrition using shallow or deep models.

Another study \cite{EJHRMS1220} which examined the factors influencing employee retention in the healthcare sector, delves into the specific context of private healthcare institutions in Malaysia, identifying key determinants such as job satisfaction, work environment, compensation and benefits, work-life balance and employee recognition. However, the author’s study offered valuable insights into the unique challenges faced by private healthcare organizations, particularly in terms of employee recognition and work-life balance. Their methodology relied on traditional statistical methods like regression analysis, correlation analysis and structural equation modeling.

\subsection{\textcolor{black}{Predictive models}}
Rahul et. al. \cite{yedida2018employee} focused on predicting voluntary employee attrition using supervised machine learning techniques for binary classification (leave or not leave). It uses a dataset from Kaggle that includes 14,999 employee records and 10 features, such as average monthly hours worked, number of projects completed, years spent in the company and recent promotions. To predict employee attrition, the authors experimented with several models, including Naive Bayes, Logistic Regression, Multi-Layer Perceptron (MLP) and K-Nearest Neighbors (KNN). Among these, KNN emerged as the most effective, outperforming the other classifiers with 0.9697 in  ROC-AUC, 0.9432 in accuracy and 0.8826 in F1 score. The authors did not explore deep learning architectures beyond the MLP, nor did they incorporate explainability techniques.

Aasheesh et. al. \cite{barvey2018proactive} focused on developing a solution that predicts which employees are likely to leave a company (employee attrition) and suggests actions that managers or the organization can take to prevent this from happening. The IBM Watson dataset was used as the base, but since the data alone didn't fully capture the business needs, feature engineering was applied to make it more useful. The model employed an ensemble classification approach to distinguish between employees who left and stayed. After identifying patterns among those who had already left, the model predicted future cases of attrition, resulting in 91\% accuracy in predicting which employees might leave. Additionally, linear regression was used to estimate how long each employee might stay, helping management plan for those likely to leave. The authors did not incorporate deep models and interpretability in their paper.

Alao et. al. \cite{alao2013analyzing} focused on identifying the factors related to employees that can help predict whether they will leave an organization. The research used personnel records from a higher institution in South-West Nigeria, analyzing data from 309 employees out of 432 to identify key factors like Sex, State of Origin, Length of Service, Rank, Salary and Reasons for leaving. It tested decision tree algorithms, including C4.5 (J48), REPTree and CART (SimpleCart), to predict employee attrition, with C4.5 (J48) performing the best. A pipeline for an Employee Attrition Prediction System was designed, with See5 Decision Tree showing the best overall accuracy. In this work, the authors explored only decision tree-based algorithms, not exploring other training techniques and explainability. 

Sandeep et. al. \cite{Yadav2018} examined employee data from a Kaggle dataset containing 14,999 records with 12 features related to employees' professional lives. The paper conducted experiments using four approaches: brute force (Logistic  Regression, SVM, Random Forest, Decision Tree, AdaBoost), one-hot encoding (Logistic  Regression, SVM, Random Forest, Decision Tree, AdaBoost), feature selection (Random Forest, AdaBoost) and experienced employee data (Logistic  Regression, SVM, Random Forest, Decision Tree, AdaBoost) and evaluated their performance using accuracy, precision, recall and F1-score metrics. Besides, both Random Forest and AdaBoost were trained with updated features on feature selection approaches and experienced employee data. While exploring various preprocessing steps and models, they lacked the incorporation of deep learning and interpretability of their model’s findings.

In this work, Shobhanam et. al. \cite{Krishna2022} evaluated an HR dataset from Kaggle through feature engineering and data visualization. They then used machine learning techniques such as Random Forest and AdaBoost and compared their results, finding that Random Forest performed better than the other. ROC score, specificity, recall, CV score, accuracy and precision were used to assess the models. The authors did not address explainability or the use of alternative learning strategies; instead, they solely used ensemble models to complete their prediction task.

Jungryeol et. al. \cite{Park2024} focused on predicting the employees who were eager to change their jobs or possess a willingness towards it. They used data from the Korea Employment Information Service's Job Movement Path Survey for college graduates. To employ their machine learning for classification, they used Logistic Regression (LR), k-nearest neighbor (KNN) and extreme gradient boosting (XGB). To confirm the influence of predictors, they use OLS regression. XGB shows the highest accuracy rate with 78.5\%. From their findings, job security and organizational satisfaction were determined to be the most important factors. They also found reversed influence on some traditional factors like workload importance and the relevance of the job to one's educational field. The authors only explored limited machine learning models for their study.

Nesreen et. al. \cite{ElRayes2020} analyzed a dataset of 5,550 employee job transitions between 2007 and 2016, sourced from an extensive proprietary database of resumes shared anonymously through Glassdoor's platform, to gain insights into employee job transition decisions and develop predictive models. The study showed approximately three-quarters of employees leave their employer during a job transition and that salary increases are a primary motivation for job transitions. Tree-based models, particularly the light gradient boosted tree classifier, outperform other models in predicting employee attrition with an AUC value of 75. The study only focused on tree-based models and did not comprise exploring the interpretability of the models.

Chung et. al. \cite{CHUNG2023119364} worked on employee attrition prediction, which emphasizes the importance of human resources and the challenges associated with predicting employee turnover. Traditional methods like regression and correlation analysis have been used to identify factors contributing to attrition, but these often lack accuracy and struggle with complex data patterns. This study proposes a stacking ensemble model using Logistic Regression, Random Forest, XGBoost, SVM and ANN, demonstrating improved performance compared to single models. The ensemble model's ability to handle complex relationships and improve prediction accuracy underscores its potential in employee attrition prediction. The research did not comprise exploring the interpretability of the applied models.

Raza et. al. \cite{app12136424} addressed the critical issue of employee attrition by analyzing its causes and proposing a predictive framework using machine learning. Recognizing the substantial financial impact of attrition, with the Society for Human Resource Management (SHRM) estimating an average cost-per-hire of USD 4129, the study identifies key factors such as monthly income, hourly rate, job level and age through Employee Exploratory Data Analysis (EEDA). To predict attrition, four machine learning techniques were employed: Logistic Regression (LR), Random Forest (RF), Gradient Boosting (GB) and the proposed optimized Extra Trees Classifier (ETC). Among these, the ETC achieved the highest accuracy of 93\%, outperforming other methods and recent state-of-the-art studies. The authors did not employ the explainability of their performing models.

Jain et. al. \cite{Jain2020} highlighted the growing importance of machine learning techniques in human resource management (HRM), drawing from various scholarly works that explore the relationship between HRM and organizational performance. The authors noted that machine learning techniques are increasingly useful in prediction across various domains, including source code changes, crop suitability and productivity prediction. Specifically in employee attrition, previous studies have employed various machine learning classification techniques such as Support Vector Machines (SVM), Decision Trees (DT), Logistic Regression and Neural Networks, with different studies finding varying levels of predictive accuracy depending on the specific approach and dataset used. The authors did not incorporate ensemble machine learning techniques and did not introduce explainability to their models’ predictions.

\subsection{\textcolor{black}{XAI-Based framework}}
Mariam et. al. \cite{Akasheh2024} proposed a novel framework that leverages knowledge graphs and explainable AI techniques to predict employee turnover. For the study, dataset from IBM analytics consisting of approximately 1500 samples and 35 features was used. Conventional ML techniques such as LR, RF, LGBM, GB and L-SVM are studied and explored. First, with the default setting, LR shows the best output with an accuracy of 0.88. Later, by transforming conventional tabular data into a knowledge graph format and employing graph convolutional networks, L-SVM demonstrated the best results in forecasting turnover with 0.925 in accuracy. While exploring the interpretability of their best model, the authors provided only local instance-level explanations and lacked global interpretability for a more comprehensive overview. 

Gabriel et. al. \cite{Diaz2023} explored the use of explainable AI (XAI) to predict employee turnover and develop data-driven strategies to deal with this challenging issue. They used IBM HR open dataset consisting of 1471 samples. They perform EDA to process their data, apply SMOTE-Tomek process to reduce imbalancing of the data and then train ML models such as XGBoost, RF, KNN, SVM, DT, LR and Gaussian NB. Among these, XGBoost showed the best results with an accuracy of 83\% in the balanced dataset. After getting their final model, interpretability is applied, both locally and globally, which includes techniques like PDP, LIME, ELI5 and SHAP. The study did not include deep models for prediction tasks and corrective measures from the found interpretability.

Karthik et. al. \cite{Karthik2022} employed LIME and SHAP models to explain the interoperability and explainability of the ML models used for forecasting employee attrition. The IBM HR Analytics Attrition dataset is preprocessed using Label Encoder to convert categorical inputs into numerical vectors. The LightGBM model is used as the baseline model and the preprocessed dataset is divided into training and testing sets using stratified cross-validation. The study achieves an AUC score of 0.837 on the test set using the LightGBM model. The authors only explored the local interpretability of the model and did not delve deeper into the corrective factors to strengthen employee retention.

\textcolor{blue}{Beyond the attrition domain, our framework builds on---and is positioned relative to---the constrained counterfactual-recourse literature. Wachter et al.~\cite{wachter2017counterfactual} established counterfactuals as model-agnostic explanations; Ustun et al.~\cite{ustun2019actionable} formalized actionable recourse as an optimization problem under mutability constraints for linear classifiers; Karimi et al.~\cite{karimi2021algorithmic} extended recourse from explanations to causal interventions; and Mothilal et al.~\cite{mothilal2020dice} provide the DiCE library, which already supports diversity objectives and user-specified restrictions on which features may vary. We state plainly what LSC-XAI does and does not take from this literature: it does not claim to invent constrained recourse, mutability constraints or restricted feature variation---all are prior art, and our counterfactual engine is DiCE itself (Section~\ref{sec:method}). What it contributes is (i) a domain-operationalized feasibility judgement that extends the standard mutability notion with organizational controllability and ethical permissibility, encoded as auditable per-feature assignments (\cref{tab:feasibility_complete}); (ii) an agreement gate under which counterfactual generation is restricted to the features jointly endorsed by two independent explainers (LIME and SHAP); and (iii) an aggregation layer that pools per-instance counterfactuals into frequency-ranked, population-level recommendation sets rather than returning isolated per-instance counterfactuals. The evaluation of these claims---explanation-quality metrics with confidence intervals, the Saudi cross-domain transfer protocol and the subgroup fairness audit---constitutes the empirical contribution reported in Sections~5.6--5.8.}

\subsection{Research Gap}
Despite significant advancements in employee attrition prediction, several critical gaps remain unaddressed in the existing body of literature:

\begin{enumerate}
    \item \textbf{Lack of integrated model comparison:} Existing studies predominantly focus on either shallow models (e.g., Logistic Regression, decision trees), ensemble models (e.g., random forest, gradient boosting), or deep learning models (e.g., neural networks) in isolation. There is a scarcity of comparative frameworks that simultaneously evaluate all three categories to holistically assess performance in real-world contexts.

    \item \textbf{Neglect of healthcare domain-specific analysis:} The majority of research on attrition prediction has been conducted in corporate or general industrial contexts. Very few studies have focused on the healthcare sector, which has unique stressors, job dynamics and retention challenges. Also, none of the prior studies evaluated their proposed frameworks using real-world employee attrition in healthcare data. 

    \item \textbf{Limited emphasis on actionable interpretability:} While some research employs local explainability techniques (e.g., LIME, SHAP), global interpretability—necessary for high-level strategic decision-making—remains largely underexplored in this domain.

    \item \textbf{Insufficient intervention strategies:} Prior models often focus on prediction accuracy without translating their findings into actionable interventions. There is a gap in developing frameworks that bridge the prediction layer with human-centric retention strategies.
    \item \textbf{Constrained recourse without domain grounding:} \textcolor{blue}{while counterfactual libraries and the actionable-recourse literature can restrict which features may vary, no prior attrition study grounds that restriction in institution-specific controllability and ethical-permissibility criteria with an auditable per-feature record, gates counterfactual generation on agreement between two independent explainers, or aggregates per-instance counterfactuals into frequency-ranked population-level recommendations for human decision-making.}
\end{enumerate}

\newpage
\section{Methodology}
\label{sec:method}
This section covers the acquisition of a healthcare employee attrition dataset and details the preprocessing steps applied. It then describes various shallow, ensemble and deep models trained for this research, along with their evaluation criteria. The methodology diagram is shown in figure \ref{fig:workflow}.

% \vspace{0.5cm}
% \pagebreak[2] % Suggests a page break here if it's REALLY tight, but not forced
% \noindent
% \begin{minipage}{\linewidth}
%     \centering
%     \includegraphics[width=0.95\textwidth, height=0.65\textheight]{Attrition_Methodology-updated.png}
%     \captionof{figure}{Overview of the research methodology. The process begins with defining research objectives, followed by data preprocessing for model training. The best-performing model is subsequently analyzed using a hybrid explainability framework to extract actionable insights.}
%     \label{fig:workflow}
% \end{minipage}
\begin{figure}[pos=H]
    \centering
    \includegraphics[width=0.95\textwidth, height=0.60\textheight, keepaspectratio]{Attrition_Methodology-updated.png}
    \caption{Overview of the research methodology. The process begins with defining research objectives, followed by data preprocessing for model training. The best-performing model is subsequently analyzed using a hybrid explainability framework to extract actionable insights.}
    \label{fig:workflow}
\end{figure}
\newpage

\subsection{Data Acquisition}
In this study, the Employee Attrition for Healthcare dataset has been used, which is sourced directly from Kaggle \cite{KaggleHealthcareAttrition}. 
It should be noted that this dataset is synthetically derived from the IBM Watson HR Analytics dataset \cite{ibm_hr_attrition_2017} and does not represent records of real healthcare employees. Findings from the primary experiment should therefore be interpreted as a proof-of-concept demonstration of the LSC-XAI framework's methodological capabilities, with real-world applicability examined in \cref{real-world}. This dataset \cite{KaggleHealthcareAttrition} comprises 1676 entries, with \textcolor{blue}{34 raw attributes} and a target column, 'Attrition' which signifies if an employee has left the organization. \textcolor{blue}{These 34 raw attributes expand to 39 one-hot encoded candidate features after encoding (Section~3.2.4). Throughout the paper, feature counts therefore refer to one of three spaces, and are labelled accordingly: the 34 raw attributes, the 39 encoded candidate features, and the 23 features retained by the fold-wise selection in the Fold-1 model used for the interpretability analysis (Section~3.2.4).} Key features include demographic details (age, gender, marital status), job specifics (role, department, years with the company) and satisfaction, performance and work-life balance metrics. Additional attributes like monthly income, training frequency, among others offer a holistic view of the workforce, aiding in attrition risk analysis.

\subsubsection*{Real-World Validation Dataset}
To validate the predictive generalisability, the Saudi Employee Attrition Dataset \cite{alqahtani2025saudi} a survey-based dataset of 1,200 private-sector employees across Saudi Arabia was employed. As the dataset's feature set did not fully align with the primary dataset, only common columns were retained, yielding a reduced but consistent feature space for cross-dataset evaluation. Medical-sector records were subsequently extracted as the target sub-population, yielding 304 data points. Three experimental setups of progressive distributional alignment were evaluated, detailed in \cref{real-world}.\\

\subsection{Data Preprocessing}
This subsection includes visual analysis, encoding of the features, various types of feature selection methods, handling of the outliers to ensure better preprocessing for the training of the models for this study.

\subsubsection{EDA and Visualization}
\textcolor{black}{
The dataset \cite{KaggleHealthcareAttrition} comprises no missing values. The target variable is quite imbalanced: 11.87\% (n=199) of the instances represent attrition, while 88.13\% (n=1,477) represent retention. Preliminary exploratory analysis of the key features revealed the following trends:}

\textcolor{black}{Job Metrics: 'JobLevel' (ranked 1–5) showed an inverse relationship with attrition; employees at lower levels exhibited a higher propensity to leave. Regarding 'JobRole', nurses and those categorized as "others" were more frequently associated with turnover.}

\textcolor{black}{Demographics: While all 'MaritalStatus' categories (single, married, divorced) reported attrition, it was most prevalent among single employees.}

\textcolor{black}{Professional History: Most employees had worked at at least one prior company ('NumCompaniesWorked'), with multi-company experience correlating with a higher likelihood of switching.}

\textcolor{black}{Work Conditions: Although the majority did not work 'OverTime', those who did were notably more likely to leave. Conversely, employees with fewer 'Shift' assignments showed higher retention rates.}

\textcolor{black}{Tenure: Both 'TotalWorkingYears' and 'YearsAtCompany' were strong predictors; attrition was concentrated among those with 1–10 years of experience. Furthermore, employees in the early stages of their 'YearsInCurrentRole' demonstrated a higher propensity for turnover.}

\begin{figure}[pos=H]
    \centering
        \begin{subfigure}{0.3\textwidth}
            \centering
            \includegraphics[width=\linewidth, height=100pt]{distribution_job_level.png}
            \caption{Job level}
        \end{subfigure} 
        \begin{subfigure}{0.3\textwidth}
            \centering
            \includegraphics[width=\linewidth, height=100pt]{distribution_job_role.png}
            \caption{Job role}
        \end{subfigure} 
        \begin{subfigure}{0.3\textwidth}
            \centering
            \includegraphics[width=\linewidth, height=100pt]{distribution_marital_status.png}
            \caption{Marital status}
        \end{subfigure} \\
        \begin{subfigure}{0.3\textwidth}
            \centering
            \includegraphics[width=\linewidth, height=100pt]{distribution_num_of_companies_worked.png}
            \caption{Number of companies worked}
        \end{subfigure} 
        \begin{subfigure}{0.3\textwidth}
            \centering
            \includegraphics[width=\linewidth, height=100pt]{distribution_overtime.png}
            \caption{Over time}
        \end{subfigure} 
        \begin{subfigure}{0.3\textwidth}
            \centering
            \includegraphics[width=\linewidth, height=100pt]{distribution_shift.png}
            \caption{Shift}
        \end{subfigure} \\
        \begin{subfigure}{0.3\textwidth}
            \centering
            \includegraphics[width=\linewidth, height=100pt]{distribution_total_working_years.png}
            \caption{Total working years}
        \end{subfigure} 
        \begin{subfigure}{0.3\textwidth}
            \centering
            \includegraphics[width=\linewidth, height=100pt]{distribution_years_at_company.png}
            \caption{Years at company}
        \end{subfigure} 
        \begin{subfigure}{0.3\textwidth}
            \centering
            \includegraphics[width=\linewidth, height=100pt]{distribution_years_in_current_role.png}
            \caption{Years in current role}
        \end{subfigure}\\
    \caption{Distributions of various features are displayed, revealing a significant imbalance with 'No Attrition' representing the majority class.}
    \label{fig:grid}
\end{figure}

% \begin{figure}[pos=H]
%     \centering
%         \begin{subfigure}{0.18\textwidth}
%             \centering
%             \includegraphics[width=\linewidth, height=73pt]{distribution_job_level.png}
%             \caption{Job level}
%         \end{subfigure} 
%         \begin{subfigure}{0.18\textwidth}
%             \centering
%             \includegraphics[width=\linewidth, height=73pt]{distribution_job_role.png}
%             \caption{Job role}
%         \end{subfigure} 
%         \begin{subfigure}{0.18\textwidth}
%             \centering
%             \includegraphics[width=\linewidth, height=73pt{distribution_marital_status.png}
%             \caption{Marital status}
%         \end{subfigure} 
%         \begin{subfigure}{0.18\textwidth}
%             \centering
%             \includegraphics[width=\linewidth, height=73pt]{distribution_overtime.png}
%             \caption{Over time}
%         \end{subfigure} 
%         \begin{subfigure}{0.18\textwidth}
%             \centering
%             \includegraphics[width=\linewidth, height=73pt]{distribution_shift.png}
%             \caption{Shift}
%         \end{subfigure} \\
%         \begin{subfigure}{0.25\textwidth}
%             \centering
%             \includegraphics[width=\linewidth, height=73pt]{distribution_num_of_companies_worked.png}
%             \caption{Num. of companies worked}
%         \end{subfigure} 
%         \begin{subfigure}{0.20\textwidth}
%             \centering
%             \includegraphics[width=\linewidth, height=73pt]{distribution_total_working_years.png}
%             \caption{Total working years}
%         \end{subfigure} 
%         \begin{subfigure}{0.20\textwidth}
%             \centering
%             \includegraphics[width=\linewidth, height=73pt]{distribution_years_at_company.png}
%             \caption{Years at company}
%         \end{subfigure} 
%         \begin{subfigure}{0.20\textwidth}
%             \centering
%             \includegraphics[width=\linewidth, height=73pt]{distribution_years_in_current_role.png}
%             \caption{Years in current role}
%         \end{subfigure}
%     \caption{Distributions of various features are displayed, revealing a significant imbalance with 'No Attrition' representing the majority class.}
%     \label{fig:grid}
% \end{figure}

\subsubsection{Data Encoding}
\textcolor{black}{Data encoding in machine learning preprocessing constitutes a critical step that can influence model performance. The dataset under consideration comprised both numerical and categorical values. Among the categorical types, both nominal and ordinal variables are present. Nominal variables are addressed using one-hot encoding, a widely employed technique in machine learning for representing categorical variables as binary vectors \cite{Bishop2006}. Ordinal variables, along with some nominal ones that are represented textually, have been manually encoded to ensure that their original order is preserved during the transformation from categorical to numerical format.}

\begin{table}[width=1\linewidth,cols=2,pos=h]
    \centering
    % \vspace{-0.3cm}
    \caption{Attributes 'BusinessTravel', 'Gender' and 'Overtime' are manually transformed and their values are shown.}
    \renewcommand{\arraystretch}{1.5} % Adjust row height
    \label{tab:sample}
    \begin{tabular}{C|C} % Defines a 3-column table with vertical borders
        \hline
        \textbf{Attribute} & \textbf{Manually transformed values} \\ \hline
                           & 'Non\_Travel': 0            \\
         BusinessTravel    & 'Travel\_Rarely': 1         \\
                           & 'Travel\_Frequently': 2     \\ \hline
        Gender             & 'Male': 1, 'Female': 0     \\ \hline
        Overtime           & 'Yes': 1, 'No': 0          \\ \hline
    \end{tabular}
\end{table}

\subsubsection{\textcolor{black}{Feature Engineering}}
\textcolor{black}{To optimize model training efficiency and mitigate the risk of overfitting, we employed a multi-stage feature selection pipeline utilizing Pearson Correlation, Mutual Information (MI), Chi-squared tests and Recursive Feature Elimination with Cross-Validation (RFECV). Initially, Pearson correlation was applied to identify multicollinearity. Following standard practices for high-dimensional stability \cite{Pearson1895}, a threshold of $|r| > 0.8$ was established for redundancy (see Fig. \ref{fig:pearson}). For instance, although 'MonthlyIncome' and 'JobLevel' exhibited a high correlation ($r = 0.95$), \textcolor{blue}{only 'MonthlyIncome' was retained, because the implemented rule discards the member of a highly correlated pair with the lower variance; 'JobLevel' is accordingly absent from the 23 features of the Fold-1 model used for the interpretability analysis (Section~3.2.4)}. Subsequently, a "majority-vote" pruning strategy was implemented: features identified as non-significant by at least two of the three remaining techniques (MI, Chi-squared and RFECV) were removed from the final feature set.}\\
Following the initial correlation analysis, the remaining features were evaluated using Mutual Information (MI), Chi-squared ($\chi^2$) tests and Recursive Feature Elimination with Cross-Validation (RFECV) to ensure the inclusion of both linear and non-linear dependencies \cite{CoverThomas1991, Agresti2018}. MI was utilized to rank features by informational gain, while $\chi^2$ tests assessed statistical dependence between categorical attributes and attrition. To prevent selection bias and improve model generalizability, RFECV with stratified cross-validation was employed \cite{Guyon2002, Kohavi1995}.

% \textcolor{black}{The quantitative distributions for feature correlations, Mutual Information rankings, Chi-squared scores and RFECV are visualized in figures \ref{fig:pearson}, \ref{fig:mutual_info}, \ref{fig:chi2}, \ref{fig:rfecv} of the Appendix. These results provided the statistical basis for the following pruning decisions.}

\subsubsection{Feature Elimination}
% Features were permanently discarded only if they were identified as non-significant by at least two of these three independent methodologies, ensuring a robust and empirically justified feature set. In accordance with this: 'Gender', 'Department\_Maternity', 'EducationField\_Medical', 'EducationField\_TechnicalDegree', 'RelationshipSatisfaction', 'PerformanceRating', 'Education' and 'EducationField\_Other' were removed.
\textcolor{black}{To prevent information leakage, mutual information, chi-squared testing and RFECV were performed independently within each training fold of the 5-fold stratified cross-validation described in Section 3.3, rather than once on the complete dataset. Pearson correlation-based collinearity pruning, which does not use the target variable, was likewise computed per training fold for consistency, though its target-independence means leakage risk from this step alone is minimal. Across the five folds, the number of features removed by the majority-vote pruning procedure ranged from \textcolor{blue}{8 to 21 (mean = 14.4, with per-fold counts of 15, 15, 21, 13 and 8)} out of a pool of 39 candidate features\textcolor{blue}{, i.e., the 34 raw attributes of Section~3.1 expanded by one-hot encoding (Section~3.2.3)}. Table \ref{tab:removed_features} in appendix reports the per-feature removal frequency across folds. \textcolor{blue}{38.5\%} of features were retained in every fold, \textcolor{blue}{10.3\%} were removed in every fold and the remaining features showed fold-dependent selection status, with a mean pairwise Jaccard similarity of \textcolor{blue}{0.45} across the five folds' removed-feature sets. \textcolor{blue}{In the Fold-1 training fold --- the fold used throughout the interpretability analysis of Sections~5.4--5.6 --- one further feature was removed by the Pearson collinearity screen ('JobLevel', correlated with 'MonthlyIncome' at $r = 0.95$) and 15 by the majority-vote procedure, leaving 23 features. This 23-feature set is the feature space of the interpretability model and the denominator of the sparsity metric reported in Table~10, where $1 - 12/23 = 0.4783$ for LIME and SHAP and $1 - 3/23 = 0.8696$ for LSC-XAI.} This variability confirms that several features flagged as non-significant under a single global selection pass are in fact borderline and supports the necessity of fold-wise selection to avoid overstating the reliability of the reported feature set.}
% \begin{figure}[pos=H]
%     \centering
%     \includegraphics[width=0.7\linewidth]{common_feature_eliminated.png}
%     \caption{Less significant features indicated by two or more feature selection methods are shown. These features were removed for the proceeding steps.}
%     \label{fig:enter-label}
% \end{figure}

\subsubsection{\textcolor{black}{Distributional Transformation}}
To mitigate the influence of extreme values and correct for non-normality—common issues in attrition datasets—we applied the Box-Cox transformation \cite{Box1964}. This approach stabilizes variance and ensures the data more closely approximates a Gaussian distribution, thereby enhancing the robustness of subsequent statistical modeling. The transformation parameter $\lambda$ was optimized using Maximum Likelihood Estimation (MLE) to determine the most effective power transformation for each feature. \textcolor{black}{$\lambda$ was estimated using only each of the 5-fold's training data and applied without refitting to the corresponding held-out test fold}. \textcolor{blue}{The transformation was applied only to continuous features (e.g., Age, DailyRate, HourlyRate, MonthlyIncome, MonthlyRate, PercentSalaryHike, DistanceFromHome and the tenure measures TotalWorkingYears, YearsAtCompany, YearsInCurrentRole, YearsSinceLastPromotion and YearsWithCurrManager); binary and categorical attributes, including all one-hot encoded variables, were excluded. Because Box-Cox requires strictly positive inputs, exact zeros — which occur in the tenure variables (e.g., YearsAtCompany = 0 for recent hires; YearsInCurrentRole = 0 for employees recently moved into their role; YearsSinceLastPromotion = 0 for employees recently promoted; YearsWithCurrManager = 0 for employees recently assigned a manager) — were replaced with a small fixed constant (0 $\rightarrow$ 0.1); this substitution is a fixed, non-data-derived rule applied identically to the training and held-out folds. No feature contained negative values, so no additional shift constant was required; the implementation nevertheless includes a guard that would add a constant of $|min| + 1$ before fitting if any feature contained non-positive values, applied unchanged to the held-out fold.} The resulting shifts in distribution are illustrated in \cref{fig:boxtrans} of the Appendix. \textcolor{black}{No records were removed as a result of this procedure; Box-Cox modifies the distributional shape of the affected features without excluding observations.}

\subsubsection{\textcolor{black}{Addressing Class Imbalance}}
\textcolor{black}{To address potential model bias arising from the imbalanced class distribution, we employed Synthetic Minority Over-sampling Technique for Nominal and Continuous (SMOTE-NC) on the training set. Unlike standard SMOTE, SMOTE-NC effectively handles the dataset's hybrid feature space by interpolating synthetic instances while preserving the statistical structure of categorical attributes \cite{Chawla2002}. This approach ensures that the minority class representation is enhanced without distorting the underlying relationships between nominal and continuous variables \cite{Batista2004}. \textcolor{blue}{Within the cross-validation experiments, SMOTE-NC was applied to each training fold only, after fold-wise feature selection and Box-Cox transformation and before Z-score normalization (Section~3.3).}}

\subsubsection{\textcolor{black}{Data Normalization}}
\textcolor{black}{Prior to model training, the feature set was standardized using Z-score normalization. This procedure rescales the data to a zero-mean and unit-variance distribution, ensuring that features with larger numerical ranges do not disproportionately influence the model's loss function \cite{Iglewicz1993}. The transformation is defined as:\begin{equation}Z = \frac{X-\mu}{\sigma} \label{eq:eq5}\end{equation}where $X$, $\mu$ and $\sigma$ denote the original data point, the feature mean and the standard deviation, respectively. \textcolor{blue}{In the cross-validation experiments, the continuous features were standardized using parameters ($\mu$ and $\sigma$) estimated exclusively on the SMOTE-resampled training folds and applied, without refitting, to the corresponding held-out folds (Section~3.3).}}

\textcolor{black}{\subsection{Data-splitting for Training and Testing Set}}
\textcolor{black}{To split the data into training and testing sets, stratified K-Fold Cross Validation \cite{Kohavi1995} was applied. It ensures minority classes are properly distributed between the training and the testing sets. The value of K was set to 5. For the training set 4 folds were used and 1 fold was used for the testing set, making a 80\%-20\% split. Data distribution for training and testing set across the folds are shown in table \ref{tab:smoted}. In each iteration, \textcolor{blue}{every preprocessing transform was executed strictly inside the fold, so that no transform's parameters were ever estimated on the held-out fold, in the following order: (1) the majority-vote feature selection of Section~3.2.4, including the target-independent Pearson collinearity screen, was computed using the four training folds only, and the retained feature set was then applied to the held-out fold; (2) the Box-Cox parameters ($\lambda$) were estimated by MLE on the four training folds only and applied, without refitting, to the held-out fold (Section~3.2.2); (3) SMOTE-NC was applied to the four training folds only, after feature selection and Box-Cox, so that synthetic minority instances were generated in the selected, transformed feature space and the held-out fold remained untouched by resampling; (4) Z-score normalization was applied to the continuous features, with the feature means ($\mu$) and standard deviations ($\sigma$) estimated on the SMOTE-resampled training data only and then applied, without refitting, to the held-out fold; and (5) the model was trained on the resampled and standardized training data and evaluated on the untouched held-out fold.}}

\begin{table}[width=0.55\linewidth,cols=5,pos=h]
    \centering
    \caption{\textcolor{black}{Data quantities for the training and testing sets are showcased. The number of data denoted under 'After Upsampling' were used for training in the mentioned iteration}}
    \renewcommand{\arraystretch}{1.25} % Adjust row height
    \begin{tabular*}{0.55\linewidth}{C|C|C|C|C}
        \cline{1-5} 
        \textbf{\textcolor{black}{Iteration}} & \textbf{\textcolor{black}{Attrition State}} & \textbf{\textcolor{black}{Initial Count}} & \textbf{\textcolor{black}{After Upsampling}} & \textbf{\textcolor{black}{Testing}}\\ \cline{1-5}
        
        \multirow{2}{*}{\textcolor{black}{1}} & \textcolor{black}{Yes} & \textcolor{black}{159} & \textcolor{black}{1181} & \textcolor{black}{40}\\ \cline{2-5} 
                            & \textcolor{black}{No}  & \textcolor{black}{1181} & \textcolor{black}{1181} & \textcolor{black}{296}\\ \cline{1-5} 
        \multirow{2}{*}{\textcolor{black}{2}} & \textcolor{black}{Yes} & \textcolor{black}{160} & \textcolor{black}{1181} & \textcolor{black}{39}\\ \cline{2-5} 
                            & \textcolor{black}{No}  & \textcolor{black}{1181} & \textcolor{black}{1181} & \textcolor{black}{296}\\ \cline{1-5} 
        \multirow{2}{*}{\textcolor{black}{3}} & \textcolor{black}{Yes} & \textcolor{black}{159} & \textcolor{black}{1182} & \textcolor{black}{40}\\ \cline{2-5} 
                            & \textcolor{black}{No}  & \textcolor{black}{1182} & \textcolor{black}{1182} & \textcolor{black}{295}\\ \cline{1-5}  
        \multirow{2}{*}{\textcolor{black}{4}} & \textcolor{black}{Yes} & \textcolor{black}{159} & \textcolor{black}{1182} & \textcolor{black}{40}\\ \cline{2-5} 
                            & \textcolor{black}{No}  & \textcolor{black}{1182} & \textcolor{black}{1182} & \textcolor{black}{295}\\ \cline{1-5}
        \multirow{2}{*}{\textcolor{black}{5}} & \textcolor{black}{Yes} & \textcolor{black}{159} & \textcolor{black}{1182} & \textcolor{black}{40}\\ \cline{2-5} 
                            & \textcolor{black}{No}  & \textcolor{black}{1182} & \textcolor{black}{1182} & \textcolor{black}{295}\\ \cline{1-5}
    \end{tabular*}
    \label{tab:smoted}
\end{table} 

\subsection{\textcolor{black}{Predictive Modeling Framework}}
\textcolor{black}{We evaluated several predictive architectures, ranging from baseline statistical models to advanced ensemble and deep learning techniques, to identify the most robust approach for predicting healthcare worker attrition.}

\subsubsection{\textcolor{black}{Baseline and Shallow Learners}}
\textcolor{black}{We implemented Logistic Regression (LR) \cite{logisRegr}, Support Vector Machines (SVM) \cite{cortes1995support}, K-Nearest Neighbors (KNN) \cite{fix1951knn} and Decision Trees (DT) \cite{breiman1984classification} as baseline classifiers. LR provided a linear benchmark for attrition prediction, while SVM was used for its ability to identify separating hyperplanes in high-dimensional feature spaces. For KNN, we used the library-default configuration, including the default number of neighbors and distance metric, to maintain consistency with the baseline-comparison setting. DTs served as an interpretable tree-based baseline and also provided the conceptual foundation for subsequent ensemble learning methods.}

\subsubsection{\textcolor{black}{Ensemble Learning Architectures}}
\textcolor{black}{To enhance predictive accuracy and mitigate overfitting, we employed several boosting and bagging variants:
\textbf{Bagging-based}: Random Forest (RF) was implemented using bootstrap aggregation and random feature selection to improve model diversity \cite{breiman2001random}.\\
\textbf{Boosting-based}: We compared Gradient Boosting (GB) \cite{friedman2001greedy}, AdaBoost \cite{freund1997decision}, XGBoost \cite{chen2016xgboost}, LightGBM \cite{ke2017lightgbm} and CatBoost \cite{prokhorenkova2018catboost}. Notably, LightGBM was selected for its leaf-wise growth strategy to handle large-scale data efficiently, while CatBoost was utilized for its improved handling of the categorical features present in our healthcare dataset without the need for extensive manual encoding.\\
\textbf{Stacking Ensemble}: A multi-level stacking classifier was constructed \cite{dzeroski2004stacking}. Based on individual performance, SVM, RF and LightGBM were integrated as base learners, with Logistic Regression serving as the meta-model to aggregate predictions due to its low computational overhead and interpretability.} 

\subsubsection{\textcolor{black}{Deep Learning for Tabular Data}}
\textcolor{black}{For non-linear feature extraction, we implemented three specialized architectures:\\
\textbf{Multi-layer Perceptron (MLP)}: The MLP architecture was manually configured using fully connected layers, ReLU activation, batch normalization, sigmoid output activation, binary cross-entropy loss and the Adam optimizer. A 10\% validation split from the training data was used to monitor training behavior. The final architecture and training configuration are reported in \cref{tab:software_implementations} and illustrated in \cref{fig:mlp}.\\
\textbf{1D-Convolutional Neural Networks (1D-CNN)}: Although typically used for sequential data, 1D-CNNs were employed to capture local patterns across the feature vectors \cite{yin2017survey}. Studies by Borisov et al. \cite{borisov2022} and Baosenguo \cite{baosenguo2021} have shown that 1D convolutional operations can extract meaningful local patterns from tabular feature representations when features are appropriately ordered or grouped. \textcolor{black}{In this study, the default arbitrary feature ordering was maintained without grouping, which is noted as a limitation. Our 1D-CNN model was manually configured using one convolutional layer, ReLU activation, max pooling, a fully connected output layer, sigmoid activation, binary cross-entropy loss and the Adam optimizer. The final architecture and training configuration are reported in \cref{tab:software_implementations} and illustrated in \cref{fig:cnn}.}\\
\textbf{TabNet}: To bridge the gap between tree-based models and deep learning, we utilized TabNet, which applies a sequential attention mechanism for sparse feature selection \cite{arik2020tabnet}. The TabNet model was trained using the default TabNetClassifier configuration, with the maximum number of epochs and early-stopping patience reported in \cref{tab:software_implementations}.}


\subsection{Model Evaluation}
\textcolor{black}{The predictive performance of the models was evaluated using a suite of standard classification metrics, including Accuracy, Precision, Recall and F1-score. For further evaluation, Wilcoxon signed-rank test \cite{wilcoxon1945} and model complexity was also considered.} 

\section{Interpretable Models}
\label{sec:interpretable_models}
The importance of trust and transparency in artificial intelligence systems' decision-making has raised demand for transparent machine learning models. When it comes to these models, it is critical to comprehend and analyze the logic behind their choices. Interpretability refers to the ability to comprehend the mechanisms and reasoning underlying an ML model's conclusions, offering a visual, quantitative and/or qualitative explanation of the factors that affect predictions. Because of their intrinsic design, 'white-box' models are capable of establishing unambiguous relationships between the variables and the output. In contrast, 'black-box' models are considerably more difficult to interpret due to their intricate internal structures \cite{prabhu2023interpretation}. Importantly, even in the domain of white-box models (which are often referred to as interpretable), interpretability is still debated \cite{lipton2018mythos}.\\ 
\textbf{LIME:} One of the models that addresses the complex nature of interpreting in a well-established way is LIME (Local Interpretable Model Agnostic Explainer)\cite{prabhu2023interpretation}. LIME analyzes the outcomes to determine how the predictions alter when perturbations are added to the input data. LIME generates new datasets with these variances to obtain sets of predictions. Essentially, the explanations are based on local surrogate models. Surrogate models, such as decision trees or linear models, are interpretable models that gain knowledge of the predictions made by the underlying black-box model. Instead of trying to fit a global surrogate model, LIME concentrates on fitting local surrogate models to explain why the predictions were made independently \cite{ribeiro2016should}.\\
\textbf{SHAP:} Another well-known technique for deciphering black-box models is Shapely Additive Explanations, or SHAP. The goal of the SHAP interpretability model is to explain a case based on how each feature contributed to the prediction \cite{lundberg2017unified}. SHAP values try to take into consideration the result of a function as the sum of the impacts of every characteristic that is conditionally entered.\\
\textbf{Counterfactual:} The counterfactual approach is a method for determining key elements that determine the model's prediction. Counterfactuals are scenarios created by making only minor changes to input characteristics in order to see how they affect the output. Their goal is to provide end users with clear, actionable explanations \cite{wang2020counterfactual}. In areas such as finance, healthcare and criminal justice, counterfactuals can help stakeholders understand why a particular decision was made and what changes could possibly lead to a different outcome. This can help in addressing biases and ensuring fairness \cite{caruana2021counterfactual}. Despite their utility, counterfactual explanations grapple with challenges like computational efficiency, ensuring relevance and maintaining fidelity to the original model \cite{guidotti2018survey}.
\newline

\subsection{Proposed LSC-XAI: A Hybrid Explainability Framework}
\label{sec:lsc_xai}

% \textcolor{black}{Highlight this with better transitions}\newline
\textcolor{black}{Through this research, we propose a hybrid interpretability framework}, \textbf{LIME–SHAP–Counterfactual Explainable AI (LSC-XAI)}, which integrates local explanations, global attributions and actionable counterfactual reasoning in a sequential and coherent manner. This technique addresses the limitations inherent in those individual XAI techniques. \textcolor{black}{The stand-alone XAI frameworks are described in \cref{sec:interpretable_models}.}

Each component contributes uniquely: LIME provides local, instance-specific explanations but may suffer from instability across different perturbations and samples. SHAP delivers globally consistent attributions based on Shapley values but can overlook localized predictive nuances. Counterfactual explanations offer actionable insights by suggesting minimal input changes to alter model predictions, yet they often propose implausible or contextually infeasible modifications when used in isolation. Figure~\ref{fig:lcsxai_flow} illustrates the overview of the full framework. 
% \begin{figure}[pos=H]
%     \centering
%     \includegraphics[width=0.75\linewidth, height=10em]{xai_flow_2.png}
%     \caption{
%         Overview of the proposed LSC-XAI framework. The pipeline integrates three stages of interpretability: local explanations via LIME for true positive predictions, global feature attributions via SHAP for both true positives and false negatives and actionable counterfactual reasoning on the intersection of feasible features. The framework outputs frequency-ranked, modifiable feature suggestions to support transparent, instance-level and population-level insights for attrition mitigation.
%     }
%     \label{fig:lcsxai_flow}
% \end{figure}

To facilitate understanding of the proposed hybrid interpretability method, we first define the key notation used throughout this section.
% \begin{itemize}
%     \item \( \text{LIME}(f, \mathbf{x}^*, i) \): Local importance score for feature \( i \)
%     \item \( \text{SHAP}(f, \mathbf{x}^*, i, y) \): SHAP value for feature \( i \) and class \( y \)
%      \item \( \text{FEASIBLE}(f, \mathbf{x}^*, i) \): Feasible feature selection along feature \( i \)
%     \item \( \text{CF}(f, \mathbf{x}^*, i) \): Counterfactual generation along feature \( i \)
%     \item \( \delta \in \mathbb{R}^d \): Feature modification direction vector
%     \item \( \phi(\delta) \in [0, 1] \): Normalized frequency of modification \( \delta \)
% \end{itemize}

\begin{itemize}
    \item $\text{LIME}(f, \mathbf{x}^*, i)$: Local importance score for feature $i$
    \item $\text{SHAP}(f, \mathbf{x}^*, i, y)$: SHAP value for feature $i$ and class $y$
    \textcolor{black}{\item $\mathcal{F}(f_i) \in \{0,1\}$: Feasibility indicator for feature $f_i$ under mutability $\mu(f_i)$, organizational controllability $\omega(f_i)$ and ethical permissibility $\eta(f_i)$}
    \textcolor{black}{\item $\mathcal{F}^* \subseteq \mathcal{C}$: Feasible feature subset derived from the LIME--SHAP intersection $\mathcal{C}$}
    \item $\text{CF}(f, \mathbf{x}^*, i)$: Counterfactual generation along feature $f_i \in \mathcal{F}^*$
    \item $\boldsymbol{\delta} \in \mathbb{R}^d$: Feature modification direction vector
    \item $\phi(\boldsymbol{\delta}) \in [0, 1]$: Normalised frequency of modification $\boldsymbol{\delta}$
\end{itemize}

\subsubsection*{\textcolor{black}{Formal Feasibility Constraint Framework}}
\label{subsec:formal_feasibility}
\textcolor{black}{Feature feasibility was determined through a structured constraint function grounded in the actionable recourse literature~\cite{ustun2019actionable, karimi2021algorithmic}. Formally, a feature $f_i$ is designated feasible if and only if it satisfies all three of the following binary criteria:}

\textcolor{black}{\begin{equation}
    \mathcal{F}(f_i) = \mathbf{1}\left[\mu(f_i) = 1\right] \cdot 
    \mathbf{1}\left[\omega(f_i) = 1\right] \cdot 
    \mathbf{1}\left[\eta(f_i) = 1\right]
\end{equation}}


\textcolor{black}{\noindent where:}
\begin{itemize}
    \item \textcolor{black}{$\mu(f_i) \in \{0,1\}$ denotes \textbf{mutability} — whether $f_i$ is an actionable, non-immutable attribute (i.e., not a fixed characteristic such as age or historical employment records)~\cite{wachter2017counterfactual},}

    \textcolor{black}{\item $\omega(f_i) \in \{0,1\}$ denotes \textbf{organizational controllability} — whether the employing institution has direct authority to modify $f_i$, consistent with the job characteristics model distinguishing employee-controlled from organization-controlled attributes~\cite{hackman1976motivation},}

    \textcolor{black}{\item $\eta(f_i) \in \{0,1\}$ denotes \textbf{ethical permissibility} — whether modification of $f_i$ is free from fairness concerns or legal constraints under employment law~\cite{barocas2023fairness}.}
\end{itemize}

\textcolor{black}{\noindent The feasible feature set $\mathcal{F}^*$ used for counterfactual generation is then defined as:}
\textcolor{black}{\begin{equation}
    \mathcal{F}^* = \left\{ f_i \in \mathcal{C} \;\middle|\; 
    \mathcal{F}(f_i) = 1 \right\}
\end{equation}}

\noindent where $\mathcal{C} = \mathcal{F}_{\text{LIME}} \cap 
\mathcal{F}_{\text{SHAP}}$ denotes the intersection of top-$n$ features identified by LIME and SHAP respectively.\\

In the proposed algorithm \ref{alg:lscxai}, LIME is employed initially to obtain the local importance scores \(\text{LIME}(f, \mathbf{x}^*, i)\) of features contributing to explainability specific to detecting true positive outcomes—in this study, whether an employee predicted to leave the job actually leaves. Subsequently, SHAP values \(\text{SHAP}(f, \mathbf{x}^*, i, y)\) are computed to identify important features that contribute to both true positive and false negative predictions, where false negative refers to an employee predicted to stay but who actually leaves. \textcolor{black}{The intersection $\mathcal{C}$ of features identified by both methods yields a subset of consistently important attributes. The feasibility constraint $\mathcal{F}^*$ is then applied to $\mathcal{C}$, retaining only features satisfying mutability, organizational controllability and ethical permissibility criteria (Equation~(2)), thereby grounding counterfactual generation in organizationally actionable modifications. Counterfactuals $\text{CF}(f, \mathbf{x}^*, i)$ are generated exclusively over $f_i \in \mathcal{F}^*$, exploring plausible feature modification directions $\boldsymbol{\delta} \in \mathbb{R}^d$. Finally, the normalized modification frequency $\phi(\boldsymbol{\delta}) \in [0,1]$ is computed across all instances and the highest-ranked modifications are surfaced as the actionable recommendations of the LSC-XAI framework. The LSC-XAI algorithm is formally structured in the following.}

\textcolor{blue}{Counterfactual generation is implemented with the dice-ml library (DiCE v0.12) using its random-search method (\texttt{dice\_ml.Dice(..., method="random")}, scikit-learn backend) over the trained base model, with the SMOTE-resampled, fold-wise preprocessed training fold passed as the data interface. For each instance, the generator is requested to produce three counterfactuals (\texttt{total\_CFs} = 3) of the opposite class (\texttt{desired\_class} = "opposite"), with \texttt{features\_to\_vary} restricted to the feasible feature set $\mathcal{F}^*$; a candidate is valid when the base model's predicted class flips, and generations that fail to produce a valid counterfactual are skipped rather than forced. In the headline pipeline, the recommendation loop attempts every correctly classified attrition employee in the held-out test fold of the representative model, and Section~5.5 reports the sensitivity of the recommendations when enlarged feasible feature sets are permitted.}

\textcolor{blue}{We clarify the relationship of this formalism to existing constrained-recourse methods, since the notion of restricting which features may vary is not new. Counterfactual libraries such as DiCE~\cite{mothilal2020dice} already support user-specified restrictions on feature variation, and the actionable-recourse literature already formalizes mutability and actionability constraints~\cite{ustun2019actionable,karimi2021algorithmic}. The novelty claimed here is therefore deliberately narrower than constrained recourse itself: the mutability indicator $\mu$ reuses that standard notion, while the organizational-controllability indicator $\omega$ and the ethical-permissibility indicator $\eta$ operationalize it for the healthcare-HR setting, where the decision-relevant question is not only whether an attribute \emph{can} in principle change but whether the \emph{employing institution} may legitimately change it, in which direction, and with what organizational meaning. The per-feature assignments, rationales and permitted directions reported in \cref{tab:feasibility_complete} are the auditable artifact of that operationalization; the assignments themselves remain author judgement informed by the domain literature (see the Limitations, Section~6.5) rather than learned or empirically validated quantities. Read as a whole, the framework does not replace an existing constrained generator: it composes an existing one (DiCE) with a cross-method agreement gate ($\mathcal{C}$) and a frequency-aggregation step, so that recourse is generated only over features that are jointly endorsed by two independent explainers and institutionally actionable, and is reported as ranked population-level recommendations rather than isolated per-instance counterfactuals. To our knowledge, this particular composition, and its empirical evaluation in the healthcare-attrition domain, has not previously been reported.}
\subsection{Comparison of Other Frameworks with LSC-XAI}
In the following table \ref{tab:lsc_comparison}, comparison among LIME, SHAP, Counterfactuals and LSC-XAI is provided. \textcolor{blue}{The first three rows describe the design scope of each technique and the limitations documented in the literature, rather than measured performance differences; the final row reports the quantitative status of each method as measured in Section~5.6, and the claimed benefits of LSC-XAI are restricted accordingly.}
\begin{table}[width=1\linewidth,cols=5,pos=H]
\centering
\caption{Based on various attributes, comparison between existing interpretability techniques and LSC-XAI is shown. \textcolor{blue}{The first three rows describe design scope and documented limitations; the final row reports the quantitative status of each method measured in Section~5.6 (\cref{tab:xai_metrics}, \ref{tab:cf_metrics}).}}
\label{tab:lsc_comparison}
\renewcommand{\arraystretch}{1.075} % Adjust row height
\begin{tabular*}{\tblwidth}{p{0.12\linewidth}p{0.16\linewidth}p{0.16\linewidth}p{0.16\linewidth}p{0.27\linewidth}}
\hline
\textbf{Feature} & \textbf{LIME} & \textbf{SHAP} & \textbf{Counterfactuals} & \textbf{LSC-XAI} \\
\hline
Robustness and Reliability &  
Suffers from instability across different samples and perturbation strategies & 
Can be computationally expensive and may miss local nuances & 
\textcolor{blue}{May generate suggestions that are unrealistic or organizationally infeasible} & \textcolor{blue}{Requires a feature to be ranked important by both local (LIME) and global (SHAP) explanations before it is used for counterfactual generation (a design-level cross-method gate rather than a measured robustness gain)}
\\\hline
Comprehensive Perspective &  
Only local explanations & 
Primarily global view & 
Instance-specific suggestions without broader context &  \textcolor{blue}{Couples local (LIME), global (SHAP) and actionable (counterfactual) views within a single pipeline and returns a ranked set of feasible modifications rather than a per-instance explanation (design scope)}
\\\hline
Actionable Recommendations &  
Provide importance scores (locally) but limited actionable guidance & 
Provide importance scores (globally) but limited actionable guidance & 
May suggest changes to less relevant features &  \textcolor{blue}{Restricts modifications to the feasibility-filtered LIME--SHAP intersection and ranks them by modification frequency, yielding population-level recommendations (design scope)}
\\\hline
\textcolor{blue}{Empirical status (Section~5.6)} & \textcolor{blue}{Stability 0.8034; fidelity 0.9375; sparsity fixed at 0.4783} & \textcolor{blue}{Stability 1.0000 (ceiling, expected for a linear base model); fidelity 0.9375; sparsity fixed at 0.4783} & \textcolor{blue}{Proximity 6.1206; plausibility 0.2308} & \textcolor{blue}{Stability 1.0000 (inherited through the intersection, not an independent gain); fidelity 0.9405, not measurably different from either standalone method; sparsity 0.8696 and proximity 1.2449 follow by construction; plausibility 0.2143, not statistically distinguishable from the standalone generator} \\\hline
\end{tabular*}
\end{table}
\textcolor{blue}{Each of the other methods may provide different explanations for the same instance, which can complicate interpretation and reduce trust. LSC-XAI addresses this by construction: it returns only features endorsed by both explainers and restricts modifications to organizational-actionable attributes. We deliberately do not claim that this reduces explanation variance, increases explanation reliability or reduces false positives relative to the standalone methods. Section~5.6 shows that LSC-XAI matches SHAP's stability ceiling and incurs no measurable fidelity cost, whereas its sparsity and counterfactual proximity follow deterministically from the restricted feature set and its counterfactual plausibility is statistically indistinguishable from that of the standalone generator. Establishing reliability or variance-reduction claims would require a dedicated repeated-run study beyond the scope of the present evaluation.}

\begin{algorithm}[H]
\caption{LSC-XAI: LIME--SHAP--Counterfactual Explainable AI}
\label{alg:lscxai}
\vspace{1mm}
\textbf{Input:} Trained model $f : \mathbb{R}^d \rightarrow \{0,1\}$, instance $\mathbf{x}^* \in \mathbb{R}^d$, top-$n$ threshold\\
\textbf{Output:} Frequency-ranked actionable feature modifications $\boldsymbol{\delta}$
\vspace{1mm}
\hrule
\vspace{1mm}

\textbf{Step 1: Feature Importance Extraction}\\
Compute LIME scores for the positive class:
\begin{equation*}
    \mathbf{l} = [l_1, \dots, l_d], \quad l_i = \text{LIME}(f, \mathbf{x}^*, i)
\end{equation*}
Compute SHAP values for both classes:
\begin{equation*}
    \mathbf{s}^{(0)} = [s_1^{(0)}, \dots, s_d^{(0)}], \quad
    \mathbf{s}^{(1)} = [s_1^{(1)}, \dots, s_d^{(1)}], \quad
    s_i^{(y)} = \text{SHAP}(f, \mathbf{x}^*, i, y)
\end{equation*}
Extract top-$n$ feature sets:
\begin{equation*}
    \mathcal{F}_{\text{LIME}} = \{i \mid l_i \in \text{top-}n\}, \quad
    \mathcal{F}_{\text{SHAP}} = \{i \mid s_i^{(0)} \in \text{top-}n \ \vee \ s_i^{(1)} \in \text{top-}n\}
\end{equation*}

\vspace{1mm}
\textbf{\textcolor{black}{Step 2: Feasibility-Constrained Feature Identification}}\\
\textcolor{black}{Apply the three-criterion feasibility constraint (Equation~(10)) to the LIME--SHAP intersection:}
\textcolor{black}{\begin{equation*}
    \mathcal{F}^* = \left(\mathcal{F}_{\text{LIME}} \cap \mathcal{F}_{\text{SHAP}}\right) \cap 
    \left\{f_i \mid \mathcal{F}(f_i) = 1\right\}
\end{equation*}}
\textcolor{black}{where $\mathcal{F}(f_i) = \mathbf{1}[\mu(f_i){=}1] \cdot \mathbf{1}[\omega(f_i){=}1] \cdot \mathbf{1}[\eta(f_i){=}1]$}

\vspace{5mm}
\textbf{Step 3: Counterfactual Generation}\\
For each $f_i \in \mathcal{F}^*$, generate counterfactual instances:
\begin{equation*}
    \mathcal{C}_i = \left\{\tilde{\mathbf{x}}_j^{(i)}\right\}_{j=1}^{k}, \quad
    \tilde{\mathbf{x}}_j^{(i)} = \text{CF}(f, \mathbf{x}^*, f_i)
\end{equation*}
Aggregate all counterfactuals:
\begin{equation*}
    \mathcal{C} = \bigcup_{i} \mathcal{C}_i
\end{equation*}

\vspace{1mm}
\textbf{Step 4: Frequency-Based Actionable Insights}\\
Compute modification vectors and their normalised frequencies:
\begin{equation*}
    \boldsymbol{\delta}_j = \tilde{\mathbf{x}}_j - \mathbf{x}^*, \quad \forall\, \tilde{\mathbf{x}}_j \in \mathcal{C}
\end{equation*}
\begin{equation*}
    \phi(\boldsymbol{\delta}) = 
    \frac{\left|\left\{\tilde{\mathbf{x}} \in \mathcal{C} \mid \tilde{\mathbf{x}} - \mathbf{x}^* = \boldsymbol{\delta}\right\}\right|}{|\mathcal{C}|}, 
    \quad \phi(\boldsymbol{\delta}) \in [0,1]
\end{equation*}
\vspace{1mm}
\textbf{Return:} Set of $\boldsymbol{\delta}$ ranked by $\phi(\boldsymbol{\delta})$ in descending order
\end{algorithm}

\section{Experiments and Results}
\label{sec:exp}
This section provides a systematic evaluation of the proposed methodology, beginning with a comparative performance analysis of the machine learning models across the established evaluation metrics. To ensure the reliability of these findings, a Wilcoxon signed-rank test is performed to statistically validate the superiority of the models. A model is then selected, followed by an ablation study to quantify the impact of specific preprocessing components. Subsequently, we deploy the LSC-XAI framework to the selected model. The robustness of the framework is then assessed through a quantitative analysis of its explanatory fidelity and finally, the framework is validated using an empirical application to a real-world healthcare dataset.

\subsection{\textcolor{black}{Computational Environment}}
\textcolor{black}{To ensure reproducibility, we report the computational environment, software implementations, model configurations, random seeds and tuning procedures used in the experiments. Classical learners and stacking were implemented using scikit-learn v1.6.1. XGBoost, LightGBM, CatBoost, TabNet and the neural-network models were implemented using their respective Python libraries.
Most classical and ensemble models were trained using library-default configurations to establish a reproducible baseline. The KNN model was not hyperparameter-tuned in the submitted implementation and used the default KNeighborsClassifier() configuration. No automated grid search, randomized search, or Bayesian optimization procedure was used in the current implementation. Therefore, the final settings, software implementations, random seeds and tuning procedures are reported in \cref{tab:software_implementations}. \textcolor{blue}{For the neural models, the PyTorch seed is fixed per fold before model construction and training (\texttt{torch.manual\_seed(42 + fold)}; \cref{tab:software_implementations}), with \texttt{num\_workers}=0 and CPU execution, and the classical and ensemble learners receive the corresponding per-fold \texttt{random\_state} values; the reported neural-model results are therefore exactly reproducible within the archived environment.}
The experiments were conducted using the Google Colab cloud-based environment with a Python 3 runtime. The computational workload was handled by a Google Compute Engine backend with 12.7 GB RAM. No GPU was used for the shallow or ensemble learners, as these models were trained using CPU-based execution.
To reduce data leakage, the train-test split was performed before resampling and model training. SMOTE-NC upsampling was applied only to the training data. The reported results should be interpreted as performance from fixed model configurations rather than from an exhaustive hyperparameter-optimized pipeline. \textcolor{blue}{The protocol by which the base model, the preprocessing configuration, the representative fold and the top-$N$ explanation threshold were selected is described in Section~5.4, and its consequences for interpretation are discussed in Section~6.5.}}

\subsection{\textcolor{black}{Model Comparison Results}}
\label{model-comparison}
\textcolor{blue}{In this section, precision, recall and F1-score denote macro-averages over the two classes, whereas accuracy denotes the overall fraction of correct predictions. Macro-recall is $(\mathrm{TPR}+\mathrm{TNR})/2$, equivalent to balanced accuracy for this binary task; it is not attrition-class sensitivity. The reported Logistic Regression recall of 83.11\% is therefore a macro-average, not the percentage of attriting employees identified.} The mean of the performance metrics (Accuracy, Precision, Recall, F1-score) and \textcolor{black}{standard deviation} are presented in \cref{tab:model_comparison}. As evident from the results, CatBoost emerged as the top-performing model across three out of four performance metrics. It achieved the highest mean accuracy of \textcolor{black}{91.41 ± 1.63\%}, a precision of \textcolor{black}{80.11 ± 4.54\%} and F-1 score of \textcolor{black}{78.57 ± 3.72\%}. In terms of recall, the Logistic Regression model performed the best with \textcolor{black}{83.11 ± 1.34\%} score.

\begin{table}[width=1\linewidth,cols=5,pos=H]
\centering
\caption{Model performance metrics, including mean and standard deviation across 5-fold stratified
cross-validation, computed under the fold-wise (leakage-safe) preprocessing pipeline described in
Section 3.2. \textcolor{blue}{Precision, recall and F1-score are macro-averaged over the attrition and non-attrition classes; accuracy is overall classification accuracy.} The highest value in each column is highlighted in bold.}
\label{tab:model_comparison}
\renewcommand{\arraystretch}{1.075} % Adjust row height
\begin{tabular}{lcccc}
\toprule
\textbf{Model} & \textbf{Accuracy} & \textbf{Precision} & \textbf{Recall} & \textbf{F1-score} \\
\midrule
Logistic Regression        & 87.47 $\pm$ 1.25\% & 72.54 $\pm$ 1.62\% & \textbf{83.11 $\pm$ 1.34\%} & 76.05 $\pm$ 1.72\% \\
K-Nearest Neighbors         & 74.70 $\pm$ 2.27\% & 61.89 $\pm$ 1.09\% & 74.12 $\pm$ 2.13\%          & 62.36 $\pm$ 1.71\% \\
Decision Trees              & 85.38 $\pm$ 1.61\% & 67.72 $\pm$ 2.21\% & 72.54 $\pm$ 3.58\%          & 69.46 $\pm$ 2.45\% \\
Support Vector Machines     & 89.32 $\pm$ 1.32\% & 74.84 $\pm$ 2.71\% & 81.73 $\pm$ 4.98\%          & 77.54 $\pm$ 3.46\% \\
Random Forest               & 90.99 $\pm$ 1.24\% & 78.60 $\pm$ 2.98\% & 77.25 $\pm$ 4.52\%          & 77.82 $\pm$ 3.75\% \\
Gradient Boosting           & 90.10 $\pm$ 1.46\% & 76.46 $\pm$ 3.25\% & 76.98 $\pm$ 2.44\%          & 76.68 $\pm$ 2.78\% \\
Adaptive Boosting           & 88.78 $\pm$ 1.92\% & 73.85 $\pm$ 3.39\% & 79.28 $\pm$ 3.72\%          & 76.09 $\pm$ 3.57\% \\
Extreme Gradient Boosting   & 90.75 $\pm$ 1.80\% & 78.81 $\pm$ 5.11\% & 76.25 $\pm$ 3.06\%          & 77.21 $\pm$ 3.30\% \\
Light Gradient Boosting     & 91.23 $\pm$ 0.75\% & 79.49 $\pm$ 1.77\% & 76.96 $\pm$ 3.44\%          & 78.05 $\pm$ 2.37\% \\
Categorical Boosting        & \textbf{91.41 $\pm$ 1.63\%} & \textbf{80.11 $\pm$ 4.54\%} & 77.29 $\pm$ 3.18\% & \textbf{78.57 $\pm$ 3.72\%} \\
Stacking                    & 90.99 $\pm$ 1.49\% & 78.58 $\pm$ 3.66\% & 78.35 $\pm$ 4.03\%          & 78.38 $\pm$ 3.62\% \\
Multi-layer Perceptron      & 86.22 $\pm$ 1.33\% & 70.95 $\pm$ 1.95\% & 81.95 $\pm$ 4.01\%          & 74.27 $\pm$ 2.16\% \\
Convolutional Neural Networks & 86.93 $\pm$ 1.25\% & 71.59 $\pm$ 1.93\% & 81.46 $\pm$ 4.84\%        & 74.77 $\pm$ 2.47\% \\
TabNet                      & 78.88 $\pm$ 4.74\% & 62.63 $\pm$ 2.86\% & 70.81 $\pm$ 5.08\%          & 63.71 $\pm$ 3.16\% \\
\bottomrule
\end{tabular}
\end{table}

% \begin{table}[width=1\linewidth,cols=5,pos=H]
% \centering
% \caption{\textcolor{black}{Model performance metrics, including mean accuracy and \textcolor{black}{absolute mean deviation} across \textcolor{black}{5 folds of stratified cross-validation}. The highest values from each matrices are highlighted in bold.}}
% \label{tab:model_comparison}
% \renewcommand{\arraystretch}{1.075} % Adjust row height
% \begin{tabular*}{\tblwidth}{L|C|C|C|C}
% \hline
% \textbf{\textcolor{black}{Model}} & \textbf{\textcolor{black}{Accuracy}} & \textbf{\textcolor{black}{Precision}} & \textbf{\textcolor{black}{Recall}} & \textbf{\textcolor{black}{F1-score}} \\
% \hline
% \textcolor{black}{Logistic Regression} & \textcolor{black}{88.72 $\pm$ 1.78\%} & \textcolor{black}{74.30 $\pm$ 2.87\%} & \textbf{\textcolor{black}{83.61 $\pm$ 3.08\%}} & \textcolor{black}{77.65 $\pm$ 2.99\%} \\\hline
% \textcolor{black}{K-Nearest Neighbors} & \textcolor{black}{73.03 $\pm$ 1.99\%} & \textcolor{black}{61.49 $\pm$ 1.54\%} & \textcolor{black}{74.25 $\pm$ 3.57\%} & \textcolor{black}{61.30 $\pm$ 1.86\%} \\\hline
% \textcolor{black}{Decision Trees} & \textcolor{black}{86.99 $\pm$ 0.85\%} & \textcolor{black}{69.75 $\pm$ 2.10\%} & \textcolor{black}{72.82 $\pm$ 3.70\%} & \textcolor{black}{71.07 $\pm$ 2.70\%} \\\hline
% \textcolor{black}{Support Vector Machines} & \textcolor{black}{89.98 $\pm$ 2.21\%} & \textcolor{black}{76.04 $\pm$ 4.29\%} & \textcolor{black}{80.63 $\pm$ 5.67\%} & \textcolor{black}{78.00 $\pm$ 4.79\%} \\\hline
% \textcolor{black}{Random Forest} & \textcolor{black}{91.05 $\pm$ 0.81\%} & \textcolor{black}{79.42 $\pm$ 2.15\%} & \textcolor{black}{74.91 $\pm$ 2.64\%} & \textcolor{black}{76.85 $\pm$ 2.24\%} \\\hline
% \textcolor{black}{Gradient Boosting} & \textcolor{black}{90.99 $\pm$ 0.61\%} & \textcolor{black}{78.49 $\pm$ 1.48\%} & \textcolor{black}{78.13 $\pm$ 2.60\%} & \textcolor{black}{78.28 $\pm$ 2.00\%} \\\hline
% \textcolor{black}{Adaptive Boosting} & \textcolor{black}{90.27 $\pm$ 1.05\%} & \textcolor{black}{76.73 $\pm$ 2.04\%} & \textcolor{black}{82.31 $\pm$ 3.76\%} & \textcolor{black}{78.99 $\pm$ 2.33\%} \\\hline
% \textcolor{black}{Extreme Gradient Boosting} & \textcolor{black}{91.83 $\pm$ 0.80\%} & \textcolor{black}{81.48 $\pm$ 2.88\%} & \textcolor{black}{77.73 $\pm$ 2.30\%} & \textcolor{black}{79.31 $\pm$ 1.80\%} \\\hline
% \textcolor{black}{Light Gradient Boosting} & \textcolor{black}{91.77 $\pm$ 1.04\%} & \textcolor{black}{81.27 $\pm$ 2.76\%} & \textcolor{black}{77.04 $\pm$ 3.45\%} & \textcolor{black}{78.87 $\pm$ 3.02\%} \\\hline
% \textcolor{black}{Categorical Boosting} & \textbf{\textcolor{black}{92.24 $\pm$ 1.09\%}} & \textbf{\textcolor{black}{82.67 $\pm$ 3.15\%}} & \textcolor{black}{78.20 $\pm$ 2.44\%} & \textbf{\textcolor{black}{80.16 $\pm$ 2.56\%}} \\\hline
% \textcolor{black}{Stacking} & \textcolor{black}{91.83 $\pm$ 1.40\%} & \textcolor{black}{80.93 $\pm$ 3.54\%} & \textcolor{black}{79.04 $\pm$ 3.39\%} & \textcolor{black}{79.89 $\pm$ 3.26\%} \\\hline
% \textcolor{black}{Multi-layer Perceptron} & \textcolor{black}{88.49 $\pm$ 2.00\%} & \textcolor{black}{74.04 $\pm$ 3.56\%} & \textcolor{black}{82.82 $\pm$ 2.04\%} & \textcolor{black}{77.16 $\pm$ 3.03\%} \\\hline
% \textcolor{black}{Convolutional Neural Networks} & \textcolor{black}{88.84 $\pm$ 2.27\%} & \textcolor{black}{74.24 $\pm$ 4.05\%} & \textcolor{black}{81.06 $\pm$ 3.59\%} & \textcolor{black}{76.92 $\pm$ 4.01\%} \\\hline
% \textcolor{black}{TabNet} & \textcolor{black}{81.68 $\pm$ 6.27\%} & \textcolor{black}{65.29 $\pm$ 4.90\%} & \textcolor{black}{71.57 $\pm$ 2.96\%} & \textcolor{black}{66.51 $\pm$ 5.55\%} \\\hline
% \end{tabular*}
% \end{table}

% \textcolor{black}{While \cref{tab:model_comparison} presents mean performance across 5-fold stratified cross-validation runs, point estimates alone are insufficient to establish statistical superiority between the models. To address this, pairwise Wilcoxon signed-rank tests were conducted on the recall scores obtained from 5-fold stratified cross-validation across all 14 models at a significance threshold of $\alpha$ = 0.05.}

% \textcolor{black}{Recall was adopted as the primary evaluation criterion for this comparison because, in a workforce retention context, false negatives — meaning employees predicted to stay who subsequently leave — represent the operationally critical error. An attrition prediction framework that misses at-risk employees undermines the very purpose of any downstream retention intervention, including the LSC-XAI recommendations generated in section \ref{subsec:lsc_res}.}

\textcolor{black}{\textcolor{blue}{The recall comparisons discussed here concern macro-averaged recall, not attrition-class sensitivity.} While table \ref{tab:model_comparison} presents mean performance across 5-fold stratified cross-validation, point estimates alone are insufficient to establish statistical differences between models and unadjusted pairwise testing is underpowered at this sample size. We therefore first conducted an omnibus Friedman test across all fourteen models' recall scores, followed by Holm-corrected post-hoc pairwise Wilcoxon signed-rank tests. The Friedman test detected a statistically significant difference in recall among the fourteen models ($\chi^2(13) = 43.79$, $p < 0.001$), with a large effect size (Kendall's W = 0.674), indicating that model choice meaningfully affects recall. However, none of the ninety-one pairwise post-hoc comparisons remained significant after Holm correction. Notably, several pairs including Ada Boost vs. Logistic Regression, CNN vs. Random Forest and Logistic Regression vs. XGBoost exhibited a rank-biserial effect size of $|r| = 1.0$, indicating that one model outperformed the other in every one of the five folds with no exceptions; even this maximal, perfectly consistent effect cannot reach significance after Holm correction at n=5, since the Wilcoxon signed-rank test cannot produce a two-sided raw p-value below 0.0625 regardless of the magnitude of the underlying difference. We therefore revise our conclusion accordingly: the omnibus test provides clear evidence that real performance differences exist among the fourteen models, but the current 5-fold evaluation design lacks the statistical power to resolve which specific pairs differ, even in cases of perfectly consistent directional effects. This finding should not be interpreted as evidence that model choice has no effect on recall; rather, resolving pairwise differences reliably would require a larger number of cross-validation folds or repetitions. Table \ref{tab:recall_ci} reports mean recall with 95\% confidence intervals for all fourteen models and table \ref{tab:posthoc_top} reports the top pairwise comparisons ranked by effect size (full results in Appendix \cref{app:pairwise_statistical}). Figure \ref{fig:wilcoxon} illustrates the p-value among all the models and figure \ref{fig:recall_vs_complexity} demonstrates the relationship between mean recall values and model complexity.}

\begin{table}[pos=h]
\centering
\caption{\textcolor{black}{\textcolor{blue}{Mean macro-averaged recall (the unweighted mean of attrition and non-attrition recall)} with 95\% confidence intervals ($t$-distribution, $n=5$ folds) across all fourteen models, computed under the fold-wise leakage-safe pipeline. Models are sorted by mean recall.}}
\label{tab:recall_ci}
\begin{tabular}{lccc}
\toprule
\textbf{Model} & \textbf{Mean Recall} & \textbf{Std Dev} & \textbf{95\% CI} \\
\midrule
Logistic Regression & 0.8311 & 0.0134 & [0.8145, 0.8477] \\
Multi-layer Perceptron & 0.8195 & 0.0401 & [0.7697, 0.8693] \\
Support Vector Machines & 0.8173 & 0.0498 & [0.7554, 0.8791] \\
Convolutional Neural Networks & 0.8146 & 0.0484 & [0.7545, 0.8746] \\
Adaptive Boosting & 0.7928 & 0.0372 & [0.7466, 0.8390] \\
Stacking & 0.7835 & 0.0403 & [0.7334, 0.8335] \\
Categorical Boosting & 0.7729 & 0.0318 & [0.7333, 0.8124] \\
Random Forest & 0.7725 & 0.0452 & [0.7163, 0.8286] \\
Gradient Boosting & 0.7698 & 0.0244 & [0.7395, 0.8001] \\
Light Gradient Boosting & 0.7696 & 0.0344 & [0.7269, 0.8124] \\
Extreme Gradient Boosting & 0.7625 & 0.0306 & [0.7245, 0.8005] \\
K-Nearest Neighbors & 0.7412 & 0.0213 & [0.7148, 0.7676] \\
Decision Trees & 0.7254 & 0.0358 & [0.6810, 0.7699] \\
TabNet & 0.7081 & 0.0508 & [0.6451, 0.7711] \\
\bottomrule
\end{tabular}
\end{table}

\begin{table}[pos=h]
\centering
\caption{\textcolor{black}{Top pairwise Wilcoxon signed-rank post-hoc comparisons across all fourteen models' \textcolor{blue}{macro-averaged recall} scores, ranked by absolute rank-biserial effect size, Holm-corrected
for multiple comparisons across all 91 pairs. All raw $p$-values are bounded below by
0.0625, the minimum achievable two-sided $p$-value for a paired test with $n=5$;
consequently no pair survives Holm correction, including pairs with maximal effect
size ($|r|=1.0$). Full 91-pair table in \cref{tab:app_statistical_comparisons}.}}
\label{tab:posthoc_top}
\begin{tabular}{llcccc}
\toprule
\textbf{Model 1} & \textbf{Model 2} & \textbf{Rank-biserial $r$} & \textbf{$p$ (raw)} & \textbf{$p$ (Holm-adj.)} & \textbf{Significant} \\
\midrule
Ada Boost & Logistic Regression & $-1.00$ & 0.0625 & 1.00 & No \\
Ada Boost & XGBoost & $1.00$ & 0.0625 & 1.00 & No \\
Ada Boost & TabNet & $1.00$ & 0.0625 & 1.00 & No \\
CNN & Random Forest & $1.00$ & 0.0625 & 1.00 & No \\
CNN & Stacking & $1.00$ & 0.0625 & 1.00 & No \\
CNN & XGBoost & $1.00$ & 0.0625 & 1.00 & No \\
CNN & TabNet & $1.00$ & 0.0625 & 1.00 & No \\
CatBoost & Decision Tree & $1.00$ & 0.0625 & 1.00 & No \\
% CNN & Decision Tree & $1.00$ & 0.0625 & 1.00 & No \\
% MLP & Random Forest & $1.00$ & 0.0625 & 1.00 & No \\
% Logistic Regression & XGBoost & $1.00$ & 0.0625 & 1.00 & No \\
% LightGBM & SVM & $-1.00$ & 0.0625 & 1.00 & No \\
% LightGBM & Logistic Regression & $-1.00$ & 0.0625 & 1.00 & No \\
% Logistic Regression & TabNet & $1.00$ & 0.0625 & 1.00 & No \\
% Gradient Boosting & Logistic Regression & $-1.00$ & 0.0625 & 1.00 & No \\
\bottomrule
\end{tabular}
\end{table}

\begin{figure}[pos=H]
    \centering
    \includegraphics[width=0.8\linewidth, height=0.45\linewidth]{wilcoxon_signed_rank_test.jpg}
    \caption{Statistical comparison using Wilcoxon Signed-Rank Test}
    \label{fig:wilcoxon}
\end{figure}

% \textcolor{black}{The Wilcoxon signed-rank tests' score shown in \cref{fig:wilcoxon} revealed that no model achieved a statistically significant improvement in recall over any other model (all pairwise p-values > 0.05). This finding is practically significant: it demonstrates that in the healthcare employee attrition prediction setting, increasing model complexity — from linear models through ensembles to deep learning — does not yield a statistically measurable improvement in the metric that matters most operationally.
% Given the absence of statistically significant recall differences, final model selection was determined by two secondary criteria applied in order of priority: (1) recall value and (2) model complexity, favoring simpler models under equivalent performance. 

\begin{figure}[pos=H]
    \centering
    \includegraphics[width=1\linewidth]{mean_recall_vs_complexity_graph.PNG}
    \caption{\textcolor{black}{Mean Recall vs Model Complexity plot. Logistic Regression is algorithmically lesser complex model among the other models, with the highest recall value from the evaluation results.}}
    \label{fig:recall_vs_complexity}
\end{figure}

% \pagebreak[2] % Suggests a page break here if it's REALLY tight, but not forced
% \noindent
% \begin{minipage}{\linewidth}
%     \centering
%     \includegraphics[width=1\linewidth, height=0.3\linewidth]{mean_recall_vs_complexity_graph.PNG}
%     \captionof{figure}{Mean recall vs model complexity plot. It denotes Logistic Regression to be in superior position under these alignments}
%     \label{fig:recall_vs_complexity}
% \end{minipage}

\textcolor{black}{This approach is consistent with the framework's emphasis on deployability and transparency in healthcare HR settings. Under these criteria, Logistic Regression emerged as the preferred model. Additionally, the confusion matrices and Receiver Operating Characteristics (ROC) curve outputs of 1 set of the K-Fold CV training are showcased in Appendix \cref{appendix:confusion}, \cref{appendix:roc}. Figure \ref{fig:roc_subfigures} shows the ROC-AUC curves of the trained classifiers. The CatBoost and Ada Boost models showed the best performance with 0.92 in the ROC-AUC score.}

% \newline
% Among the shallow models, SVM showed better performance with 89\% in accuracy, 74\% in precision, 82\% in recall and 77\% in F1-score, beating all the other shallow models in each of the evaluation criteria. Among the ensemble models, CatBoost showed best results in all the criteria, with 91\% in accuracy, 80\% in precision, recall and F1-score. And among the deep models, CNN performed better, achieving 89\% in accuracy, 75\% in precision, 85\% in recall and 79\% in F1-score. Considering all the evaluation scores, CatBoost outperformed all the other models, showcasing stable training and performance. Throughout the training, not much variance in training of the models is observed.


\subsection{Ablation study}
% To systematically evaluate the contributions of various preprocessing steps, an ablation study using logistic regression as the baseline machine learning model was conducted. Logistic regression was chosen as the baseline model due to its well-documented interpretability and established role as a reference in classification tasks \cite{hastie2009elements, james2013introduction}. The study proceeded in a stepwise manner, isolating the effects of feature selection, outlier removal and data upsampling techniques on key performance metric: F1-score. Each step was designed to assess the impact of individual components on the overall predictive capability of the model. The findings are shown in \cref{tab:abl_stud}.

% To systematically evaluate the contributions of various preprocessing steps, an ablation study was conducted using \textcolor{black}{the logistic regression algorithm} as the baseline machine learning model. The study proceeded in a stepwise manner, isolating the effects of feature selection, distributional transformation and data upsampling techniques on key performance metric: Recall. Each step was designed to assess the impact of individual components on the overall predictive capability of the model. The findings are shown in \cref{tab:ablation_stud}

% Initially, in state (i) the reference model was trained on the original encoded dataset to establish a baseline. This yielded a reference performance metric that provided a point of comparison for subsequent experiments. In state (ii), feature elimination techniques were applied, as discussed in the previous section, to identify and remove columns deemed non-essential to the predictive task. Finally, in state (iii) to address outliers and class imbalance in the dataset, the box-cox transformation and SMOTE-NC were applied on top of the previous state. The logistic regression model was trained on this new state of the dataset and the resulting performance metrics were compared with the baseline and feature-reduced configurations. 

% \begin{table}[width=1\linewidth,cols=7,pos=H]
% \centering
% \caption{\textcolor{black}{Performance metrics for different preprocessing steps are shown. The last state showcases improved overall improvement in the scores.}}
% \label{tab:ablation_stud}
% \renewcommand{\arraystretch}{1} % Adjust row height
%     \begin{tabular*}{\tblwidth}{CCCCCCC}
%     \hline
%     \textbf{State} & \textbf{Dataset dimension} & \textbf{Encoding} & \textbf{Feature elimination} & \textbf{Outlier handling} & \textbf{Upsampling} & \textbf{\textcolor{black}{Recall(\%)}}\\\hline
%     i & 1340 x 38 & \checkmark & x & x & x & \textcolor{black}{76.14}\\\hline
%     ii & 1340 x 32 & \checkmark & \checkmark & x & x & \textcolor{black}{76.14}\\\hline
%     iii & 2362 x 32 & \checkmark & \checkmark & \checkmark & \checkmark & \textcolor{black}{82.70}\\
%     \hline
%     \end{tabular*}
% \end{table}

To rigorously evaluate the individual contributions of our preprocessing pipeline, we conducted an expanded ablation study isolating the effects of Feature Selection (FS), distributional transformation (Box-Cox) and class balancing (SMOTE-NC). Table \ref{tab:ablation_study} details the performance metrics across these distinct configurations.\\
\textcolor{blue}{In this ablation, precision, recall and F1-score are macro-averaged over both classes. The baseline (Step 1) achieved macro-recall of 0.7805 and accuracy of 0.9278. Feature selection alone (Step 2) yielded macro-precision of 0.8251, macro-recall of 0.7304 and macro-F1 of 0.7653, all below their baseline point estimates.}
\textcolor{blue}{Adding Box-Cox to feature selection (Step 3) increased macro-precision to 0.8366, accuracy to 0.9195 and macro-recall to 0.7411 relative to Step 2. Adding SMOTE-NC instead (Step 4) produced higher macro-recall (0.8197), but lower macro-precision (0.7218) and accuracy (0.8735). These are descriptive comparisons of the reported fold means, not tests of significance.}
\textcolor{blue}{The full pipeline (Step 5: feature selection, Box-Cox and SMOTE-NC) achieved the highest reported mean macro-recall among the evaluated configurations (0.8311), with macro-F1 of 0.7605 and accuracy of 0.8747. Its results match the Logistic Regression row in \cref{tab:model_comparison}. Macro-recall weights both classes equally and does not by itself establish the highest attrition-class sensitivity; these results also do not establish that this combination is strictly necessary for identifying at-risk employees.}

\begin{table}[pos=h]
\centering
    \caption{Ablation study isolating individual preprocessing effects. \textcolor{blue}{Precision, recall and F1-score are macro-averaged over the attrition and non-attrition classes; accuracy is overall classification accuracy. Values are mean $\pm$ standard deviation across 5-fold cross-validation.}}
\label{tab:ablation_study}
\begin{tabular}{l l c c c c}
\hline
\textbf{Ablation Step} & \textbf{Accuracy} & \textbf{Precision} & \textbf{Recall} & \textbf{F1} \\
\hline
1. Baseline & $92.78 \pm 1.84\%$ & $84.76 \pm 5.11\%$ & $78.05 \pm 5.32\%$ & $80.83 \pm 5.26\%$ \\
2. FS & $91.59 \pm 0.74\%$ & $82.51 \pm 2.84\%$ & $73.04 \pm 3.10\%$ & $76.53 \pm 2.73\%$ \\
3. FS + Box-Cox & $91.95 \pm 0.89\%$ & $83.66 \pm 3.15\%$ & $74.11 \pm 3.55\%$ & $77.63 \pm 2.95\%$ \\
4. FS + SMOTENC & $87.35 \pm 1.57\%$ & $72.18 \pm 2.33\%$ & $81.97 \pm 2.99\%$ & $75.52 \pm 2.62\%$ \\
5. FS + Box-Cox + SMOTENC & $87.47 \pm 1.25\%$ & $72.54 \pm 1.62\%$ & $83.11 \pm 1.34\%$ & $76.05 \pm 1.72\%$ \\
\hline
\end{tabular}
\end{table}

\subsection{\textcolor{black}{LSC-XAI Interpretation Results}}
\label{subsec:lsc_res}
Based on its complexity and performance on recall metric, Logistic Regression was selected as the foundational model for interpretability analysis. \textcolor{blue}{We state this selection protocol explicitly, because it conditions how the reported results should be read. No validation split distinct from the reported cross-validation folds exists in this study. All comparisons in \cref{tab:model_comparison} are computed on the same five stratified folds (Section~3.3), and the same folds inform every subsequent selection: (i) the fourteen classifiers were compared on cross-validated recall, and Logistic Regression was selected as the base model on the basis of that comparison together with its parsimony; (ii) the preprocessing configuration was selected from the ablation study of \cref{tab:ablation_study}, whose configurations were evaluated on the same folds; (iii) the interpretability analysis, the explanation-quality metrics of Tables~10--11 and the counterfactual evaluation are all computed on a single representative fold (Fold~1), chosen because its Logistic Regression recall lies closest to the mean recall across the five folds; and (iv) the top-$N$ threshold defining the LIME and SHAP subsets was fixed a priori at $N = 12$ and then examined by the sensitivity analysis of Section~5.5, which was likewise run on this fold. The reported performance figures are therefore internal estimates obtained from the very partitions that guided these choices, rather than results from a held-out, selection-free split; they should be read as an optimistic-leaning indication of performance, and confirming them would require a nested or repeated-split protocol in which selection and evaluation use disjoint partitions (Section~6.5).} Before implementing the integrated LSC-XAI framework, we independently applied LIME and SHAP to the model to establish baseline local and global explanations. This stepwise approach allows for a comparative assessment of how the proposed LSC-XAI methodology synthesizes these techniques to provide more robust insights into the drivers of attrition.
\begin{figure}[pos=H]
    \centering
    \includegraphics[width=0.9\linewidth, height=17em]{LIME_output_new.jpg}
    \caption{\textcolor{black}{Local output for LIME is shown for an instance with positive attrition prediction}}
    \label{fig:lime_output}
\end{figure}
\begin{figure}[pos=H]
    \centering
    \includegraphics[width=.8\linewidth, height=20em]{SHAP_output_new.PNG}
    \caption{\textcolor{black}{Global output for SHAP is shown}}
    \label{fig:shap_output}
\end{figure}
From the output in \cref{fig:lime_output}, only local instance's important features and their respective values could be observed. And from the output in \cref{fig:shap_output}, globally important features and their respective significance could be observed. No actionable guidelines for improving the attrition outcome were identified from either method independently.

According to the proposed algorithm of LSC-XAI, at first LIME was implemented. In cases where the model perfectly predicted the true positive attrition outcome, i.e. instances where the model accurately predict if an employee actually left the job, the features contributing to the decisions were recorded and a subset of features had been made. Consequently, SHAP has been implemented to get insight into the features contributing to the explainability on a global level. For SHAP, instances where the model was able to predict both true-positive and true-negative cases were considered. Important features from SHAP have been recorded on a separate subset of features. To reduce computation cost, \textcolor{blue}{a threshold of $N = 12$ features was set a priori and applied to both subsets. This threshold is a fixed design choice rather than a fraction of a single canonical count: it corresponds to approximately one third of the 34 raw attributes of the dataset and to roughly half of the 23 features retained in the Fold-1 model (Section~3.2.4), which is the feature space on which the explainers were computed}. \textcolor{black}{A sensitivity analysis on different thresholds for selecting the number of common features is provided in \cref{sens_analysis}.} 
\begin{table}[width=1\linewidth, pos=H]
    \centering
    \caption{\textcolor{black}{Global important features identified by LIME and SHAP have been shown. Common features identified by both the techniques are separately marked (*)}}
    \label{tab:common_feat}
    \renewcommand{\arraystretch}{1} % Adjust row height
    \begin{tabular}{C|C}
        \hline
        \textbf{\textcolor{black}{Subset of features from LIME}} & \textbf{\textcolor{black}{Subset of features from SHAP}} \\\hline
        \textcolor{black}{MonthlyRate} & \textcolor{black}{OverTime*} \\
        \textcolor{black}{DailyRate} & \textcolor{black}{Shift} \\
        \textcolor{black}{TotalWorkingYears*} & \textcolor{black}{Age} \\
        \textcolor{black}{Shift*} & \textcolor{black}{TotalWorkingYears*} \\
        \textcolor{black}{NumCompaniesWorked*} & \textcolor{black}{JobInvolvement*} \\
        \textcolor{black}{OverTime*} & \textcolor{black}{NumCompaniesWorked*} \\
        \textcolor{black}{JobRole\_Therapist} & \textcolor{black}{JobRole\_Other} \\
        \textcolor{black}{JobInvolvement*} & \textcolor{black}{YearsInCurrentRole*} \\
        \textcolor{black}{EducationField\_Other} & \textcolor{black}{JobRole\_Nurse} \\
        \textcolor{black}{JobRole\_Administrative} & \textcolor{black}{DistanceFromHome*} \\
        \textcolor{black}{YearsInCurrentRole*} & \textcolor{black}{EnvironmentSatisfaction} \\
        \textcolor{black}{DistanceFromHome*} & \textcolor{black}{JobLevel} \\
        \hline
    \end{tabular}
\end{table}

\textcolor{black}{Obtaining both the subsets of important features from the two explainable techniques (LIME and SHAP), common features were marked in \cref{tab:common_feat} and among them, feasible features were selected following the equations (2) \& (3) to undergo the counterfactual technique for suggesting possible changes. Table \ref{tab:feasibility_complete} showcases the markings (1/0) obtained by the common features in accordance to the formal feasibility constraints.}

\begin{table}[width=1\linewidth, pos=H]
\centering
\caption{\textcolor{black}{Feasibility assessment for common features identified by LIME and SHAP. Columns show the three binary feasibility criteria (mutability $\mu$, organizational controllability $\omega$, ethical permissibility $\eta$), the resulting feasibility indicator ($\Psi$), the rationale for each assignment, the permitted
direction/range of modification and its organizational meaning. Features satisfying all three criteria are shown in bold.}}
\label{tab:feasibility_complete}
\begin{tabular}{p{3cm}ccccp{3.25cm}p{2.7cm}p{3.25cm}}
\toprule
\textbf{Feature} & $\mu$ & $\omega$ & $\eta$ & $\Psi$ & \textbf{Rationale} & \textbf{Permitted direction/range} & \textbf{Organizational meaning} \\
\midrule
DistanceFromHome & 1 & 0 & 1 & 0 &
Mutable in principle (employee could relocate) but outside institutional authority &
N/A (excluded) & N/A \\
\midrule
TotalWorkingYears & 0 & 0 & 0 & 0 &
Retrospective, cumulative career history; cannot be altered by any organizational intervention &
N/A (excluded) & N/A \\
\midrule
NumCompaniesWorked & 0 & 0 & 0 & 0 &
Historical employment record; immutable by definition &
N/A (excluded) & N/A \\
\midrule
YearsInCurrentRole & 0 & 0 & 0 & 0 &
Tenure-based attribute; accumulates only with time, not through organizational action &
N/A (excluded) & N/A \\
\midrule
\textbf{JobInvolvement} & \textbf{1} & \textbf{1} & \textbf{1} & \textbf{1} &
Reflects assigned engagement/responsibility level, which is adjustable via task design &
Increase or decrease within ordinal scale (1--4) &
Redistributing task scope or engagement level assigned to the employee \\
\midrule
\textbf{Shift} & \textbf{1} & \textbf{1} & \textbf{1} & \textbf{1} &
Scheduling category set by the organization, not a fixed employee attribute &
Increase or decrease within ordinal scale (0--3) &
Reassignment to a different shift category, not simply more or fewer shifts worked \\
\midrule
\textbf{OverTime} & \textbf{1} & \textbf{1} & \textbf{1} & \textbf{1} &
Directly controlled by scheduling and staffing policy &
Binary: enable or disable overtime requirement &
Adjusting whether overtime is required of the employee \\
\bottomrule
\end{tabular}
\end{table}

The exclusion of 'DistanceFromHome' from $\mathcal{F}$ reflects its failure on organizational controllability alone. Although mutable and ethically permissible, an employee's place of residence lies outside institutional authority, rendering counterfactual recommendations on this feature operationally not-actionable.
'TotalWorkingYears', 'NumCompaniesWorked', 'YearsInCurrentRole' are excluded from $\mathcal{F}$ as all three criteria fail simultaneously. These features represent retrospective, temporally accumulated attributes that are immutable by definition — no organizational intervention can alter an employee's career history, prior employment record, or tenure in a current role. 

\textcolor{black}{To address the inherent subjectivity of these assignments, table \ref{tab:feasibility_complete} provides the complete rationale underlying each feasibility decision, including the permitted direction of change and its organizational interpretation. We note that these assignments reflect author judgment informed by the cited job-characteristics and actionable-recourse literature \cite{ustun2019actionable, karimi2021algorithmic, hackman1976motivation} and have not been validated by healthcare HR professionals, occupational-health specialists, employee representatives, or employment-law experts; this is acknowledged as a limitation in Section \ref{sec:disc}.}

For the feasible features selected:
\begin{itemize}
    \item \textcolor{black}{\textit{\textbf{JobInvolvement}}: Assigned duties to an individual can be increased or decreased}
    \item \textcolor{black}{\textit{\textbf{Shift}}: Working shifts for an individual can be increased or decreased}
    \item \textcolor{black}{\textit{\textbf{OverTime}}: Extra hours of duties for an individual can be increased or decreased}
\end{itemize}

Following the generation of counterfactual explanations based on predefined feasible constraints, the system identified optimal feature subsets and their respective directional shifts. These shifts—categorized as increases or decreases—delineate the specific adjustments required to flip a predicted attrition outcome to a retention state. The frequency and composition of these suggested feature modifications are illustrated in \cref{fig:cf_final}, providing a prioritized road-map for organizational intervention.

\begin{figure}[pos=H]
    \centering
    \includegraphics[width=.8\linewidth, height=20em]{top_12_recomm.png}
    \caption{\textcolor{black}{Frequency count of suggested subset of feature changes are shown. Higher count indicates the concerning subset of features are important for the consideration of changing in the mentioned direction (increase/decrease).}}
    \label{fig:cf_final}
\end{figure}

\subsection{\textcolor{black}{Sensitivity Analysis for Different Feasible Feature Sets}}
\label{sens_analysis}
\textcolor{black}{To evaluate the robustness of the LSC-XAI framework, a sensitivity analysis was performed by varying the parameter $N$, representing the number of top features selected from the LIME and SHAP base explanations. We tested six configurations: $N \in \{32, 16, 14, 12, 10, 8\}$. Illustrated in \cref{fig:sens_analy}, the left sub-figures depict the generated recommendation sets, while the right sub-figures display the feature frequencies across 1000 bootstrap iterations. \textcolor{black}{In the supplementary tables \ref{tab:cf_freq_individual}, \ref{tab:cf_freq_combinations}-- exact recommendation frequency (both individual and combined) and confidence intervals are provided.}  The results indicate that the most significant predictors remain highly stable across varying feature thresholds. For $N \ge 14$, the features Shift[inc], JobLevel[inc] and OverTime[dec] consistently emerged as the primary drivers of attrition. While JobLevel was excluded at $N \le 12$ due to stricter thresholding constraints, the framework maintained high internal consistency. Throughout the experimentation, LSC-XAI repeatedly generated stable recommendation sets—most notably $\{$JobLevel[inc], OverTime[dec]$\}$ and $\{$JobLevel[inc], Shift[inc]$\}$—underscoring the reliability of the proposed methodology in identifying actionable organizational interventions. }    

\begin{figure}[pos=H]
    \centering
    \begin{subfigure}[t]{1\textwidth}
        \centering
        \includegraphics[width=0.9\linewidth, height=25em]{top_all_recomm.png}
        \caption{N = 32 (recommendation sets)}
    \end{subfigure}
    \hfill
    \begin{subfigure}[t]{1\textwidth}
        \centering
        \includegraphics[width=0.9\linewidth, height=20em]{top_all_freq.png}
        \caption{N = 32 (indv. feature frequency)}
    \end{subfigure}
    \hfill
    
    \vspace{0.2cm}
    
    \begin{subfigure}[t]{0.49\textwidth}
        \centering
        \includegraphics[width=0.9\linewidth, height=15em]{top_16_recomm.png}
        \caption{N = 16 (recommendation sets)}
    \end{subfigure}
    \hfill
    \begin{subfigure}[t]{0.49\textwidth}
        \centering
        \includegraphics[width=0.9\linewidth, height=15em]{top_16_freq.png}
        \caption{N = 16 (indv. feature frequency)}
    \end{subfigure}
\end{figure}

\begin{figure}[pos=H]
    \ContinuedFloat
    \centering
    \begin{subfigure}[t]{0.49\textwidth}
        \centering
        \includegraphics[width=0.9\linewidth, height=15em]{top_14_recomm.png}
        \caption{N = 14 (recommendation sets)}
    \end{subfigure}
    \hfill
    \begin{subfigure}[t]{0.49\textwidth}
        \centering
        \includegraphics[width=0.9\linewidth, height=15em]{top_14_freq.png}
        \caption{N = 14 (indv. feature frequency)}
    \end{subfigure}
    
    \vspace{0.2cm}

    \begin{subfigure}[t]{0.49\textwidth}
        \centering
        \includegraphics[width=0.9\linewidth, height=15em]{top_12_recomm.png}
        \caption{N = 12 (recommendation sets)}
    \end{subfigure}
    \hfill
    \begin{subfigure}[t]{0.49\textwidth}
        \centering
        \includegraphics[width=0.9\linewidth, height=15em]{top_12_freq.png}
        \caption{N = 12 (indv. feature frequency)}
    \end{subfigure}
    
    \vspace{0.2cm}

    \begin{subfigure}[t]{0.49\textwidth}
        \centering
        \includegraphics[width=0.9\linewidth, height=10em]{top_10_recomm.png}
        \caption{N = 10 (recommendation sets)}
    \end{subfigure}
    \hfill
    \begin{subfigure}[t]{0.49\textwidth}
        \centering
        \includegraphics[width=0.9\linewidth, height=10em]{top_10_freq.png}
        \caption{N = 10 (indv. feature frequency)}
    \end{subfigure}
    
    \vspace{0.2cm}

    \begin{subfigure}[t]{0.49\textwidth}
        \centering
        \includegraphics[width=0.9\linewidth, height=10em]{top_8_recomm.png}
        \caption{N = 8 (recommendation sets)}
    \end{subfigure}
    \hfill
    \begin{subfigure}[t]{0.49\textwidth}
        \centering
        \includegraphics[width=0.9\linewidth, height=10em]{top_8_freq.png}
        \caption{N = 8 (indv. feature frequency)}
    \end{subfigure}
    
    \caption{\textcolor{black}{For various number of thresholds selecting the top features from LIME and SHAP, resultant changes in recommendation sets (left-side) and individual feature's frequency (right-side) is shown.}}
    \label{fig:sens_analy}
\end{figure}

\subsection{\textcolor{black}{Quantitative Evaluation of LSC-XAI Against Standalone XAI Methods}}
\label{subsec:quant_eval_lsc_xai}
\textcolor{black}{To address RQ3, LSC-XAI is evaluated against standalone LIME \cite{ribeiro2016should}, SHAP \cite{lundberg2017shap} and counterfactual explanations \cite{wachter2017counterfactual} across four established XAI evaluation metrics \cite{nauta2023evaluationsurvey}: explanation stability, surrogate fidelity, feature sparsity and counterfactual quality (\cref{tab:xai_metrics}, \ref{tab:cf_metrics}).}

\begin{table}[pos=H]
\caption{\textcolor{black}{Quantitative evaluation of explanation methods across stability, fidelity, and sparsity, with 95\% bootstrap confidence intervals, computed using the Fold~1 model and preprocessed data
used throughout the interpretability analysis in Section~5.4. Stability and sparsity are computed over $n=31$ true-positive test instances; fidelity is computed over the full held-out test set ($n=336$). Zero-width intervals for LIME/SHAP sparsity reflect that both methods return exactly
the top-$N=12$ features specified a priori for every instance, so sparsity for these two methods is a fixed quantity rather than an empirical distribution. \textcolor{blue}{Sparsity is computed as $1 - (\text{number of selected features} / 23)$, where 23 is the number of features retained by the fold-wise selection in the Fold-1 training fold (Section~3.2.4); this gives $1 - 12/23 = 0.4783$ for LIME and SHAP, and $1 - 3/23 = 0.8696$ for LSC-XAI.}}}
\label{tab:xai_metrics}
\renewcommand{\arraystretch}{1.5}
\begin{tabular}{lccc}
\hline
\textbf{Method} & \textbf{Stability} & \textbf{Fidelity} & \textbf{Sparsity} \\
\hline
LIME     & 0.8034 [0.7867, 0.8201] & 0.9375 [0.9107, 0.9613] & 0.4783 [0.4783, 0.4783] \\
SHAP     & 1.0000 [1.0000, 1.0000] & 0.9375 [0.9107, 0.9613] & 0.4783 [0.4783, 0.4783] \\
LSC-XAI  & 1.0000 [1.0000, 1.0000] & 0.9405 [0.9137, 0.9643] & 0.8696 [0.8696, 0.8696] \\
\hline
\end{tabular}
\end{table}

\textcolor{black}{Stability, measured as mean pairwise Jaccard similarity across 10 repeated runs on the same instances, showed that SHAP and LSC-XAI both reached ceiling performance (Jaccard = 1.0000), while LIME was lower and showed a narrow but distinguishable confidence interval (0.8034 [0.7867, 0.8201]). \textcolor{blue}{For transparency, we report the explainer configuration and its determinism properties. SHAP values were computed with the shap.Explainer API (shap v0.52.0) applied to the trained Logistic Regression's predict function, with the fold's processed test set as the background sample. Because the model is supplied to SHAP as a black-box predict function rather than as a linear model object, the API's automatic selection resolves to its model-agnostic Permutation SHAP approximation — an estimator based on randomly ordered feature permutations — rather than the analytically exact Linear SHAP explainer; this sampling component is stochastic and no random seed was fixed. Attribution rankings for a linear base model are nevertheless highly stable at the library's automatic evaluation budget, and the zero-width interval indicates that no top-12 ranking flip occurred across the ten runs. The SHAP and LSC-XAI stability values should therefore be read as the expected behavior of (near-)deterministic attributions for a linear base model rather than as empirical findings: an exact linear explainer would be deterministic by construction, and the sampling-based configuration used here exhibited no ranking variation in practice. LSC-XAI inherits this ceiling through the intersection — its feature set is recomputed in each run from the LIME and SHAP outputs, and because the three feasible features dominate every LIME ranking, the intersected, feasibility-filtered set was identical in every run. LIME, whose perturbation sampling is likewise unseeded, is the only method whose stability value (0.8034) carries empirical information; the stability comparisons for SHAP and LSC-XAI are accordingly uninformative by construction rather than advantages.}}

\textcolor{black}{Fidelity, assessed as surrogate agreement with the Logistic Regression base model on the held-out test set (n = 336), showed overlapping confidence intervals across all three methods (LIME 0.9375 [0.9107, 0.9613]; SHAP 0.9375 [0.9107, 0.9613]; LSC-XAI 0.9405 [0.9137, 0.9643]). \textcolor{blue}{We define fidelity formally. For each method, the surrogate's feature subset is the method's top-12 features (LIME: the features with the highest aggregated local importance across the evaluated instances, obtained from its perturbation-based local sampling; SHAP: the features with the highest mean absolute attribution). A fresh Logistic Regression surrogate (scikit-learn, max\_iter = 500, random\_state = 42) is then trained on the training fold restricted to that subset, with the base model's own training-fold predictions as the fitting target, so that the surrogate mimics the base model rather than the ground truth. The agreement statistic is the fraction of the 336 held-out test instances on which the surrogate's predicted class equals the base model's predicted class, and the intervals are 95\% bootstrap confidence intervals over these 336 binary agreement outcomes. The identical LIME and SHAP values (315/336) and identical intervals follow mechanically from the linear base model: because both methods recover the same underlying linear decision function, their top-12 feature subsets coincided, so the two surrogates were the same model trained on the same features with the same random seed and produced identical agreement vectors. We agree that the two computations therefore collapse to the same quantity; the identical values are a consequence of the base model's linearity and do not constitute independent corroboration, and near-perfect fidelity is expected for a linear base model rather than informative.} We do not interpret LSC-XAI's marginally higher point estimate as evidence of superior fidelity, since the intervals overlap substantially; the honest reading is that all three methods are statistically indistinguishable in fidelity at this sample size.}

\textcolor{black}{Sparsity, defined as 1 - (|selected features| / |total features|), is higher for LSC-XAI (0.8696) than for either standalone method (0.4783 for both LIME and SHAP, which return exactly the pre-specified top-12 features for every instance). \textcolor{blue}{We emphasize that this difference is a deterministic consequence of the framework's design rather than an empirical finding: LIME and SHAP return exactly twelve features by construction, whereas LSC-XAI's intersection and feasibility filtering retain three feasible features, so the two sparsity values follow arithmetically from these fixed choices (11 of 23 vs. 20 of 23 features of the Fold-1 interpretability model). The sparsity column therefore documents the intended compactness of the recommendation set — a design property, not a measured advantage over standalone methods.}}

% \begin{table}[width=0.6\linewidth,cols=3,pos=H]
% \centering
% \caption{\textcolor{black}{Evaluation of counterfactual quality. \\For Proximity, lower is better ($\downarrow$); for Plausibility, higher is better ($\uparrow$).}}
% \label{tab:cf_metrics}
% \renewcommand{\arraystretch}{1.075}
%     \begin{tabular*}{\tblwidth}{L|C|C}
%     \hline
%     \textbf{\textcolor{black}{Method}} & \textbf{\textcolor{black}{Proximity ($\downarrow$)}} & \textbf{\textcolor{black}{Plausibility ($\uparrow$)}} \\
%     \hline
%     \textcolor{black}{Standalone CF} & \textcolor{black}{7.0268} & \textcolor{black}{0.0769} \\\hline
%     \textcolor{black}{LSC-XAI CF} & \textbf{\textcolor{black}{2.9384}} & \textbf{\textcolor{black}{0.1429}} \\\hline
%     \end{tabular*}
% \end{table}
\begin{table}[pos=H]
\caption{\textcolor{black}{Evaluation of counterfactual quality with 95\% bootstrap confidence intervals. The standalone generator converged to a valid counterfactual for 13 of the 15 attempted instances (86.7\% convergence); LSC-XAI converged for 14 of the 15 attempted instances (93.3\% convergence). For Proximity, lower is better ($\downarrow$); for Plausibility, higher is better ($\uparrow$).}}
\label{tab:cf_metrics}
\renewcommand{\arraystretch}{1.5}
\begin{tabular}{lcc}
\hline
\textbf{Method} & \textbf{Proximity} ($\downarrow$) & \textbf{Plausibility} ($\uparrow$) \\
\hline
Standalone CF ($n=13$) & 6.1206 [5.6895, 6.5294] & 0.2308 [0.0000, 0.4615] \\
LSC-XAI CF ($n=14$)    & 1.2449 [0.9732, 1.5555] & 0.2143 [0.0000, 0.4286] \\
\hline
\end{tabular}
\end{table}

\textcolor{black}{Counterfactual quality was evaluated via proximity and plausibility. \textcolor{blue}{Proximity is defined as the $L_1$ distance between the original instance and the generated counterfactual, $d_1(\mathbf{x}, \mathbf{x}') = \sum_{j} |x'_j - x_j|$, computed on the Z-score standardized feature values (standardized with the training-fold mean and standard deviation), so that each feature's contribution is expressed in units of its training-fold standard deviation; smaller values indicate fewer and smaller changes. Counterfactuals are generated by a random-search procedure: in each of up to 50 attempts, every permitted feature --- any feature for the standalone baseline, only the three feasibility-filtered features for LSC-XAI --- is perturbed by an offset drawn uniformly from $\pm 1.5$ training-fold standard deviations, and the smallest-distance candidate that flips the base model's prediction is retained. Proximity is averaged over the instances for which a valid counterfactual was found: of the 15 attempted instances (drawn from the true-positive evaluation set), the standalone generator converged for 13 (86.7\%) and the LSC-XAI generator for 14 (93.3\%); the remaining attempts did not produce a prediction flip within the 50-attempt budget. The evaluation generator is a lightweight random-search implementation of the same feasibility-restricted design, so that the standalone baseline differs from LSC-XAI only in the permitted feature set; the production recommendations analyzed in Section~5.4 use the DiCE generator described in Section~4. Plausibility is operationalized as a binary indicator per counterfactual: a counterfactual is deemed plausible if every one of its feature values lies within two standard deviations of the training-fold mean ($|z| < 2$ for all features), i.e., if the suggested feature profile remains inside the 2$\sigma$ envelope of the training distribution; plausibility is the fraction of valid counterfactuals satisfying this criterion, with 95\% bootstrap confidence intervals computed over the per-instance indicators.} \textcolor{blue}{LSC-XAI's feasibility-constrained counterfactuals showed smaller proximity (1.2449 [0.9732, 1.5555]) than standalone counterfactuals (6.1206 [5.6895, 6.5294]). We emphasize that smaller proximity is an expected, largely deterministic consequence of the design rather than an empirical advantage: restricting perturbations to the three feasible features necessarily bounds the magnitude of feature modification, so a lower proximity value follows from the restriction itself and cannot be read as evidence of superior counterfactual generation. What the comparison does establish empirically is that this restriction does not degrade plausibility:} standalone (0.2308 [0.0000, 0.4615]) and LSC-XAI (0.2143 [0.0000, 0.4286]) confidence intervals overlap almost entirely, and the standalone point estimate is nominally higher. We therefore do not claim a plausibility advantage for LSC-XAI; this comparison is inconclusive at the current sample size (n = 13–14) and would require a larger counterfactual evaluation set to resolve. \textcolor{blue}{We emphasize that the low absolute plausibility values do not, by themselves, indicate unrealistic counterfactuals, for three reasons. First, the joint 2$\sigma$ criterion is highly conservative: on the full encoded feature set, only 25.5\% of the real training instances themselves satisfy the all-feature 2$\sigma$ envelope (the rate is higher still on the model's retained 23-feature subset), because six of the sixteen binary or one-hot attributes take their minority value at $|z| \ge 2$, so that many genuine employee records fail the criterion automatically. Second, a counterfactual that flips the model's prediction necessarily lies near the class decision boundary — a low-density region of the feature space — so any distributional-envelope criterion will reject a substantial fraction of valid counterfactuals for any generator. Third, for the binary and ordinal attributes that dominate the LSC-XAI modification set (e.g., a shift category or an overtime flag), a Gaussian 2$\sigma$ test is only a crude proxy for categorical rarity. The plausibility metric is therefore interpreted comparatively — as a check that the feasibility restriction does not further degrade the realism of the generated points, which the overlapping intervals confirm — rather than as an absolute certificate of realism. Consistent with this reading, the actionable character of the framework's recommendations rests on which features are modified (organizationally controllable, ethically permissible attributes with permitted directions of change) rather than on the counterfactual point itself being a realistic employee profile, and the recommendations remain hypotheses for HR review and prospective evaluation (Section~6.5).}}

\textcolor{black}{Taken together, these results indicate that LSC-XAI matches SHAP's stability ceiling and shows no measurable fidelity cost relative to either standalone method\textcolor{blue}{, while its higher sparsity and smaller counterfactual proximity follow deterministically from the framework's design choices — the fixed LIME–SHAP intersection, feasibility filtering to three feasible features and perturbation restriction — rather than constituting measured advantages}. LSC-XAI's counterfactuals do not show a distinguishable plausibility advantage over standalone counterfactuals at the current sample size (n = 13–14); this comparison remains open pending a larger evaluation set. \textcolor{blue}{The empirical findings are therefore confined to stability and fidelity — matching SHAP's ceiling at no measurable fidelity cost — whereas the sparsity and proximity figures reflect the intended by-design compactness and boundedness of the recommendations.}}

\subsection{\textcolor{black}{Real-World Validation Experiment}}
\label{real-world}
\textcolor{black}{This experiment evaluates the predictive transfer of the selected classifier under progressive distributional alignment; it does not assess whether the complete LSC-XAI explainability pipeline — the LIME–SHAP intersection, feasibility filtering and counterfactual generation described in Section \ref{sec:lsc_xai} — produces reliable or actionable explanations on real-world data. Extending the full XAI pipeline to the Saudi subset is identified as future work (Section \ref{sec:conclusion}).}
Setup 1 established a cross-distribution baseline by training in the synthetic dataset \cite{KaggleHealthcareAttrition} and testing entirely in the Saudi medical-sector subset. 
Setup 2 introduced partial alignment by augmenting the synthetic training set with 50\% of the Saudi medical-sector records and testing on the remaining 50\% in a stratified-setup. 
Setup 3 applied an 80/20 stratified train-test split exclusively within the Saudi medical-sector subset.
\textcolor{blue}{Two properties of this experiment constrain how the three setups may be compared. First, the three test sets are overlapping subsets of the same 304 Saudi records: Setup 1 evaluates all 304, Setup 2 a single fixed 50\% half (n = 152), and Setup 3 a single fixed 20\% split (n = 61), because all three partitions are drawn with the same random seed; training and test roles interleave correspondingly across setups (Setup 3, for example, trains on records that Setup 2 tests on). Cross-setup comparisons are therefore correlated rather than independent, and apparent differences between setups partly reflect which records were included. Second, each setup is evaluated on a single fixed split with no repeated partitioning: the reported bootstrap confidence intervals quantify sampling variation within that split but not variation across independent data partitions. The Setup 1/2/3 contrasts should therefore be read as descriptive, same-sample comparisons rather than as confirmatory evidence.}

% \begin{table}[width=1\linewidth,cols=6,pos=H]
% \centering
% \caption{\textcolor{black}{Experimental performance across different training setups, comparing Recall for both classes and the overall Macro Recall.}}
% \label{tab:experimental_setups}
% \renewcommand{\arraystretch}{1.2}
%     \begin{tabular*}{\tblwidth}{L|C|C|C|C|C}
%     \hline
%     \textbf{\textcolor{black}{Setup}} & \textbf{\textcolor{black}{Training Data}} & \textbf{\textcolor{black}{Test Data}} & \textbf{\begin{tabular}[c]{@{}c@{}}\textcolor{black}{Recall} \\ \textcolor{black}{(No Attrition)}\end{tabular}} & \textbf{\begin{tabular}[c]{@{}c@{}}\textcolor{black}{Recall} \\ \textcolor{black}{(Attrition)}\end{tabular}} & \textbf{\textcolor{black}{Macro Recall}} \\ \hline
%     \textcolor{black}{1} & \textcolor{black}{Synthetic dataset} & \textcolor{black}{Saudi medical subset (100\%)} & \textcolor{black}{1.00} & \textcolor{black}{0.00} & \textcolor{black}{0.50} \\ \hline
%     \textcolor{black}{2} & \textcolor{black}{Synthetic + 50\% Saudi} & \textcolor{black}{Remaining 50\% Saudi} & \textcolor{black}{0.20} & \textcolor{black}{0.68} & \textcolor{black}{0.44} \\ \hline
%     \textcolor{black}{3} & \textcolor{black}{80\% Saudi medical} & \textcolor{black}{20\% Saudi medical} & \textcolor{black}{0.66} & \textbf{\textcolor{black}{0.70}} & \textbf{\textcolor{black}{0.68}} \\ \hline
%     \end{tabular*}
% \end{table}

\begin{table}[pos=h]
\centering
\caption{\textcolor{black}{Performance metrics across different training configurations}}
\label{tab:performance_metrics}
\resizebox{\textwidth}{!}{
\begin{tabular}{l c c c c c c c c c c c c}
\toprule
\textbf{Setup} & \textbf{n (test)} & \textbf{TP} & \textbf{FP} & \textbf{FN} & \textbf{TN} & \textbf{Precision} & \textbf{Recall (Sens.)} & \textbf{Specificity} & \textbf{F1} & \textbf{Balanced Accuracy} & \textbf{PR-AUC} & \textbf{PR-AUC no-skill} \\ 
\midrule

\begin{tabular}[c]{@{}l@{}}1 \end{tabular} 
& 304 & 62 & 123 & 51 & 68 
& 0.335 [0.266, 0.402]  
& 0.549 [0.458, 0.643] 
& 0.356 [0.290, 0.425] 
& 0.416 [0.341, 0.483] 
& 0.452 [0.397, 0.507] 
& 0.362 [0.291, 0.445] 
& 0.372 \\ 
\midrule

\begin{tabular}[c]{@{}l@{}}2 \end{tabular} 
& 152 & 34 & 74 & 22 & 22 
& 0.315 [0.232, 0.408] 
& 0.607 [0.480, 0.736] 
& 0.229 [0.146, 0.316] 
& 0.415 [0.320, 0.509] 
& 0.418 [0.344, 0.498] 
& 0.355 [0.266, 0.469] 
& 0.368 \\ 
\midrule

\begin{tabular}[c]{@{}l@{}}3 \end{tabular} 
& 61 & 18 & 13 & 5 & 25 
& 0.581 [0.394, 0.750] 
& 0.783 [0.591, 0.945] 
& 0.658 [0.514, 0.800] 
& 0.667 [0.491, 0.807] 
& 0.720 [0.602, 0.833] 
& 0.574 [0.398, 0.783] 
& 0.377 \\ 
\bottomrule
\end{tabular}
}
\end{table}

\textcolor{black}{Table \ref{tab:performance_metrics} reports the full confusion matrix, precision, recall (sensitivity), specificity, F1-score, balanced accuracy and PR-AUC for all three validation setups, together with 95\% bootstrap confidence intervals. In Setup 2, incorporating 50\% of the Saudi sample into training raised attrition recall to 0.607 [0.480, 0.736], compared with 0.549 in the fully cross-domain Setup 1. However, this gain in sensitivity came at a substantial cost to specificity, which fell to 0.229 [0.146, 0.316] meaning the model correctly identified fewer than one in four employees who did not leave. Balanced accuracy, which weights both classes equally, was in fact lower in Setup 2 (0.418 [0.344, 0.498]) than in Setup 1 (0.452 [0.397, 0.507]) and both intervals sit at or below the 0.5 chance baseline. Of the 96 employees in the Setup 2 test set who did not leave, 74 (77.1\%) were incorrectly flagged as attrition risks. This indicates that partial in-domain retraining shifts the model toward over-predicting attrition rather than improving its overall discriminative ability and the earlier framing of Setup 2 as an improvement understated this trade-off. Notably, PR-AUC for both Setup 1 (0.362) and Setup 2 (0.355) was statistically indistinguishable from the no-skill baseline defined by each test set's attrition base rate (0.372 and 0.368, respectively), indicating that under cross-domain transfer the model's ranking of employees by attrition risk carried little information beyond the class prior. By contrast, Setup 3 (Saudi-only training and testing) achieved the strongest performance across all metrics — balanced accuracy 0.720 [0.602, 0.833], specificity 0.658 [0.514, 0.800] and a PR-AUC of 0.574 [0.398, 0.783] that clears its own no-skill baseline (0.377) by a wide margin — suggesting that the feature set itself carries meaningful in-domain signal. \textcolor{blue}{The degradation observed in Setups 1 and 2 reflects a compound of two distinct shifts that the current design cannot separate: a change in the attrition class prior (11.9\% in the training data versus roughly 37\% in the Saudi test sets) and any change in the conditional feature distributions. No threshold recalibration or prior correction was attempted, so the threshold-sensitive results should not be read as evidence of feature-distribution mismatch alone; they establish that the model does not transfer without institution-specific recalibration, not the mechanism of that failure.} \textcolor{blue}{All reported metrics use the default 0.5 decision threshold: \texttt{class\_weight='balanced'} compensates for the training prior during fitting but does not recalibrate the operating point for the 37\% Saudi test prior, so the specificity collapse in Setups 1 and 2 confounds prior shift with any change in the conditional feature distributions. Conversely, the threshold-free PR-AUC --- which measures ranking quality independent of the operating point --- is statistically indistinguishable from the no-skill baseline in Setups 1 and 2, so threshold choice alone cannot explain the degradation either.}}

\subsection{\textcolor{black}{Empirical Fairness Evaluation}}
\textcolor{black}{To empirically address potential algorithmic bias, we evaluated the best-performing predictive baseline across key demographic (Gender, Age Group, Marital Status) and organizational (Department, Job Role) cohorts using the holdout test set ($n = 336$). Algorithmic fairness was quantified using Equal Opportunity Difference ($\Delta\text{TPR}$) and False Positive Rate Difference ($\Delta\text{FPR}$) across sub-cohorts. The subgroup evaluation metrics are detailed in \cref{tab:empirical_fairness_results}.}

\begin{table}[pos=H]
\centering
\caption{\textcolor{black}{Empirical Subgroup Fairness and Error Disparity Analysis Across Demographic and Organizational Cohorts}}
\label{tab:empirical_fairness_results}
% \resizebox{\textwidth}{!}{
\scriptsize
\begin{tabular}{l l c c c c c c}
\toprule
\textbf{Attribute} & \textbf{Subgroup} & \textbf{$n$} & \textbf{Positives ($n_1$)} & \textbf{Recall (TPR)} & \textbf{FPR} & \textbf{$\Delta\text{TPR}$ (Max Diff)} & \textbf{$\Delta\text{FPR}$ (Max Diff)} \\
\midrule
\multirow{2}{*}{Gender} 
& Female & 133 & 17 & 0.765 & 0.103 & \multirow{2}{*}{0.018} & \multirow{2}{*}{0.013} \\
& Male & 203 & 23 & 0.783 & 0.117 & & \\
\midrule
\multirow{4}{*}{Age Group} 
& 18–30 & 83 & 24 & 0.875 & 0.119 & \multirow{4}{*}{0.875} & \multirow{4}{*}{0.110} \\
& 31–40 & 129 & 11 & 0.727 & 0.136 & & \\
& 41–50 & 83 & 3 & 0.667 & 0.113 & & \\
& 51–60 & 41 & 2 & 0.000 & 0.026 & & \\
\midrule
\multirow{3}{*}{Marital Status} 
& Divorced & 74 & 5 & 0.600 & 0.043 & \multirow{3}{*}{0.240} & \multirow{3}{*}{0.129} \\
& Married & 156 & 10 & 0.700 & 0.110 & & \\
& Single & 106 & 25 & 0.840 & 0.173 & & \\
\midrule
\multirow{3}{*}{Department} 
& Cardiology & 98 & 14 & 0.929 & 0.119 & \multirow{3}{*}{0.329} & \multirow{3}{*}{0.011} \\
& Maternity & 169 & 21 & 0.714 & 0.108 & & \\
& Neurology & 69 & 5 & 0.600 & 0.109 & & \\
\midrule
\multirow{5}{*}{Job Role} 
& Administrative & 19 & 0 & N/A\textsuperscript{*} & 0.000 & \multirow{5}{*}{0.150} & \multirow{5}{*}{0.147} \\
& Admin (reference category)\textsuperscript{**} & 4 & 0 & N/A\textsuperscript{*} & 0.000 & & \\
& Nurse & 170 & 20 & 0.850 & 0.147 & & \\
& Other & 111 & 20 & 0.700 & 0.099 & & \\
& Therapist & 32 & 0 & N/A\textsuperscript{*} & 0.063 & & \\
\bottomrule
\multicolumn{8}{l}{\footnotesize \textsuperscript{*}Recall is undefined (N/A) due to zero positive attrition cases ($y=1$) in these specific test cohorts. Positive-case counts ($n_1$) are reported alongside total $n$ for every subgroup. \textsuperscript{**}The `Admin' subgroup is the reference (dropped) level of the one-hot encoding (\texttt{drop\_first=True}); it was previously mislabelled `Base Job Role' in the analysis code, which has been corrected.}
\end{tabular}
% }
\end{table}

\subsubsection*{\textcolor{black}{Demographic Cohort Performance}}
\textcolor{black}{\textbf{Gender:} The model exhibited near-perfect demographic parity across gender groups. Female ($n = 133$) and Male ($n = 203$) cohorts demonstrated comparable Recall ($0.765$ vs. $0.783$) and False Positive Rates ($0.103$ vs. $0.117$), yielding negligible disparities ($\Delta\text{TPR} = 0.018$, $\Delta\text{FPR} = 0.013$). This indicates that the model's sensitivity and error rates remain invariant to gender identity.}

\textcolor{black}{\textbf{Age Groups:} The nominal recall values fall from early-career personnel aged 18–30 ($0.875$, 21 of 24 positives) through the 31–40 ($0.727$, 8 of 11) and 41–50 ($0.667$, 2 of 3) groups, to $0.000$ for the 51–60 cohort (0 of 2 positives; $n = 41$). \textcolor{blue}{These point estimates cannot be read as a trend: the positive-case counts behind them are 24, 11, 3 and 2, so the Wilson 95\% confidence intervals span most of the unit interval for every bracket above 18–30 (31–40: [0.434, 0.903]; 41–50: [0.208, 0.939]; 51–60: [0.000, 0.658], against [0.690, 0.957] for 18–30). The headline disparity ($\Delta\text{TPR} = 0.875$) is therefore dominated by the near-empty 51–60 cell and should be treated as a small-sample instability flag rather than as evidence of age bias. The model's conservatism in the oldest bracket ($\text{FPR} = 0.026$) reflects that, with only two attrition events present, any recall estimate is uninformative.}}

\textcolor{black}{\textbf{Marital Status:} Single personnel ($n = 106$) experienced higher attrition recall ($0.840$) paired with an elevated FPR ($0.173$) relative to Married ($0.700$ Recall; $0.110$ FPR) and Divorced staff ($0.600$ Recall; $0.043$ FPR), generating a recall gap of $\Delta\text{TPR} = 0.240$ and an FPR gap of $\Delta\text{FPR} = 0.129$.}

\textcolor{black}{\subsubsection*{Organizational Cohort Performance}
\textcolor{black}{\textbf{Departmental Disparities:} True positive detection varied across specialized medical divisions. Cardiology ($n = 98$) achieved the highest sensitivity ($0.929$, 13 of 14 positives), whereas Maternity ($n = 169$) and Neurology ($n = 69$) recorded recalls of $0.714$ (15 of 21) and $0.600$ (3 of 5), respectively ($\Delta\text{TPR} = 0.329$). \textcolor{blue}{These estimates rest on 14, 21 and 5 positive cases respectively, so the Neurology figure in particular should be read with its small denominator in mind.} Conversely, False Positive Rates remained remarkably stable across departments, spanning a narrow range between $0.108$ and $0.119$ ($\Delta\text{FPR} = 0.011$).}

\textcolor{black}{\textbf{Job Roles and Sample Sparsity:} Model sensitivity for Nurses ($n = 170$) reached $0.850$ ($\text{FPR} = 0.147$), while broader roles categorized as 'Other' ($n = 111$) achieved a $0.700$ recall ($\text{FPR} = 0.099$). For specialized administrative, \textcolor{blue}{Admin (the dropped one-hot reference category)} and therapist roles ($n = 19$, $4$ and $32$, respectively), recall was mathematically undefined ($\text{N/A}$) due to the absence of positive attrition instances ($y = 1$) in the evaluation split. This highlights the operational challenge of evaluating sub-cohort equity in highly skewed domain-specific healthcare datasets.}

\section{Discussion, Findings and Limitations}
\label{sec:disc}
\subsection{\textcolor{black}{Predictive Performance: Addressing RQ1 and Contextualization Against Prior Work}}
RQ1 asked how shallow learners, ensemble methods and deep learning models compare in predicting healthcare employee attrition. \textcolor{blue}{The comparisons in this subsection use macro-averaged recall, which weights attrition and non-attrition recall equally, rather than attrition-class sensitivity alone. The results demonstrate that no single category of model achieves a statistically dominant advantage. As detailed in Section~\ref{model-comparison}, the omnibus Friedman test detected a statistically significant difference in recall among the fourteen classifiers ($\chi^2(13) = 43.79$, $p < 0.001$, Kendall's W $= 0.674$), indicating that model choice does affect recall. However, none of the ninety-one Holm-corrected pairwise Wilcoxon comparisons reached significance at $n = 5$ folds, even for pairs with perfectly consistent directional differences (rank-biserial $|r| = 1.0$), because the 5-fold design lacks the statistical power to resolve specific pairs. Consequently, the preference for simpler models articulated below rests on descriptive point estimates and parsimony, and it should not be read as evidence that increasing model complexity yields no gains in recall.}
Among ensemble methods, CatBoost outperformed all competitors across three of the four metrics (\textcolor{black}{accuracy: 91.41\%, precision: 80.11\%, F1-score: 78.57\%}), consistent with its documented advantage in natively handling the categorical richness of HR datasets \cite{prokhorenkova2018catboost}. While Raza et al. \cite{app12136424} reported 93\% accuracy using an optimized Extra Trees Classifier, that margin reflects a more balanced dataset rather than a fundamental model superiority. The deep learning models — MLP, CNN and TabNet did not surpass simpler classifiers, corroborating the established finding that deep architectures do not inherently generalize better on small, structured tabular datasets. This challenges the assumption implicit in several prior studies \cite{app12136424, Akasheh2024} that ensemble or deep architectures are necessary for attrition prediction. \textcolor{blue}{Logistic Regression offers a parsimonious base model with the highest reported mean macro-recall; this descriptive comparison does not establish equivalent minority-class sensitivity or statistical superiority over the other models.}

\subsection{\textcolor{black}{The LSC-XAI Framework: Addressing RQ2 and RQ3}}
\textcolor{black}{RQ2 asked whether a hybrid interpretability framework integrating SHAP, LIME and counterfactual explanations can robustly identify key attrition drivers and generate actionable recommendations. The LSC-XAI framework answers this affirmatively. By identifying overtime workload, shift structure and job involvement as the primary feasible intervention targets consistent with burnout and job characteristics theory \cite{hackman1976motivation} and by aggregating counterfactual modifications into frequency-ranked recommendations, the framework translates black-box model outputs into \textcolor{blue}{candidate organizational actions, whose empirical plausibility profile is discussed in Section~5.6}. The most frequently recommended interventions increasing job level alongside shift adjustments and redistributing rather than reducing workload are grounded in occupational research demonstrating that perceived workload fairness, rather than absolute hours, more strongly predicts turnover intention in clinical populations \cite{LavoieTremblay2010}. \textcolor{blue}{Collinearity pruning and feasibility screening serve different purposes: the former removes redundant predictors before model fitting, whereas the latter assesses candidate explanation features for mutability, organizational controllability and ethical permissibility. In the Fold-1 model, MonthlyIncome was retained and JobLevel was removed by collinearity pruning (Section~3.2.4). MonthlyIncome is not among the features assessed in the reported feasibility table; its absence from the resulting recommendation set should therefore not be interpreted as a feasibility rejection or as evidence that compensation-related interventions are unimportant.}}

\textcolor{black}{RQ3 asked whether the hybrid framework can provide better direction than existing, typically single-technique, XAI approaches. Single-method explainability studies conventionally report either SHAP's global feature rankings or LIME's local approximations in isolation, with no built-in mechanism to check whether the two agree or to translate an importance score into a concrete organizational action. LSC-XAI addresses both limitations. First, requiring convergence across SHAP, LIME and the counterfactual engine functions as an internal consistency check: a feature such as overtime workload is reported with more confidence precisely because it surfaces under all three methods rather than a single one, whereas a single-technique study has no comparable way to flag disagreement between what a feature's importance score suggests and what a local explanation or a minimal-change counterfactual would suggest for the same instance. Second, and more consequentially, the framework does not rest this claim on the number of techniques combined; it quantifies the resulting explanations directly, reporting stability, fidelity and sparsity (Section \ref{subsec:quant_eval_lsc_xai}) rather than presenting interpretability as a qualitative or assumed property, as is common where a single explainability technique is applied without any accompanying evaluation of how consistent or faithful its explanations are. This distinction matters: three techniques agreeing is not evidence of quality unless that agreement, and the fidelity of each explanation to the underlying model, is itself measured. It is this quantitative evaluation, not the mere combination of techniques, that substantiates a "better direction" claim relative to single-technique frameworks and gives HR practitioners a documented basis for how much to trust a given recommendation, rather than a recommendation alone.}

It is important to be precise about what these recommendations represent. The counterfactual modifications identified in Section \ref{subsec:lsc_res} are model-generated: they describe changes to feature values that would alter the trained classifier's prediction for a given instance, not changes that have been shown to causally reduce real-world attrition. We therefore distinguish three levels of claim. First, model counterfactuals—the direct output of LSC-XAI indicate which feature changes flip the model's decision boundary. Second, candidate managerial actions are the operational translations of these counterfactuals (e.g., adjusting shift allocation or overtime policy) and require domain review of feasibility, cost and organizational constraints before adoption. Third, causally validated interventions would require prospective, controlled evaluation demonstrating that implementing a given action actually reduces attrition-evidence this study does not provide. For example, a counterfactual suggesting an increase in 'JobLevel' does not establish that promoting an employee will retain them; it only establishes that, holding other features fixed, a higher 'JobLevel' value is associated with a lower predicted attrition probability. The recommendations produced by LSC-XAI should accordingly be treated as hypotheses for HR review and pilot testing, not as validated retention interventions.

\textcolor{black}{The high false-positive rate observed in \cref{tab:performance_metrics} Setup 2 (77.1\% of non-attriting employees incorrectly flagged) carries practical consequences beyond the reported metrics. In a healthcare staffing context, an attrition-risk flag is unlikely to be actioned silently; it typically triggers some combination of manager check-ins, retention conversations, or targeted interventions. When applied to three-quarters of employees who were never at risk of leaving, this produces several downstream costs. First, it imposes an unnecessary administrative burden on HR and departmental managers, who must review and ultimately dismiss a large volume of false alerts, diluting attention available for genuine attrition risk cases. Second, employees who are flagged may experience unwarranted concern about their job security or how they are perceived by management, independent of whether they were ever considering leaving. Third, repeated or visible flagging risks stigmatization— an employee identified as an "attrition risk" may be treated differently in scheduling, promotion consideration, or informal team dynamics, a risk that is heightened in smaller departments or close-knit clinical teams where such information is difficult to keep confidential. Taken together, these costs suggest that a specificity of 0.229, as observed in Setup 2, would likely be unacceptable for deployment even though the corresponding gain in recall might appear attractive in isolation. Any operational deployment of this model would need to weigh the cost of a missed attrition case against the cumulative organizational and human cost of over-flagging and the reported balanced accuracy and specificity not recall alone — should be the metrics used to judge whether a given setup is fit for that purpose.}

\subsection{\textcolor{black}{Real-World Validation: Generalizability and Its Boundaries}}
\textcolor{black}{The real-world validation experiment on the Saudi medical-sector subset \cite{alqahtani2025saudi} surfaces a critical generalization concern applicable to the broader attrition prediction literature. \textcolor{blue}{Under fully cross-domain transfer (Setup 1), attrition-class recall was 0.549 [0.458, 0.643] with balanced accuracy of 0.452 [0.397, 0.507] — at or below the 0.5 chance level — and the model's ranking of employees by attrition risk carried no information beyond the class prior (PR-AUC 0.362 vs. a no-skill baseline of 0.372), confirming that models trained exclusively on synthetic or proxy datasets — a widespread practice given the scarcity of real-world HR data — cannot be directly deployed without domain-specific recalibration. This finding issues a methodological caution for all prior studies reporting high accuracy on held-out IBM HR data without external validation. With 50\% real-world augmentation (Setup 2), attrition-class recall rose to 0.607 [0.480, 0.736], but at a substantial cost to specificity (0.229 [0.146, 0.316]); balanced accuracy (0.418 [0.344, 0.498]) did not improve over Setup 1, so partial augmentation alone did not restore discriminative ability. Only fully in-distribution training (Setup 3) recovered genuine discriminative performance (attrition-class recall 0.783 [0.591, 0.945]; balanced accuracy 0.720 [0.602, 0.833]; PR-AUC 0.574 [0.398, 0.783] vs. a no-skill baseline of 0.377), indicating that effective re-calibration requires institution-specific training data rather than augmentation alone} — consistent with the recommendation of Di Caprio et al. \cite{DICAPRIO2025100403} that retention models must be grounded in institution-specific data.}

\subsection{\textcolor{black}{Ethical Considerations and Algorithmic Governance}}
\textcolor{black}{While the predictive power of the LSC-XAI framework offers significant utility for organizational stability, we must acknowledge the profound ethical responsibility inherent in modeling employee behavior. Using AI to assess attrition risk is not a purely technical exercise; it directly impacts human livelihoods. Consequently, this study addresses the ethical dimensions of our work through three critical lenses: fairness, intent and privacy.\\
\textit{\textbf{Mitigating Bias and Indirect Discrimination:}}
Machine learning models are susceptible to "proxy discrimination," where seemingly neutral variables (e.g., commute distance or marital status) inadvertently correlate with protected characteristics like gender or age. We contend that the "black-box" nature of traditional AI exacerbates this risk. By utilizing the LSC-XAI framework, we provide a mechanism for algorithmic auditing. Transparency allows decision-makers to scrutinize the rationale behind a high-risk score, ensuring that retention strategies are based on professional merit rather than systemic biases embedded in the data.\\
\textit{\textbf{The "Actionable Recourse" Paradigm:}}
A central ethical concern in attrition modeling is the potential for "pre-emptive" adverse actions, such as bypassing an employee for promotion because they are flagged as a "flight risk." We advocate for an Appreciative Inquiry approach: our framework is designed to facilitate supportive interventions, not punitive measures. The generated counterfactuals (e.g., suggesting a workload adjustment rather than a salary change) serve as a roadmap for actionable recourse, empowering organizations to improve the employee experience rather than penalizing individuals for predicted behaviors.\\
\textit{\textbf{Data Privacy and Human Oversight:}}
\textcolor{blue}{This study analyses only publicly released secondary datasets that contain no direct identifiers: the primary dataset is synthetic, derived from the IBM Watson HR Analytics data and released on Kaggle (Section~3.1), and the real-world validation set is a public survey release distributed through Mendeley Data (Section~5.7). Both were analysed in the de-identified form in which their providers released them; the authors neither collected data from human participants nor accessed identifiable employee records, and no attempt was made to re-identify any individual. No formal institutional determination of exemption or of whether ethics review was required has been obtained for this secondary analysis; public availability alone is not taken to establish exemption (see the Ethics Approval declaration).} Furthermore, we emphasize that this framework should function as a Human-in-the-Loop (HITL) system. AI-driven insights should serve as a consultative starting point for HR professionals, ensuring that the final "judgment call" remains a human one, grounded in empathy and context that data alone cannot capture.}\\
\textcolor{black}{\textit{\textbf{Ethical Implications and Auditing Boundaries}}
While the integration of post-hoc interpretability techniques (LIME, SHAP, counterfactuals) within the LSC-XAI framework provides actionable recourse for decision-makers, algorithmic transparency must not be conflated with algorithmic fairness. Explainability methods illuminate how a predictive model synthesizes input features into a decision score, but they do not inherently guarantee equitable predictions or eliminate underlying systematic bias.
As revealed by our empirical fairness audit, an interpretable model can maintain high surrogate fidelity while simultaneously exhibiting disparate false-positive rates across marital statuses or reduced recall in specific age brackets. Therefore, transparency should be viewed as an auditing mechanism rather than a surrogate for fairness. Explainability enables domain experts to inspect model reasoning for proxy discrimination—such as overtime hours acting as an uncalibrated proxy for parental or marital status—thereby informing targeted human oversight before automated predictions influence employment outcomes.}

\subsection{Limitations}
\label{subsec:limilations}
\textcolor{black}{
While the LSC-XAI framework demonstrates strong predictive and explanatory performance, several limitations should be considered when interpreting these results.\\
\textbf{Data provenance and scale.} The primary dataset is synthetically derived from the IBM Watson HR Analytics dataset and does not represent records of real healthcare employees (Section 3.1); results on this dataset should be read as a proof-of-concept demonstration of the framework's methodological capabilities rather than as evidence of real-world predictive accuracy. The real-world validation set is comparatively small and geographically specific (304 medical-sector records from a single-country survey; Section 5.7), and \textcolor{blue}{the near-chance discriminative performance observed in Setup~1 (attrition-class recall of 0.549, with a PR-AUC of 0.362 statistically indistinguishable from the no-skill baseline of 0.372 and balanced accuracy of 0.452 at or below chance)} indicates that findings from the synthetic dataset do not transfer without institution-specific recalibration. \textcolor{blue}{The 11.9\%-to-37\% shift in the attrition class prior between the training data and the Saudi test sets further confounds this transfer gap: without threshold recalibration or prior correction, operating-point metrics (recall, specificity and balanced accuracy) cannot separate the effect of the prior shift from the effect of any change in the conditional feature distributions, so the failure should not be attributed to feature mismatch alone.} \textcolor{blue}{Each setup further uses a single fixed data split, and the three test sets are overlapping subsets of the same 304 records (Setup 1: 100\%; Setup 2: a fixed 50\% half; Setup 3: a fixed 20\% split), so cross-setup comparisons are correlated rather than independent and should be read as provisional (Section~5.7).}\\
\textbf{No prospective, longitudinal, or causal validation.} All evaluations in this study are retrospective and cross-sectional. The framework has not been assessed prospectively against subsequent employee turnover, and no causal claims are supported by the current design. As discussed in Section~6.2, the counterfactual modifications generated by LSC-XAI indicate which feature changes would alter the trained classifier's prediction, not which organizational actions would causally reduce real-world attrition; the recommendations should be treated as hypotheses for HR review and prospective evaluation, not as validated interventions.\\
\textbf{Statistical power.} The 5-fold cross-validation design used for model comparison has limited power to resolve pairwise differences among the fourteen classifiers. As shown in Section~5.2, the omnibus Friedman test indicates that real performance differences exist (p~$<$~0.001), but no individual pairwise comparison can survive Holm correction at $n=5$ folds, even for pairs with a maximal rank-biserial effect size. Conclusions about relative model performance should therefore be read as indicative rather than statistically conclusive
at the pairwise level.\\
\textbf{Model selection protocol and absence of a validation split.} \textcolor{blue}{The study contains no validation partition distinct from the cross-validation folds whose metrics are reported. Model selection (Logistic Regression as the base model), selection of the preprocessing configuration, selection of the representative fold used for the interpretability analysis, and the choice of the top-$N$ feature threshold all draw on the same five stratified folds. Reported performance is consequently not from a held-out, selection-free split: it is an internal estimate that may be optimistic, most acutely for the single-fold (Fold-1) results of Sections~5.4--5.6. Confirming framework performance would require nested cross-validation, or a repeated-split protocol in which selection and evaluation use disjoint partitions.}\\\\
\textbf{Preprocessing and explanation instability.} Fold-wise feature selection (Section~3.2.4) shows only moderate agreement across folds (mean pairwise Jaccard similarity of removed-feature sets $=\textcolor{blue}{0.45}$), indicating that a subset of features are borderline and that the reported feature
set is somewhat sensitive to which data are used for selection. \textcolor{blue}{Similarly, standalone LIME explanations show lower, though still reasonably high, stability across repeated runs (mean pairwise Jaccard $= 0.8034$ [0.7867, 0.8201]; Section~5.6); LSC-XAI's intersection with SHAP reaches the ceiling on this evaluation sample (Jaccard $= 1.0000$), fully absorbing LIME's residual run-to-run variance; nevertheless, explanation stability should not be assumed to be perfect at the individual-feature level.}\\
\textbf{Feasibility framework validation.} The mutability, organizational controllability, and ethical permissibility judgments underlying the feasibility-constrained counterfactual generation (\cref{tab:feasibility_complete}, Section~5.4) reflect author judgment informed by the cited job characteristics and actionable-recourse literature. They have not been validated by healthcare HR professionals, occupational-health specialists, employee representatives, or employment-law experts, and this absence of domain-expert review is a limitation of the current feasibility assignments.\\
\textbf{Counterfactual plausibility.} \textcolor{blue}{The plausibility rates of the generated counterfactuals are low in absolute terms (0.21--0.23; Section~5.6), reflecting a deliberately conservative operationalization (all feature values within 2$\sigma$ of the training-fold mean) under which only about a quarter of real training instances would themselves qualify, and under which any prediction-flipping counterfactual --- which must lie near the class decision boundary --- is penalized. The plausibility figures should therefore be read comparatively rather than as absolute realism certificates, and they reinforce the positioning of the generated recommendations as model-based hypotheses for HR review and prospective evaluation rather than deployable interventions.}\\
\textbf{Fairness auditing scope.} The empirical fairness evaluation reported in Section~5.8 is limited to a single held-out test split of the primary (synthetic) dataset, using Equal Opportunity Difference and False Positive Rate Difference across five attributes. It does not constitute continuous fairness monitoring, was not repeated on the real-world Saudi subset, and several subgroups (e.g., older age brackets, specialized job roles) had too few positive attrition cases for recall to be defined, limiting the granularity of the audit. \textcolor{blue}{A `Base Job Role' label appearing in an earlier draft of the audit was an encoding artefact (the dropped `Admin' reference level of the one-hot encoding) and has been corrected to `Admin' throughout; it denotes four test-fold records, all non-attriting.} \textcolor{blue}{In particular, the headline age disparity ($\Delta\text{TPR} = 0.875$) is driven by a 51–60 cohort containing only two positive attrition cases (recall 0/2) against 24 in the 18–30 cohort, with Wilson 95\% confidence intervals of [0.000, 0.658] and [0.690, 0.957] respectively; it reflects small-sample instability rather than demonstrated age bias.}\\
\textbf{Risks of employee-level risk scoring.} As discussed in Section~6.2, deploying an
attrition-risk classifier at the individual employee level carries operational risk independent of model accuracy: false positives (Section~5.7, Setup~2) can generate unnecessary managerial interventions, employee concern, and stigmatization, particularly in small or close-knit clinical teams where flagged status is difficult to keep confidential.\\
\textbf{Generalizability.} The framework was developed and validated exclusively in a healthcare
staffing context. While its modular design may be adaptable to other sectors, this has not been tested, and cross-institution or cross-industry generalizability should be treated as an open question rather than an established property of the framework (Section~7).
}

\section{Conclusion and Future Work}
\label{sec:conclusion}
In this study, we presented LSC-XAI, a novel hybrid explainable AI framework that synergistically integrates LIME, SHAP and counterfactual techniques to identify actionable feature subsets for organizational decision-making. Leveraging a healthcare employee attrition dataset, we employed rigorous preprocessing steps—including feature selection via majority voting, distributional transformation and upsampling—to ensure data quality. Machine learning models, encompassing shallow, ensemble and deep learning approaches, were evaluated, with the best-performing model selected for explainability analysis. \textcolor{black}{Our proposed LSC-XAI system provided frequency-based, actionable feature recommendations grounded in model counterfactuals; these represent candidate managerial actions for HR review rather than validated retention interventions and their real-world effectiveness would need to be established through prospective evaluation. Relative to single-technique explainability approaches, the framework's cross-method agreement and its quantified stability, fidelity and sparsity give practitioners an evidentiary basis for trusting a recommendation, rather than a recommendation alone.}
The findings of this study carry significant practical implications for policymakers and organizational leaders. \textcolor{blue}{The framework identifies a concise set of features through LIME--SHAP agreement and feasibility screening. These model-derived recommendations are descriptive: no individual-feature hypothesis tests or statistical validation of their intervention effects were performed. They should therefore be treated as candidate actions for HR review and prospective evaluation, not as established means of improving employee satisfaction or reducing attrition.} While the framework was validated in the healthcare sector, its modular design \textcolor{black}{may be adaptable to other sectors, subject to domain-specific validation} facing similar retention challenges. Additionally, the integration of explainability techniques makes AI-driven insights more interpretable, fostering trust and transparency in organizational decision-making.

Future research should explore scaling LSC-XAI to larger and more complex datasets, incorporating more healthcare institutions and global datasets for improved generalization, potentially using federated learning approaches. Relatedly, the current real-world validation (Section \ref{real-world}) evaluated only the predictive transfer of the selected classifier; a natural extension is to apply the complete LSC-XAI pipeline: LIME–SHAP intersection, feasibility filtering and counterfactual generation to real-world institutional data such as the Saudi medical-sector subset and to assess recommendation stability, plausibility and consistency with domain-expert judgment in that setting. \textcolor{black}{The integration of contemporary large language models could enhance textual HR analytics by processing unstructured employee records with interpretability.} Automated feature engineering via meta-learning or reinforcement learning could further optimize feature selection dynamically. Real-world deployment within HR analytics platforms, including interactive dashboards for risk assessment and retention strategy development, would enhance organizational adoption. Additionally, human-in-the-loop AI approaches should be considered to validate AI-driven recommendations with HR professionals for increased trust and effectiveness. Expanding the framework to other industries like education, finance and IT would assess its adaptability across different workforce sectors. Lastly, ensuring fairness and reducing bias in model predictions is critical, requiring continuous refinement of training data and interpretability methods to maintain ethical AI-driven HR solutions. By pursuing these advancements, LSC-XAI can evolve into a scalable and transparent decision-support system that enhances workforce retention across various domains.
\textcolor{black}{Ultimately, the LSC-XAI framework demonstrates that combining robust machine learning algorithms with local, global and counterfactual explanations significantly enhances operational transparency in HR analytics. However, our empirical subgroup audit emphasizes that explainability alone cannot ensure meritocratic outcomes; establishing truly trustworthy decision support systems in healthcare management requires pairing local explanations with continuous, multi-metric empirical fairness evaluations across diverse demographic and institutional cohorts.
}

\section*{Declaration of Interests}
The authors declare no competing interests.

\section*{Funding Statement}
This research received no specific grant from any funding agency.

\section*{Data Availability Statement}
\textcolor{black}{The source code, software environment specifications, model configurations, and the exact dataset used in this study (sourced from Kaggle and Mendeley) have been archived in a persistent public repository to support full reproducibility. The frozen release corresponding to this manuscript is available at https://doi.org/10.5281/zenodo.21594721}
\section*{Ethics Approval}
\textcolor{blue}{This study did not recruit human participants and did not collect primary data from individuals. All analyses were performed on publicly available secondary datasets: (i) the synthetic \emph{Employee Attrition for Healthcare} dataset derived from the IBM Watson HR Analytics data (Kaggle; Section~3.1) and (ii) the publicly released \emph{Saudi Employee Attrition Dataset} (Mendeley Data; Section~5.7). Both datasets were obtained in the de-identified form in which their providers released them; the authors had no access to identifiable records and made no attempt at re-identification. No formal institutional determination of exemption or of whether ethics review was required has been obtained for this secondary analysis, and no applicable institutional policy establishing exemption has been confirmed. We therefore do not claim institutional ethics approval or exemption on the basis of public data availability. Consent and ethics-review arrangements for the original survey collection have not been independently verified for this manuscript, and we do not assert approval or consent on behalf of its providers. The framework is positioned as decision support subject to human oversight rather than automated decision-making, consistent with the governance principles discussed in Section~6.4.}

\printcredits

%% Loading bibliography style file
\bibliographystyle{model1-num-names}
% \bibliographystyle{cas-model2-names}

% Loading bibliography database
\bibliography{cas-refs}


% \appendix
\appendix
\label{sec:appendix}
\renewcommand{\thefigure}{\thesection.\arabic{figure}}
\setcounter{figure}{0}
\section{Supplementary Feature Analysis Plots}
\label{appendix:feature_analysis}

\begin{figure}[pos=H]
    \centering
    \includegraphics[width=0.95\linewidth]{pearson_correlation.png}
    \caption{Pearson correlation values are displayed, highlighting the strong positive collinearity (0.95) between 'MonthlyIncome' and 'JobLevel'.}
    \label{fig:pearson}
\end{figure}
The heatmap in the figure illustrates the linear dependency between all numerical features. It was primarily used to identify multicollinearity, such as the near-perfect correlation between 'JobLevel' and 'MonthlyIncome'.

% \begin{figure}[pos=H]
%     \centering
%     \includegraphics[width=.7\linewidth]{mutual_information.png}
%     \caption{Mutual information for the features against target variable 'Attrition' is shown}
%     \label{fig:mutual_info}
% \end{figure}
% \textcolor{black}{Unlike linear correlation, fig. \ref{fig:mutual_info} captures non-linear relationships between features and the attrition target. It highlights that 'TotalWorkingYears' provide the highest individual predictive power, justifying  inclusion in the final feature set.}

% \begin{figure}[pos=H]
%     \centering
%     \includegraphics[width=0.7\linewidth, height=27em]{ch2_test.png}
%     \caption{Chi-squared test values for the features are shown}
%     \label{fig:chi2}
% \end{figure}
% \textcolor{black}{Fig. \ref{fig:chi2} bar chart displays the results of independence testing for categorical variables. Higher scores for features  confirm strong statistical association with employee turnover.}

% \begin{figure}[pos=H]
%     \centering
%     \includegraphics[width=0.6\linewidth, height=20em]{rfecv.png}
%     \caption{Cross-validation score for RFECV is shown. The highest score is seen when feature number is 8 and 33.}
%     \label{fig:rfecv}
% \end{figure}
% \textcolor{black}{Fig. \ref{fig:rfecv} graph plots the model's accuracy against the number of features selected. The "elbow" of the curve identifies the 33-feature threshold where predictive performance stabilizes, allowing for a more parsimonious model without sacrificing accuracy.}

\begin{figure}[pos=H]
    \centering
        \begin{subfigure}{0.23\textwidth}
            \centering
            \includegraphics[width=\linewidth, height=2cm]{box_total_working_years.png}
            \caption{Total working years}
        \end{subfigure} 
        \begin{subfigure}{0.23\textwidth}
            \centering
            \includegraphics[width=\linewidth, height=2cm]{box_trans_total_working_years.png}
            \caption{Total working years Transformed}
        \end{subfigure} 
        \begin{subfigure}{0.23\textwidth}
            \centering
            \includegraphics[width=\linewidth, height=2cm]{box_years_at_company.png}
            \caption{Years at company}
        \end{subfigure} 
        \begin{subfigure}{0.23\textwidth}
            \centering
            \includegraphics[width=\linewidth, height=2cm]{box_trans_years_at_company.png}
            \caption{Years at company Transformed}
        \end{subfigure}\\
        \begin{subfigure}{0.23\textwidth}
            \centering
            \includegraphics[width=\linewidth, height=2cm]{box_years_in_current_role.png}
            \caption{Years in current role}
        \end{subfigure} 
        \begin{subfigure}{0.23\textwidth}
            \centering
            \includegraphics[width=\linewidth, height=2cm]{box_trans_years_in_current_role.png}
            \caption{Years in current role Transformed}
        \end{subfigure}
        \begin{subfigure}{0.23\textwidth}
            \centering
            \includegraphics[width=\linewidth, height=2cm]{box_years_since_last_promotion.png}
            \caption{Years since last promotion}
        \end{subfigure} 
        \begin{subfigure}{0.23\textwidth}
            \centering
            \includegraphics[width=\linewidth, height=2cm]{box_trans_years_since_last_promotion.png}
            \caption{Years since last promotion Transformed}
        \end{subfigure}\\
    \caption{Distributions before and after applying Box-Cox transformation are shown}
    \label{fig:boxtrans}
\end{figure}
\textcolor{black}{Fig. \ref{fig:boxtrans} shows the PDF (Probability Density Function) of skewed variables before and after transformation. This process was critical for ensuring that distance-based algorithms (like KNN) were not biased by extreme values in features like 'YearsAtCompany'.}

\section{\textcolor{black}{Cross-Validation Feature Selection}}
\label{sec:appendix_features}
\textcolor{black}{The feature selection process was applied uniformly across all evaluated models during the 5-fold stratified cross-validation. Table \ref{tab:removed_features} lists the specific features excluded in each fold to optimize model performance and reduce dimensionality.}


\begin{table}[pos=h]
\footnotesize
    \centering
    \caption{\textcolor{black}{Frequency with which each candidate feature was removed by the fold-wise majority-vote feature selection procedure (including Pearson-based collinearity
pruning) across the 5 training folds ($\checkmark$ = Removed, $\times$ = Retained).}}
    \label{tab:removed_features}
    \begin{tabular}{@{}lcccccc@{}}
        \toprule
        \textbf{Feature (non significant} & \textbf{Fold 1} & \textbf{Fold 2} & \textbf{Fold 3} & \textbf{Fold 4} & \textbf{Fold 5} & \textbf{Folds removed (out of 5)} \\
        \midrule
        Department\_Neurology & $\checkmark$ & $\checkmark$ & $\checkmark$ & $\checkmark$ & $\checkmark$ & 5 \\
        EducationField\_Other & $\checkmark$ & $\checkmark$ & $\checkmark$ & $\checkmark$ & $\checkmark$ & 5 \\
        EducationField\_Technical Degree & $\checkmark$ & $\checkmark$ & $\checkmark$ & $\checkmark$ & $\checkmark$ & 5 \\
        PerformanceRating & $\checkmark$ & $\checkmark$ & $\checkmark$ & $\checkmark$ & $\checkmark$ & 5 \\
        EducationField\_Marketing & $\checkmark$ & $\checkmark$ & $\checkmark$ & $\checkmark$ & $\times$ & 4 \\
        EducationField\_Medical & $\checkmark$ & $\checkmark$ & $\checkmark$ & $\checkmark$ & $\times$ & 4 \\
        Gender & $\checkmark$ & $\checkmark$ & $\checkmark$ & $\times$ & $\checkmark$ & 4 \\
        JobSatisfaction & $\checkmark$ & $\times$ & $\checkmark$ & $\checkmark$ & $\checkmark$ & 4 \\
        Department\_Maternity & $\times$ & $\checkmark$ & $\checkmark$ & $\checkmark$ & $\times$ & 3 \\
        EducationField\_Life Sciences & $\checkmark$ & $\checkmark$ & $\checkmark$ & $\times$ & $\times$ & 3 \\
        EnvironmentSatisfaction & $\checkmark$ & $\checkmark$ & $\times$ & $\checkmark$ & $\times$ & 3 \\
        PercentSalaryHike & $\checkmark$ & $\times$ & $\checkmark$ & $\times$ & $\checkmark$ & 3 \\
        RelationshipSatisfaction & $\checkmark$ & $\checkmark$ & $\checkmark$ & $\times$ & $\times$ & 3 \\
        WorkLifeBalance & $\times$ & $\checkmark$ & $\checkmark$ & $\times$ & $\checkmark$ & 3 \\
        BusinessTravel & $\times$ & $\checkmark$ & $\checkmark$ & $\times$ & $\times$ & 2 \\
        Education & $\times$ & $\checkmark$ & $\checkmark$ & $\times$ & $\times$ & 2 \\
        JobInvolvement & $\times$ & $\times$ & $\checkmark$ & $\checkmark$ & $\times$ & 2 \\
        JobRole\_Administrative & $\times$ & $\times$ & $\checkmark$ & $\checkmark$ & $\times$ & 2 \\
        JobRole\_Nurse & $\times$ & $\checkmark$ & $\checkmark$ & $\times$ & $\times$ & 2 \\
        JobRole\_Other & $\checkmark$ & $\times$ & $\checkmark$ & $\times$ & $\times$ & 2 \\
        MaritalStatus\_Married & $\checkmark$ & $\times$ & $\checkmark$ & $\times$ & $\times$ & 2 \\
        TrainingTimesLastYear & $\checkmark$ & $\times$ & $\checkmark$ & $\times$ & $\times$ & 2 \\
        HourlyRate & $\times$ & $\times$ & $\times$ & $\checkmark$ & $\times$ & 1 \\
        JobRole\_Therapist & $\times$ & $\times$ & $\times$ & $\checkmark$ & $\times$ & 1 \\
        \bottomrule
    \end{tabular}
\end{table}

\section{Results of Confusion Matrices}
\label{appendix:confusion}
\begin{figure}[pos=H]
    \centering
    \begin{subfigure}[t]{0.24\textwidth}
        \centering
        \includegraphics[width=0.9\linewidth]{LR_cm.png}
        \caption{Logistic Regression}
    \end{subfigure}
    \hfill
    \begin{subfigure}[t]{0.24\textwidth}
        \centering
        \includegraphics[width=0.9\linewidth]{svm_cm.png}
        \caption{Support Vector Machines}
    \end{subfigure}
    \hfill
    \begin{subfigure}[t]{0.24\textwidth}
        \centering
        \includegraphics[width=0.9\linewidth]{knn_cm.png}
        \caption{K-Nearest Neighbor}
    \end{subfigure}
    \hfill
    \begin{subfigure}[t]{0.24\textwidth}
        \centering
        \includegraphics[width=0.9\linewidth]{dt_cm.png}
        \caption{Decision Tree}
    \end{subfigure}
    
    \vspace{0.2cm}
    
    \begin{subfigure}[t]{0.24\textwidth}
        \centering
        \includegraphics[width=0.9\linewidth]{rf_cm.png}
        \caption{Random Forest}
    \end{subfigure}
    \hfill
    \begin{subfigure}[t]{0.24\textwidth}
        \centering
        \includegraphics[width=0.9\linewidth]{gb_cm.png}
        \caption{Gradient Boosting}
    \end{subfigure}
    \hfill
    \begin{subfigure}[t]{0.24\textwidth}
        \centering
        \includegraphics[width=0.9\linewidth]{ada_cm.png}
        \caption{Adaptive Boosting}
    \end{subfigure}
    \hfill
    \begin{subfigure}[t]{0.24\textwidth}
        \centering
        \includegraphics[width=0.9\linewidth]{xg_cm.png}
        \caption{Extr. Gradient Boosting}
    \end{subfigure}

    \vspace{0.2cm}
    
    \begin{subfigure}[t]{0.24\textwidth}
        \centering
        \includegraphics[width=0.9\linewidth]{light_cm.png}
        \caption{Light Gradient Boosting}
    \end{subfigure}
    \hfill
    \begin{subfigure}[t]{0.24\textwidth}
        \centering
        \includegraphics[width=0.9\linewidth]{cat_cm.png}
        \caption{Categorical Boosting}
    \end{subfigure}
    \hfill
    \begin{subfigure}[t]{0.24\textwidth}
        \centering
        \includegraphics[width=0.9\linewidth]{stack_cm.png}
        \caption{Stacking}
    \end{subfigure}
    \hfill
    \begin{subfigure}[t]{0.24\textwidth}
        \centering
        \includegraphics[width=0.9\linewidth]{mlp_cm.png}
        \caption{MLP}
    \end{subfigure}
    
    \vspace{0.2cm}
    
    % Centering the last row with two subfigures
    \hspace*{\fill} % Adds stretchable space on the left
    \begin{subfigure}[t]{0.25\textwidth}
        \centering
        \includegraphics[width=0.9\linewidth]{cnn_cm.png}
        \caption{CNN}
    \end{subfigure}
    \hfill
    \begin{subfigure}[t]{0.25\textwidth}
        \centering
        \includegraphics[width=0.9\linewidth]{tab_cm.png}
        \caption{Tab Net}
    \end{subfigure}
    \hspace*{\fill} % Adds stretchable space on the right
    
    \caption{Confusion matrices for the trained models are shown.}
    \label{fig:cm_subfigures}
\end{figure}
\textcolor{black}{Fig. \ref{fig:cm_subfigures} shows the classification results for every model tested across Fold-1. This allows for a granular comparison of Type I errors (false alarms) versus Type II errors (missed attrition cases), the latter being more costly for HR retention strategies.}

\newpage
\section{Evaluation of Receiver Operating Characteristics}
\label{appendix:roc}
\begin{figure}[pos=H]
    \centering
    \begin{subfigure}[t]{0.24\textwidth}
        \centering
        \includegraphics[width=0.9\linewidth]{LR_roc.png}
        \caption{Logistic Regression}
    \end{subfigure}
    \hfill
    \begin{subfigure}[t]{0.24\textwidth}
        \centering
        \includegraphics[width=0.9\linewidth]{svm_roc.png}
        \caption{Support Vector Machines}
    \end{subfigure}
    \hfill
    \begin{subfigure}[t]{0.24\textwidth}
        \centering
        \includegraphics[width=0.9\linewidth]{knn_roc.png}
        \caption{K-Nearest Neighbor}
    \end{subfigure}
    \hfill
    \begin{subfigure}[t]{0.24\textwidth}
        \centering
        \includegraphics[width=0.9\linewidth]{dt_roc.png}
        \caption{Decision Tree}
    \end{subfigure}
    
    \vspace{0.1cm}
    
    \begin{subfigure}[t]{0.24\textwidth}
        \centering
        \includegraphics[width=0.9\linewidth]{rf_roc.png}
        \caption{Random Forest}
    \end{subfigure}
    \hfill
    \begin{subfigure}[t]{0.24\textwidth}
        \centering
        \includegraphics[width=0.9\linewidth]{gb_roc.png}
        \caption{Gradient Boosting}
    \end{subfigure}
    \hfill
    \begin{subfigure}[t]{0.24\textwidth}
        \centering
        \includegraphics[width=0.9\linewidth]{ada_roc.png}
        \caption{Adaptive Boosting}
    \end{subfigure}
    \hfill
    \begin{subfigure}[t]{0.24\textwidth}
        \centering
        \includegraphics[width=0.9\linewidth]{xg_roc.png}
        \caption{Extr. Gradient Boosting}
    \end{subfigure}
    
    \vspace{0.1cm}
    
    \begin{subfigure}[t]{0.24\textwidth}
        \centering
        \includegraphics[width=0.9\linewidth]{light_roc.png}
        \caption{Light Gradient Boosting}
    \end{subfigure}
    \hfill
    \begin{subfigure}[t]{0.24\textwidth}
        \centering
        \includegraphics[width=0.9\linewidth]{cat_roc.png}
        \caption{Categorical Boosting}
    \end{subfigure}
    \hfill
    \begin{subfigure}[t]{0.24\textwidth}
        \centering
        \includegraphics[width=0.9\linewidth]{stack_roc.png}
        \caption{Stacking}
    \end{subfigure}
    \hfill
    \begin{subfigure}[t]{0.24\textwidth}
        \centering
        \includegraphics[width=0.9\linewidth]{mlp_roc.png}
        \caption{MLP}
    \end{subfigure}
    
    \vspace{0.1cm}
    
    % Centering the last row with two subfigures
    \hspace*{\fill}
    \begin{subfigure}[t]{0.25\textwidth}
        \centering
        \includegraphics[width=0.9\linewidth]{cnn_roc.png}
        \caption{CNN}
    \end{subfigure}
    \hfill
    \begin{subfigure}[t]{0.25\textwidth}
        \centering
        \includegraphics[width=0.9\linewidth]{tab_roc.png}
        \caption{Tab Net}
    \end{subfigure}
    \hspace*{\fill}
    
    \caption{ROC-AUC figures for the trained models are shown.}
    \label{fig:roc_subfigures}
\end{figure}
\textcolor{black}{Fig. \ref{fig:roc_subfigures} displays the Receiver Operating Characteristic (ROC) curves for all high-performing models for Fold-1. The proximity of the CatBoost and AdaBoost curves to the top-left corner (AUC = 0.92) visually validates their improved ability to distinguish between classes.}
\newpage

\section{Illustration of Deep Networks}
\begin{figure}[pos=H]
    \centering
    \includegraphics[width=0.7\linewidth, height=100pt]{MLP_diagram.png}
    \caption{Multi-layer perceptron architecture used in this study}
    \label{fig:mlp}
\end{figure}
\textcolor{black}{Fig. \ref{fig:mlp} represents the 4-layer deep learning architecture. It highlights the integration of Batch Normalization and Dropout layers, which were implemented to prevent overfitting on the relatively small dataset.}

\begin{figure}[pos=H]
    \centering
    \includegraphics[width=0.65\linewidth, height=100pt]{CNN_diagram.png}
    \caption{Convolutional neural networks architecture used in this study}
    \label{fig:cnn}
\end{figure}
\textcolor{black}{Fig. \ref{fig:cnn} details the architecture of the 1D-Convolutional Neural Network, showing how 1D filters scan the feature vector to identify local feature interactions. It includes the dimensions for the Max-Pooling and Flattening layers used before the final Softmax output.}

\section{\textcolor{black}{Hyperparameters for Training the Models}}
\begin{longtable}{p{2cm} p{7.75cm} p{3cm} p{2.25cm}}
\captionsetup{justification=raggedright, singlelinecheck=false}
\caption{\textcolor{black}{Software implementations, random seeds, hyperparameters and tuning procedures are described. \textcolor{blue}{Seeds and settings correspond to the frozen release; ``exactly reproducible'' refers to reruns within that archived environment with identical library versions and CPU execution, since bitwise identity across library versions is not guaranteed for neural-network training.}}}
\scriptsize
\label{tab:software_implementations} \\

\toprule
\textbf{Component / Model} & \textbf{Software implementation} & \textbf{Final settings / reproducibility details} & \textbf{Tuning / selection procedure} \\ 
\midrule
\endfirsthead

\multicolumn{4}{c}%
{{\tablename\ \thetable{} -- continued from previous page}} \\
\toprule
\textbf{Component / Model} & \textbf{Software implementation} & \textbf{Final settings / reproducibility details} & \textbf{Tuning / selection procedure} \\ 
\midrule
\endhead

\midrule
\multicolumn{4}{r}{{Continued on next page}} \\
\endfoot

\bottomrule
\endlastfoot

Train-test split 
& \texttt{sklearn.model\_selection.train\_test\_split} 
& 80:20 split; test\_size=0.2; stratify=y; random\_state=42 
& Fixed split before model training \\ 
\midrule

SMOTE-NC 
& \texttt{imblearn.over\_sampling.SMOTENC} 
& Categorical features identified as columns with fewer than 10 unique values; random\_state=42; other settings default 
& Applied only to training data \\ 
\midrule

Logistic Regression 
& \texttt{sklearn.linear\_model.LogisticRegression} 
& random\_state=42; other parameters default 
& No hyperparameter tuning \\ 
\midrule

SVM 
& \texttt{sklearn.svm.SVC} 
& probability=True; random\_state=42; other parameters default 
& No hyperparameter tuning \\ 
\midrule

KNN 
& \texttt{sklearn.neighbors.KNeighborsClassifier} 
& Default settings: n\_neighbors=5, weights='uniform', metric='minkowski' 
& No hyperparameter tuning \\ 
\midrule

Decision Tree 
& \texttt{sklearn.tree.DecisionTreeClassifier} 
& random\_state=42; other parameters default 
& No hyperparameter tuning \\ 
\midrule

Random Forest 
& \texttt{sklearn.ensemble.RandomForestClassifier} 
& random\_state=42; default model uses default settings; stacking base learner uses n\_estimators=100 
& No hyperparameter tuning \\ 
\midrule

Gradient Boosting 
& \texttt{sklearn.ensemble.GradientBoostingClassifier} 
& random\_state=42; other parameters default 
& No hyperparameter tuning \\ 
\midrule

AdaBoost 
& \texttt{sklearn.ensemble.AdaBoostClassifier} 
& random\_state=42; other parameters default 
& No hyperparameter tuning \\ 
\midrule

XGBoost 
& \texttt{xgboost.XGBClassifier} 
& random\_state=42; other parameters default 
& No hyperparameter tuning \\ 
\midrule

LightGBM 
& \texttt{lightgbm.LGBMClassifier} 
& random\_state=42; other parameters default 
& No hyperparameter tuning \\ 
\midrule

CatBoost 
& \texttt{catboost.CatBoostClassifier} 
& verbose=0; random\_state=42; other parameters default 
& No hyperparameter tuning \\ 
\midrule

Stacking Ensemble 
& \texttt{sklearn.ensemble.StackingClassifier} 
& Base learners: Random Forest, CatBoost and SVC; final estimator: Logistic Regression; random-state settings inherited from seeded base/meta models 
& Base learners manually selected from strong individual models \\ 
\midrule

MLP 
& PyTorch 
& Hidden units: 20 and 20; ReLU; BatchNorm; sigmoid output; BCE loss; Adam optimizer; learning rate = 0.0001; batch size = 16; epochs = 30; threshold = 0.5; \textcolor{blue}{fold partition uses \texttt{StratifiedKFold(random\_state=42)}; per-fold PyTorch seed set via \texttt{torch.manual\_seed(42 + fold)} before model construction and training; CPU execution; \texttt{DataLoader} shuffling with \texttt{num\_workers}=0 therefore draws from the seeded global RNG} 
& No additional validation monitoring performed \\ 
\midrule

1D-CNN 
& PyTorch 
& Conv1D filters = 32; kernel size = 4; ReLU; max-pooling size = 2; sigmoid output; BCE loss; Adam optimizer; learning rate = 0.001; batch size = 16; epochs = 30; threshold = 0.5; \textcolor{blue}{fold partition uses \texttt{StratifiedKFold(random\_state=42)}; per-fold PyTorch seed set via \texttt{torch.manual\_seed(42 + fold)} before model construction and training; CPU execution; \texttt{DataLoader} shuffling with \texttt{num\_workers}=0 therefore draws from the seeded global RNG} 
& No additional validation monitoring performed \\ 
\midrule

TabNet 
& \texttt{pytorch\_tabnet.tab\_model.TabNetClassifier} 
& \textcolor{blue}{\texttt{TabNetClassifier(verbose=0, random\_state=42 + fold)}; max\_epochs=30; patience=10; early stopping monitored on the evaluation set supplied at fit time} 
& No explicit hyperparameter search; early stopping used \\ 

\midrule

LIME / SHAP explainers
& \texttt{lime.lime\_tabular.LimeTabularExplainer}; \texttt{shap.Explainer} (shap v0.52.0)
& LIME: perturbation sampling without a fixed random\_state; SHAP: predict-function interface with the processed Fold-1 test set as background, auto-selected Permutation SHAP approximation (sampling-based), no random seed fixed
& Interpretability analysis on the Fold-1 model (Sections 5.4 and 5.6) \\
\midrule

DiCE counterfactual generator
& \texttt{dice-ml} v0.12; \texttt{dice\_ml.Dice(method="random")}, scikit-learn backend
& \texttt{total\_CFs}=3 per instance; \texttt{desired\_class}="opposite"; \texttt{features\_to\_vary} = feasible feature set; generations without a predicted-class flip are skipped
& Recommendation generation for correctly classified attrition cases of the Fold-1 model (Section 5.4) \\

\end{longtable}

\newpage

\section{\textcolor{black}{Pairwise Statistical Significance and Non-Parametric Model Comparisons}}
\label{app:pairwise_statistical}
\textcolor{black}{Pairwise Wilcoxon signed-rank tests were conducted across all evaluated model combinations ($N = 91$) to assess performance significance. Rank-biserial correlation ($r$) was used to measure effect size, and $p$-values were adjusted using the Holm-Bonferroni correction ($p_{\text{adj}}$). As detailed in table A.1, after controlling for multiple comparisons, no pairwise model difference achieved statistical significance ($\alpha = 0.05$, $p_{\text{adj}} = 1.0000$).}
\begingroup
\captionsetup[longtable]{justification=raggedright, singlelinecheck=false}
\small

\begin{longtable}{l l c c c c}
\caption{\textcolor{black}{Pairwise Non-Parametric Model Comparisons and Statistical Significance Results ($N = 91$)}} 
\label{tab:app_statistical_comparisons} \\
\toprule
\textbf{Model 1} & \textbf{Model 2} & \textbf{Rank-biserial $r$} & \textbf{$p_{\text{raw}}$} & \textbf{$p_{\text{adj}}$} & \textbf{Significant} \\ 
\midrule
\endfirsthead

\multicolumn{6}{l}{\textit{-- Continued from previous page}} \\
\toprule
\textbf{Model 1} & \textbf{Model 2} & \textbf{Rank-biserial $r$} & \textbf{$p_{\text{raw}}$} & \textbf{$p_{\text{adj}}$} & \textbf{Significant} \\ 
\midrule
\endhead

\midrule
\multicolumn{6}{r}{\textit{Continued on next page}} \\
\endfoot

\bottomrule
\endlastfoot

AdaBoost            & Logistic Regression & -1.0000 & 0.0625 & 1.0000 & No \\
CNN                 & TabNet              &  1.0000 & 0.0625 & 1.0000 & No \\
CatBoost            & Decision Tree       &  1.0000 & 0.0625 & 1.0000 & No \\
CNN                 & Random Forest       &  1.0000 & 0.0625 & 1.0000 & No \\
CNN                 & Stacking            &  1.0000 & 0.0625 & 1.0000 & No \\
CNN                 & KNN                 &  1.0000 & 0.0625 & 1.0000 & No \\
CNN                 & Gradient Boosting   &  1.0000 & 0.0625 & 1.0000 & No \\
CNN                 & Decision Tree       &  1.0000 & 0.0625 & 1.0000 & No \\
CNN                 & CatBoost            &  1.0000 & 0.0625 & 1.0000 & No \\
AdaBoost            & XGBoost             &  1.0000 & 0.0625 & 1.0000 & No \\
AdaBoost            & TabNet              &  1.0000 & 0.0625 & 1.0000 & No \\
CNN                 & XGBoost             &  1.0000 & 0.0625 & 1.0000 & No \\
LightGBM            & Logistic Regression & -1.0000 & 0.0625 & 1.0000 & No \\
Logistic Regression & TabNet              &  1.0000 & 0.0625 & 1.0000 & No \\
Gradient Boosting   & Logistic Regression & -1.0000 & 0.0625 & 1.0000 & No \\
KNN                 & Logistic Regression & -1.0000 & 0.0625 & 1.0000 & No \\
KNN                 & MLP                 & -1.0000 & 0.0625 & 1.0000 & No \\
Gradient Boosting   & MLP                 & -1.0000 & 0.0625 & 1.0000 & No \\
Gradient Boosting   & TabNet              &  1.0000 & 0.0625 & 1.0000 & No \\
Decision Tree       & Stacking            & -1.0000 & 0.0625 & 1.0000 & No \\
Decision Tree       & Random Forest       & -1.0000 & 0.0625 & 1.0000 & No \\
Decision Tree       & MLP                 & -1.0000 & 0.0625 & 1.0000 & No \\
Decision Tree       & Logistic Regression & -1.0000 & 0.0625 & 1.0000 & No \\
Decision Tree       & SVM                 & -1.0000 & 0.0625 & 1.0000 & No \\
CatBoost            & TabNet              &  1.0000 & 0.0625 & 1.0000 & No \\
CatBoost            & Logistic Regression & -1.0000 & 0.0625 & 1.0000 & No \\
CatBoost            & MLP                 & -1.0000 & 0.0625 & 1.0000 & No \\
TabNet              & XGBoost             & -1.0000 & 0.0625 & 1.0000 & No \\
MLP                 & Stacking            &  1.0000 & 0.0625 & 1.0000 & No \\
Logistic Regression & XGBoost             &  1.0000 & 0.0625 & 1.0000 & No \\
MLP                 & TabNet              &  1.0000 & 0.0625 & 1.0000 & No \\
Random Forest       & SVM                 & -1.0000 & 0.0625 & 1.0000 & No \\
Stacking            & TabNet              &  1.0000 & 0.0625 & 1.0000 & No \\
SVM                 & TabNet              &  1.0000 & 0.0625 & 1.0000 & No \\
Random Forest       & TabNet              &  1.0000 & 0.0625 & 1.0000 & No \\
Logistic Regression & Stacking            &  1.0000 & 0.0625 & 1.0000 & No \\
SVM                 & XGBoost             &  1.0000 & 0.0625 & 1.0000 & No \\
MLP                 & XGBoost             &  1.0000 & 0.0625 & 1.0000 & No \\
Logistic Regression & Random Forest       &  1.0000 & 0.0625 & 1.0000 & No \\
LightGBM            & SVM                 & -1.0000 & 0.0625 & 1.0000 & No \\
CNN                 & LightGBM            &  0.8667 & 0.1250 & 1.0000 & No \\
AdaBoost            & Gradient Boosting   &  0.8667 & 0.1250 & 1.0000 & No \\
AdaBoost            & Decision Tree       &  0.8667 & 0.1250 & 1.0000 & No \\
KNN                 & Stacking            & -0.8667 & 0.1250 & 1.0000 & No \\
KNN                 & SVM                 & -0.8667 & 0.1250 & 1.0000 & No \\
MLP                 & Random Forest       &  0.8667 & 0.1250 & 1.0000 & No \\
SVM                 & Stacking            &  0.8667 & 0.1250 & 1.0000 & No \\
Stacking            & XGBoost             &  0.8667 & 0.1250 & 1.0000 & No \\
LightGBM            & TabNet              &  0.8667 & 0.1250 & 1.0000 & No \\
Decision Tree       & XGBoost             & -0.8667 & 0.1250 & 1.0000 & No \\
LightGBM            & MLP                 & -0.8667 & 0.1250 & 1.0000 & No \\
Gradient Boosting   & KNN                 &  0.8667 & 0.1250 & 1.0000 & No \\
CatBoost            & SVM                 & -0.8667 & 0.1250 & 1.0000 & No \\
Decision Tree       & Gradient Boosting   & -0.8667 & 0.1250 & 1.0000 & No \\
AdaBoost            & KNN                 &  0.8667 & 0.1250 & 1.0000 & No \\
Gradient Boosting   & SVM                 & -0.8667 & 0.1250 & 1.0000 & No \\
CatBoost            & KNN                 &  0.8667 & 0.1250 & 1.0000 & No \\
AdaBoost            & MLP                 & -0.7333 & 0.1875 & 1.0000 & No \\
KNN                 & Random Forest       & -0.7333 & 0.1875 & 1.0000 & No \\
KNN                 & TabNet              &  0.7333 & 0.1875 & 1.0000 & No \\
Gradient Boosting   & XGBoost             &  0.7333 & 0.1875 & 1.0000 & No \\
CatBoost            & Stacking            & -0.7333 & 0.1875 & 1.0000 & No \\
Decision Tree       & LightGBM            & -0.7333 & 0.1875 & 1.0000 & No \\
Decision Tree       & TabNet              &  0.6000 & 0.3125 & 1.0000 & No \\
KNN                 & XGBoost             & -0.6000 & 0.3125 & 1.0000 & No \\
AdaBoost            & LightGBM            &  0.6000 & 0.3125 & 1.0000 & No \\
CNN                 & Logistic Regression & -0.6000 & 0.3125 & 1.0000 & No \\
AdaBoost            & CatBoost            &  0.6000 & 0.3125 & 1.0000 & No \\
AdaBoost            & CNN                 & -0.6000 & 0.3125 & 1.0000 & No \\
CatBoost            & XGBoost             &  0.6000 & 0.3125 & 1.0000 & No \\
Gradient Boosting   & Stacking            & -0.6000 & 0.3125 & 1.0000 & No \\
AdaBoost            & Random Forest       &  0.6000 & 0.3125 & 1.0000 & No \\
AdaBoost            & SVM                 & -0.4667 & 0.4375 & 1.0000 & No \\
Random Forest       & XGBoost             &  0.4667 & 0.4375 & 1.0000 & No \\
Random Forest       & Stacking            & -0.4667 & 0.4375 & 1.0000 & No \\
KNN                 & LightGBM            & -0.4667 & 0.4375 & 1.0000 & No \\
LightGBM            & XGBoost             &  0.4000 & 0.6250 & 1.0000 & No \\
Logistic Regression & SVM                 &  0.3333 & 0.6250 & 1.0000 & No \\
Gradient Boosting   & Random Forest       & -0.3333 & 0.6250 & 1.0000 & No \\
Logistic Regression & MLP                 &  0.3333 & 0.6250 & 1.0000 & No \\
MLP                 & SVM                 & -0.3333 & 0.6250 & 1.0000 & No \\
CatBoost            & Random Forest       & -0.3333 & 0.6250 & 1.0000 & No \\
AdaBoost            & Stacking            &  0.2000 & 0.8125 & 1.0000 & No \\
CNN                 & SVM                 & -0.2000 & 0.8125 & 1.0000 & No \\
CatBoost            & Gradient Boosting   &  0.2000 & 0.8125 & 1.0000 & No \\
Decision Tree       & KNN                 & -0.2000 & 0.8125 & 1.0000 & No \\
CNN                 & MLP                 & -0.0667 & 1.0000 & 1.0000 & No \\
Gradient Boosting   & LightGBM            & -0.0667 & 1.0000 & 1.0000 & No \\
CatBoost            & LightGBM            &  0.0667 & 1.0000 & 1.0000 & No \\
LightGBM            & Stacking            & -0.0667 & 1.0000 & 1.0000 & No \\
LightGBM            & Random Forest       &  0.0667 & 1.0000 & 1.0000 & No \\

\end{longtable}
\endgroup

% =====================================================================
% Supplementary Table: Counterfactual feature-change frequencies
% Read off from bar charts (top_8 / top_10 / top_12 / top_14 / top_16 / top_all,
% freq + recomm variants). Values approximated to the nearest gridline/unit
% visible in each chart -- please spot-check against your underlying counts
% before submission, particularly Table 2's "All features" panel (39 finely
% spaced 1-5 count bars where transcription error is most likely).
% =====================================================================

\newpage
\section{\textcolor{black}{Recommendation Frequencies for Different Feature Sets}}
% ---------------------------------------------------------------------
% TABLE 1: Individual feature changes, bootstrapped mean frequency + 95% CI
% ---------------------------------------------------------------------
\small
\begin{longtable}{llrrr}
\caption{\textcolor{black}{Mean frequency of individual features changed in generated counterfactuals, by candidate feature-set size, with 95\% bootstrap confidence intervals.}} 
\label{tab:cf_freq_individual} \\
\toprule
\textbf{Feature set} & \textbf{Feature changed} & \textbf{Mean frequency} & \textbf{CI lower} & \textbf{CI upper} \\
\midrule
\endfirsthead
\multicolumn{5}{c}{\tablename\ \thetable\ -- continued} \\
\toprule
\textbf{Feature set} & \textbf{Feature changed} & \textbf{Mean frequency} & \textbf{CI lower} & \textbf{CI upper} \\
\midrule
\endhead
\bottomrule
\endfoot

\multirow{3}{*}{Top 8}  & JobInvolvement[dec] & 3  & 1  & 5  \\
  & JobInvolvement[inc] & 51 & 42 & 60 \\
  & Shift[inc]          & 69 & 63 & 75 \\
\midrule
\multirow{3}{*}{Top 10} & JobInvolvement[dec] & 9  & 4  & 14 \\
 & JobInvolvement[inc] & 49 & 40 & 58 \\
 & Shift[inc]          & 74 & 69 & 79 \\
\midrule
\multirow{4}{*}{Top 12} & JobInvolvement[dec] & 5  & 2  & 8  \\
 & OverTime[dec]       & 39 & 31 & 47 \\
 & JobInvolvement[inc] & 40 & 30 & 49 \\
 & Shift[inc]          & 58 & 49 & 67 \\
\midrule
\multirow{6}{*}{Top 14} & JobLevel[dec]       & 4  & 1  & 7  \\
 & BusinessTravel[dec] & 9  & 4  & 14 \\
 & JobInvolvement[inc] & 16 & 9  & 23 \\
 & OverTime[dec]       & 30 & 21 & 39 \\
 & Shift[inc]          & 31 & 22 & 40 \\
 & JobLevel[inc]       & 66 & 58 & 74 \\
\midrule
\multirow{7}{*}{Top 16} & BusinessTravel[inc] & 2  & 1  & 4  \\
 & JobLevel[dec]       & 2  & 1  & 4  \\
 & BusinessTravel[dec] & 4  & 1  & 7  \\
 & JobInvolvement[inc] & 13 & 7  & 19 \\
 & OverTime[dec]       & 26 & 18 & 34 \\
 & Shift[inc]          & 36 & 27 & 45 \\
 & JobLevel[inc]       & 68 & 60 & 76 \\
\midrule
\multirow{19}{*}{All features} & PercentSalaryHike[dec]        & 3  & 1  & 4  \\
 & JobSatisfaction[dec]          & 3  & 1  & 4  \\
 & MonthlyRate[dec]              & 3  & 1  & 4  \\
 & EnvironmentSatisfaction[inc]  & 3  & 1  & 4  \\
 & YearsSinceLastPromotion[dec]  & 3  & 1  & 4  \\
 & DailyRate[dec]                & 3  & 1  & 4  \\
 & WorkLifeBalance[inc]          & 3  & 1  & 6  \\
 & TrainingTimesLastYear[inc]    & 4  & 1  & 7  \\
 & TrainingTimesLastYear[dec]    & 4  & 1  & 7  \\
 & HourlyRate[dec]               & 4  & 1  & 7  \\
 & HourlyRate[inc]               & 4  & 1  & 7  \\
 & PercentSalaryHike[inc]        & 4  & 1  & 7  \\
 & YearsSinceLastPromotion[inc]  & 4  & 1  & 7  \\
 & MonthlyRate[inc]              & 4  & 1  & 7  \\
 & BusinessTravel[dec]           & 7  & 3  & 11 \\
 & JobInvolvement[inc]           & 10 & 4  & 16 \\
 & JobLevel[inc]                 & 19 & 12 & 26 \\
 & OverTime[dec]                 & 36 & 27 & 45 \\
 & Shift[inc]                    & 36 & 28 & 44 \\
\end{longtable}

% ---------------------------------------------------------------------
% TABLE 2: Feature COMBINATIONS changed together, raw frequency (no CI --
% source charts do not report bootstrap intervals for combinations)
% ---------------------------------------------------------------------
\small
\begin{longtable}{llr}
\caption{\textcolor{black}{Raw frequency of feature combinations jointly changed in generated counterfactuals by candidate feature-set size.}} 
\label{tab:cf_freq_combinations}\\
\toprule
\textbf{Feature set} & \textbf{Feature combination changed} & \textbf{Frequency} \\
\midrule
\endfirsthead
\multicolumn{3}{c}{\tablename\ \thetable\ -- continued} \\
\toprule
\textbf{Feature set} & \textbf{Feature combination changed} & \textbf{Frequency} \\
\midrule
\endhead
\bottomrule
\endfoot

\multirow{4}{*}{Top 8}  & JobInvolvement[dec], Shift[inc]              & 3  \\
  & JobInvolvement[inc]                          & 11 \\
  & Shift[inc]                                   & 26 \\
  & JobInvolvement[inc], Shift[inc]               & 40 \\
\midrule
\multirow{4}{*}{Top 10} & JobInvolvement[dec], Shift[inc]              & 9  \\
 & JobInvolvement[inc]                          & 6  \\
 & Shift[inc]                                   & 22 \\
 & JobInvolvement[inc], Shift[inc]               & 43 \\
\midrule
\multirow{9}{*}{Top 12} & JobInvolvement[dec], Shift[inc]              & 4  \\
 & OverTime[dec], Shift[dec]                    & 1  \\
 & Shift[inc]                                   & 15 \\
 & JobInvolvement[dec], OverTime[dec]           & 1  \\
 & JobInvolvement[inc]                          & 7  \\
 & OverTime[dec]                                & 9  \\
 & JobInvolvement[inc], Shift[inc]               & 22 \\
 & OverTime[dec], Shift[inc]                    & 17 \\
 & JobInvolvement[inc], OverTime[dec]           & 11 \\
\midrule
\multirow{14}{*}{Top 14} & BusinessTravel[dec], Shift[inc]              & 1  \\
 & BusinessTravel[dec], OverTime[dec]           & 1  \\
 & JobInvolvement[inc], OverTime[dec]           & 2  \\
 & JobLevel[inc]                                & 14 \\
 & BusinessTravel[dec], JobLevel[inc]           & 4  \\
 & JobLevel[dec], OverTime[dec]                 & 2  \\
 & JobLevel[dec], Shift[inc]                    & 2  \\
 & JobInvolvement[inc], JobLevel[inc]           & 8  \\
 & BusinessTravel[dec], JobInvolvement[inc]     & 3  \\
 & JobInvolvement[inc], Shift[inc]               & 3  \\
 & JobLevel[inc], Shift[inc]                    & 18 \\
 & Shift[inc]                                   & 4  \\
 & OverTime[dec], Shift[inc]                    & 3  \\
 & JobLevel[inc], OverTime[dec]                 & 22 \\
\midrule
\multirow{16}{*}{Top 16} & BusinessTravel[dec], Shift[inc]              & 1  \\
 & BusinessTravel[dec]                          & 1  \\
 & BusinessTravel[dec], OverTime[dec]           & 1  \\
 & JobInvolvement[inc]                          & 1  \\
 & BusinessTravel[dec], JobLevel[inc]           & 1  \\
 & BusinessTravel[inc], JobLevel[inc]           & 1  \\
 & BusinessTravel[inc], OverTime[dec]           & 1  \\
 & JobLevel[inc]                                & 15 \\
 & JobLevel[dec], Shift[inc]                    & 2  \\
 & OverTime[dec]                                & 2  \\
 & JobInvolvement[inc], Shift[inc]               & 3  \\
 & Shift[inc]                                   & 4  \\
 & JobInvolvement[inc], JobLevel[inc]           & 9  \\
 & OverTime[dec], Shift[inc]                    & 3  \\
 & JobLevel[inc], Shift[inc]                    & 23 \\
 & JobLevel[inc], OverTime[dec]                 & 19 \\
\midrule
 & WorkLifeBalance[inc]                   & 1 \\
 & BusinessTravel[dec], DailyRate[dec]                & 1 \\
 & EnvironmentSatisfaction[inc], Shift[inc]           & 1 \\
 & DailyRate[inc], OverTime[dec]                      & 1 \\
 & JobSatisfaction[dec], Shift[inc]                   & 2 \\
 & HourlyRate[inc], Shift[inc]                        & 1 \\
 & JobInvolvement[inc], OverTime[dec]                 & 1 \\
 & BusinessTravel[dec], OverTime[dec]                 & 2 \\
 & HourlyRate[inc], OverTime[dec]                     & 1 \\
 & OverTime[dec], PercentSalaryHike[inc]               & 1 \\
 & HourlyRate[inc], TrainingTimesLastYear[inc]        & 1 \\
 & HourlyRate[inc]                                    & 2 \\
 & JobInvolvement[inc], PercentSalaryHike[inc]        & 1 \\
 & HourlyRate[dec], JobInvolvement[inc]               & 1 \\
 & EnvironmentSatisfaction[dec], Shift[inc]           & 1 \\
 & DailyRate[dec], Shift[inc]                         & 1 \\
 & HourlyRate[dec], Shift[inc]                        & 1 \\
 All features & HourlyRate[inc], JobSatisfaction[inc]              & 1 \\
 & JobInvolvement[dec], Shift[inc]                    & 2 \\
 & OverTime[dec], TrainingTimesLastYear[inc]          & 1 \\
 & BusinessTravel[dec], JobLevel[inc]                 & 1 \\
 & JobSatisfaction[dec], OverTime[dec]                & 1 \\
 & MonthlyRate[inc], OverTime[dec]                    & 1 \\
 & EnvironmentSatisfaction[dec], JobLevel[inc]        & 1 \\
 & JobInvolvement[inc], JobLevel[inc]                 & 1 \\
 & JobLevel[inc]                                      & 4 \\
 & TrainingTimesLastYear[inc]                         & 1 \\
 & TrainingTimesLastYear[inc], YearsSinceLastPromotion[inc] & 1 \\
 & BusinessTravel[dec], PercentSalaryHike[dec]        & 1 \\
 & JobLevel[inc], OverTime[dec]                       & 5 \\
 & MonthlyRate[dec], OverTime[dec]                    & 2 \\
 & HourlyRate[dec], OverTime[dec]                     & 5 \\
 & OverTime[dec], TrainingTimesLastYear[dec]          & 1 \\
 & MonthlyRate[dec], Shift[inc]                       & 1 \\
 & Shift[inc], YearsSinceLastPromotion[dec]           & 1 \\
 & EnvironmentSatisfaction[dec], OverTime[dec]        & 2 \\
 & Shift[inc]                                         & 4 \\
 & HourlyRate[inc], JobLevel[inc]                     & 1 \\
 & Shift[inc], YearsSinceLastPromotion[inc]           & 3 \\
 & HourlyRate[inc], JobInvolvement[inc]               & 1 \\
 & JobInvolvement[inc]                                & 2 \\
 & Shift[inc], TrainingTimesLastYear[inc]              & 4 \\
 & OverTime[dec]                                      & 3 \\
 & DailyRate[inc], Shift[inc]                         & 1 \\
 & JobLevel[inc], Shift[inc]                          & 5 \\
 & JobInvolvement[inc], Shift[inc]                    & 5 \\
 & OverTime[dec], Shift[inc]                          & 5 \\
\end{longtable}

% %\vskip3pt

% % \bio{}
% % Author biography without author photo.
% % Author biography. Author biography. Author biography.
% % \endbio

% % \bio{figs/cas-pic1}
% % Author biography with author photo.
% % Author biography. Author biography. Author biography.
% % \endbio

% % \vskip3pc

% % \bio{figs/cas-pic1}
% % Author biography with author photo.
% % Author biography. Author biography. Author biography.
% % \endbio


\end{document}

